% Les maladies et l'importance de la pratique du sport — Allemand 1ère A — Cours signé OnBuch+
% Programme officiel MINESEC. Compiler avec : tectonic ALA12-les-maladies-et-l-importance-de-la-pratique-du-sport.tex
\newcommand{\DOCMATIERE}{Allemand}
\newcommand{\DOCNIVEAU}{1ère A}
\newcommand{\DOCMODULE}{Chapitre 3 : Environnement, santé et bien-être}
\newcommand{\DOCLECON}{ALA12}
\newcommand{\DOCTITRE}{Les maladies et l'importance de la pratique du sport}
% OnBuch+ — préambule « cours signés » ENRICHI (compilé avec tectonic / XeLaTeX)
% Système visuel « Pop » d'OnBuch+ : fond crème (jamais blanc), bordures encre
% épaisses, ombres « dures » décalées, titres Archivo Black en capitales.
% Version MATHÉMATIQUES : graphes (pgfplots/tikz), boîtes pédagogiques
% (définition, propriété, méthode, attention, le savais-tu, exercice résolu,
% exercice, TP, bilan), page de garde et signature OnBuch+ ancrée en bas.
%
% Macros attendues dans main.tex AVANT \input{preamble} :
%   \DOCMATIERE  \DOCNIVEAU  \DOCMODULE  \DOCLECON (numéro)  \DOCTITRE
\documentclass[11pt]{article}

\usepackage{fontspec}
\usepackage[ngerman,french]{babel} % français = langue principale ; l'allemand via \all{..} / textebox
\usepackage[a4paper,margin=2cm,top=3.4cm,bottom=2.8cm,headheight=1.4cm,footskip=1.3cm]{geometry}
\usepackage{amsmath,amssymb,amsfonts}
\usepackage{mathtools} % \xmapsto, \coloneqq... souvent utilisés par les modèles
\usepackage{cancel} % simplifications barrées (bilans de forces)
% \abs{z} (module, valeur absolue) et \norm{u} (norme), utilisés par les modèles
\providecommand{\abs}{}\let\abs\relax\DeclarePairedDelimiter{\abs}{\lvert}{\rvert}
\providecommand{\norm}{}\let\norm\relax\DeclarePairedDelimiter{\norm}{\lVert}{\rVert}
\usepackage[table]{xcolor} % [table] : fournit aussi \rowcolors (alternance de lignes)
\usepackage{pagecolor}
\usepackage{tikz}
\usetikzlibrary{arrows.meta,positioning,calc,shapes.geometric,shapes.misc,shapes.arrows,decorations.pathmorphing,decorations.pathreplacing,decorations.markings,patterns,shadows,mindmap,trees,fit,backgrounds,matrix,intersections,angles,quotes}
\usepackage{pgfplots}
\usepackage[version=4]{mhchem} % équations nucléaires \ce{^{235}_{92}U}
\usepackage[european,siunitx]{circuitikz} % schémas électriques
\pgfplotsset{compat=1.18}
\usepgfplotslibrary{fillbetween,groupplots} % aires entre courbes (calcul intégral)
\usepackage{enumitem}
\usepackage{fancyhdr}
% [most] charge toutes les bibliothèques tcolorbox (listings, minted, raster,
% documentation...) et gonfle inutilement la mémoire TeX sur un document long
% (~300 boîtes) : on ne charge que ce qui est réellement utilisé.
\usepackage[skins,breakable]{tcolorbox}
\usepackage{graphicx}
\usepackage{colortbl}
\usepackage{booktabs}
\usepackage{tabularx}
\usepackage{diagbox} % case d'en-tête diagonale des tableaux à double entrée
% Colonnes extensibles centrée / gauche / droite (C, L, R), souvent utilisées par les modèles
\newcolumntype{C}{>{\centering\arraybackslash}X}
\newcolumntype{L}{>{\raggedright\arraybackslash}X}
\newcolumntype{R}{>{\raggedleft\arraybackslash}X}
\usepackage{array}
\usepackage{multicol}
\usepackage{multirow}
\usepackage{siunitx}
% parse-numbers=false : accepte aussi \SI{2,5 \times 10^{-3}}{...}
\sisetup{locale=FR,output-decimal-marker={,},group-minimum-digits=4,parse-numbers=false}
\DeclareSIUnit\ohm{\ensuremath{\symup{\Omega}}} % U+2126 absent de la police
\DeclareSIUnit\division{div} % réglages d'oscilloscope (ms/div, V/div)
\DeclareSIUnit\tr{tr} % tours (vitesse de rotation, ex. \SI{1500}{\tr\per\minute})
\DeclareSIUnit\molar{M} % concentration molaire (ex. \SI{3}{\molar}, \SI{1}{\micro\molar})
\DeclareSIUnit\an{an} % année (le modèle confond parfois avec \year, un compteur TeX) — ex. \SI{5}{\kilo\meter\per\an}
\DeclareSIUnit\FCFA{FCFA} % franc CFA (ex. \SI{500}{\FCFA\per\kilo\gram})
\DeclareSIUnit\degreeBrix{\textdegree Brix} % teneur en sucre d'un sirop/jus (réfractométrie)
\DeclareSIUnit\month{mois} % le modèle utilise parfois \month, un compteur TeX, comme unité de durée
\usepackage{qrcode}
% \degree : commande fréquemment utilisée par les modèles pour un angle ou une
% température en dehors de \SI{}{} (non fournie par les paquets chargés).
\providecommand{\degree}{^{\circ}}
% \qed (fin de démonstration), utilisé par les modèles sans amsthm
\providecommand{\qed}{\hfill$\square$}
% \barre : double barre des valeurs interdites dans un tableau de variations (tkz-tab)
\providecommand{\barre}{\ensuremath{\|}}
% environnement proof (amsthm non chargé) utilisé par les modèles
\ifdefined\proof\else\newenvironment{proof}[1][Démonstration]{\par\noindent\textit{#1.}\ }{\qed\par}\fi
\usepackage{lastpage}
\usepackage{hyperref}
\hypersetup{hidelinks}

% ============================================================
% Palette « Pop » OnBuch+ — mêmes hex que la classe P (plus_ui.dart)
% ============================================================
\definecolor{poporange}{HTML}{F59321}
\definecolor{popdark}{HTML}{E07A0C}
\definecolor{popcream}{HTML}{FFF6E8}
\definecolor{popink}{HTML}{1C1714}
\definecolor{poppurple}{HTML}{7A5AE0}
\definecolor{popgreen}{HTML}{2E9E6A}
\definecolor{poppink}{HTML}{E0457B}
\definecolor{popblue}{HTML}{2F7DE1}
\definecolor{popgold}{HTML}{C98A12}
\definecolor{popmuted}{HTML}{8A7F76}
\definecolor{popline}{HTML}{EADFD2}
\definecolor{popgray}{HTML}{8A7F76}
\colorlet{popgrey}{popgray}
% Alias tolérants (le modèle nomme parfois les couleurs différemment)
\colorlet{popwhite}{white}
\colorlet{popblack}{popink}
\colorlet{popred}{poppink}
\colorlet{popyellow}{popgold}
% Teintes claires (fonds de tableaux/encadrés) — le modèle nomme parfois
% les couleurs avec un suffixe L (ex. popblueL) sans les avoir définies.
\colorlet{poporangeL}{poporange!20}
\colorlet{popdarkL}{popdark!20}
\colorlet{poppurpleL}{poppurple!20}
\colorlet{popgreenL}{popgreen!20}
\colorlet{poppinkL}{poppink!20}
\colorlet{popblueL}{popblue!20}
\colorlet{popgoldL}{popgold!20}
\colorlet{popredL}{popred!20}
\colorlet{popgrayL}{popgray!20}
% Teintes claires pour fonds de tableaux / zones
\colorlet{popblueL}{popblue!10!white}
\colorlet{poporangeL}{poporange!14!white}
\colorlet{popgreenL}{popgreen!12!white}
\colorlet{poppurpleL}{poppurple!12!white}
\colorlet{poppinkL}{poppink!12!white}
\colorlet{popgoldL}{popgold!14!white}

\pagecolor{popcream}

% ============================================================
% Typographie
% ============================================================
\defaultfontfeatures{Ligatures=TeX}
\setmainfont{PlusJakartaSans-Medium.ttf}[Path=../../../fonts/,BoldFont=PlusJakartaSans-Bold.ttf]
% Maths en sans-serif (Fira Math) avec chiffres et lettres droites de Plus
% Jakarta, pour que les formules se fondent dans le texte.
\usepackage[math-style=ISO,bold-style=ISO]{unicode-math}
\setmathfont{FiraMath-Regular.otf}[Path=../../../fonts/]
\setmathfont{PlusJakartaSans-Medium.ttf}[Path=../../../fonts/,range={up/num,up/{Latin,latin},\mathup}]
\usepackage{icomma} % virgule décimale sans espace en mode math
% Caractères mathématiques Unicode absents des polices de texte (Plus Jakarta,
% Archivo Black) : rendus par leur équivalent mathématique, en texte comme en
% formule — sinon ils s'affichent en carré vide (ex. « Arithmétique dans ℤ »).
\usepackage{newunicodechar}
\newunicodechar{ℝ}{\ensuremath{\mathbb{R}}}
\newunicodechar{ℤ}{\ensuremath{\mathbb{Z}}}
\newunicodechar{ℂ}{\ensuremath{\mathbb{C}}}
\newunicodechar{ℕ}{\ensuremath{\mathbb{N}}}
\newunicodechar{ℚ}{\ensuremath{\mathbb{Q}}}
\newunicodechar{𝔻}{\ensuremath{\mathbb{D}}}
\newunicodechar{α}{\ensuremath{\alpha}}
\newunicodechar{β}{\ensuremath{\beta}}
\newunicodechar{γ}{\ensuremath{\gamma}}
\newunicodechar{δ}{\ensuremath{\delta}}
\newunicodechar{ε}{\ensuremath{\varepsilon}}
\newunicodechar{θ}{\ensuremath{\theta}}
\newunicodechar{λ}{\ensuremath{\lambda}}
\newunicodechar{μ}{\ensuremath{\mu}}
\newunicodechar{σ}{\ensuremath{\sigma}}
\newunicodechar{τ}{\ensuremath{\tau}}
\newunicodechar{φ}{\ensuremath{\varphi}}
\newunicodechar{ψ}{\ensuremath{\psi}}
\newunicodechar{ω}{\ensuremath{\omega}}
\newunicodechar{Σ}{\ensuremath{\Sigma}}
% majuscules produites par \MakeUppercase dans les titres (α-aminés -> Α)
\newunicodechar{Α}{\ensuremath{\alpha}}
\newunicodechar{Β}{\ensuremath{\beta}}
\newunicodechar{Γ}{\ensuremath{\Gamma}}
\newunicodechar{Δ}{\ensuremath{\Delta}}
\newunicodechar{Ε}{\ensuremath{\varepsilon}}
\newunicodechar{Θ}{\ensuremath{\Theta}}
\newunicodechar{Λ}{\ensuremath{\Lambda}}
\newunicodechar{Μ}{\ensuremath{\mu}}
\newunicodechar{Τ}{\ensuremath{\tau}}
\newunicodechar{Φ}{\ensuremath{\Phi}}
\newunicodechar{Ψ}{\ensuremath{\Psi}}
\newunicodechar{Ω}{\ensuremath{\Omega}}
\newunicodechar{↦}{\ensuremath{\mapsto}}
\newunicodechar{∈}{\ensuremath{\in}}
\newunicodechar{∉}{\ensuremath{\notin}}
\newunicodechar{∧}{\ensuremath{\wedge}}
\newunicodechar{⇔}{\ensuremath{\Leftrightarrow}}
\newunicodechar{⟺}{\ensuremath{\Longleftrightarrow}}
\newunicodechar{⇒}{\ensuremath{\Rightarrow}}
\newunicodechar{⟹}{\ensuremath{\Longrightarrow}}
\newunicodechar{∘}{\ensuremath{\circ}}
\newunicodechar{∪}{\ensuremath{\cup}}
\newunicodechar{∩}{\ensuremath{\cap}}
\newunicodechar{≡}{\ensuremath{\equiv}}
\newunicodechar{⊂}{\ensuremath{\subset}}
\newunicodechar{⊥}{\ensuremath{\perp}}
\newunicodechar{∃}{\ensuremath{\exists}}
\newunicodechar{∀}{\ensuremath{\forall}}
\newunicodechar{∛}{\ensuremath{\sqrt[3]{\,}}}
\newunicodechar{∜}{\ensuremath{\sqrt[4]{\,}}}
\newunicodechar{⃗}{\ensuremath{{}^{\to}}}
\newunicodechar{ⁿ}{\ifmmode^{n}\else\textsuperscript{\ensuremath{n}}\fi}
\newunicodechar{ˣ}{\ifmmode^{x}\else\textsuperscript{\ensuremath{x}}\fi}
\newunicodechar{ᵇ}{\ifmmode^{b}\else\textsuperscript{\ensuremath{b}}\fi}
\newunicodechar{ʳ}{\ifmmode^{r}\else\textsuperscript{\ensuremath{r}}\fi}
\newunicodechar{ᵘ}{\ifmmode^{u}\else\textsuperscript{\ensuremath{u}}\fi}
\newunicodechar{ʸ}{\ifmmode^{y}\else\textsuperscript{\ensuremath{y}}\fi}
\newunicodechar{ᵃ}{\ifmmode^{a}\else\textsuperscript{\ensuremath{a}}\fi}
\newunicodechar{ᵉ}{\ifmmode^{e}\else\textsuperscript{\ensuremath{e}}\fi}
\newunicodechar{ᵏ}{\ifmmode^{k}\else\textsuperscript{\ensuremath{k}}\fi}
\newunicodechar{ᵖ}{\ifmmode^{p}\else\textsuperscript{\ensuremath{p}}\fi}
\newunicodechar{ᴺ}{\ifmmode^{N}\else\textsuperscript{\ensuremath{N}}\fi}
\newunicodechar{ᶜ}{\ifmmode^{c}\else\textsuperscript{\ensuremath{c}}\fi}
\newunicodechar{ʰ}{\ifmmode^{h}\else\textsuperscript{\ensuremath{h}}\fi}
\newunicodechar{ᵠ}{\ifmmode^{q}\else\textsuperscript{\ensuremath{q}}\fi}
\newunicodechar{ₙ}{\ifmmode_{n}\else\textsubscript{\ensuremath{n}}\fi}
\newunicodechar{ₐ}{\ifmmode_{a}\else\textsubscript{\ensuremath{a}}\fi}
\newunicodechar{ᵢ}{\ifmmode_{i}\else\textsubscript{\ensuremath{i}}\fi}
\newunicodechar{ᵣ}{\ifmmode_{r}\else\textsubscript{\ensuremath{r}}\fi}
\newunicodechar{ₖ}{\ifmmode_{k}\else\textsubscript{\ensuremath{k}}\fi}
\newunicodechar{ᵦ}{\ifmmode_{\beta}\else\textsubscript{\ensuremath{\beta}}\fi}
\newunicodechar{ₓ}{\ifmmode_{x}\else\textsubscript{\ensuremath{x}}\fi}
\newfontfamily\popdisplay{ArchivoBlack-Regular.ttf}[Path=../../../fonts/]
\newfontfamily\poplogo{SpaceGrotesk-Bold.ttf}[Path=../../../fonts/]
\newfontfamily\popsemi{PlusJakartaSans-SemiBold.ttf}[Path=../../../fonts/]
\newfontfamily\popxbold{PlusJakartaSans-ExtraBold.ttf}[Path=../../../fonts/]
\newfontfamily\popxboldls{PlusJakartaSans-ExtraBold.ttf}[Path=../../../fonts/,LetterSpace=3.0]

\setlist[enumerate]{leftmargin=1.7em,itemsep=5pt,topsep=4pt}
\setlist[itemize]{leftmargin=1.7em,itemsep=4pt,topsep=4pt,label={\color{poporange}\textbullet}}
\setlength{\parindent}{0pt}
\setlength{\parskip}{5pt}
\renewcommand{\arraystretch}{1.25}

% pgfplots : style Pop par défaut
\pgfplotsset{
  every axis/.append style={axis line style={popink,line width=1pt},tick style={popink},
    grid style={popline,line width=0.5pt},label style={font=\small},tick label style={font=\footnotesize}},
  popcurve/.style={poporange,line width=1.8pt,no markers,smooth},
}
% Même style disponible dans un \draw TikZ simple (hors pgfplots)
\tikzset{popcurve/.style={poporange,line width=1.8pt}}
\tikzset{popgrid/.style={popline,very thin}}
\colorlet{popcurve}{poporange} % le nom du style est parfois utilisé comme couleur

% ============================================================
% Pill / squiggle
% ============================================================
\newcommand{\poppill}[3]{%
  \tikz[baseline=-0.55ex]\node[draw=popink,line width=1.2pt,rounded corners=2.6pt,%
    fill=#2,inner xsep=6.5pt,inner ysep=3pt]{%
    \popxboldls\fontsize{7}{7}\selectfont\color{#3}\MakeUppercase{#1}};%
}
\newcommand{\popsquiggle}[1]{%
  \tikz[baseline=-0.4ex]{
    \draw[#1,line width=2.6pt,line cap=round]
      plot [smooth] coordinates {(0,0.09) (0.34,0.01) (0.68,0.21) (1.02,0.01) (1.36,0.10)};
  }%
}
% Mot-clé surligné (marqueur orange sous le texte)
\newcommand{\cle}[1]{{\popxbold\color{popdark}#1}}

% ============================================================
% En-tête / pied
% ============================================================
\pagestyle{fancy}
\fancyhf{}
\renewcommand{\headrulewidth}{0pt}
\renewcommand{\footrulewidth}{0pt}
\lhead{\raisebox{-0.15ex}{\poplogo\fontsize{13}{13}\selectfont\color{popink}OnBuch\color{poporange}+}}
\rhead{\poppill{\DOCMATIERE\ \textbullet\ \DOCNIVEAU\ \textbullet\ Leçon \DOCLECON}{white}{popink}}
\lfoot{\popxbold\fontsize{7.6}{7.6}\selectfont\color{popmuted}\MakeUppercase{Cours signé OnBuch+}\ \textbullet\ {\popsemi\DOCTITRE}}
\rfoot{\tikz[baseline=-0.6ex]\node[rounded corners=8.5pt,fill=popink,minimum height=17pt,minimum width=17pt,inner xsep=5pt,inner sep=0]{\popxbold\fontsize{8}{8}\selectfont\color{popcream}\thepage\,/\,\pageref*{LastPage}};}

% ============================================================
% Titres
% ============================================================
\newcommand{\chaptitre}[2]{%
  \begin{tcolorbox}[enhanced,colback=poporange,colframe=popink,boxrule=2.6pt,arc=11pt,
      drop shadow=popink,left=15pt,right=15pt,top=12pt,bottom=11pt,before skip=3pt,after skip=18pt]
    \poppill{#2}{popink}{popcream}\par\vspace{9pt}
    {\raggedright\hyphenpenalty=10000\exhyphenpenalty=10000\popdisplay\fontsize{22}{24}\selectfont\color{popcream}\MakeUppercase{#1}\par}
  \end{tcolorbox}
}
% Section numérotée : pastille orange avec le numéro + titre Archivo
\newcounter{popsec}
\newcommand{\coursec}[1]{%
  \stepcounter{popsec}\setcounter{popsub}{0}%
  \par\vspace{16pt}\noindent
  \begin{minipage}{\linewidth}
  \tikz[baseline=-0.7ex]\node[draw=popink,line width=1.6pt,fill=poporange,rounded corners=4pt,minimum size=22pt,inner sep=0]{\popdisplay\fontsize{12}{12}\selectfont\color{popcream}\thepopsec};%
  \hspace{8pt}{\popdisplay\fontsize{15}{17}\selectfont\color{popink}\MakeUppercase{#1}}\par
  \vspace{4pt}\noindent{\color{poporange}\rule{36pt}{3.6pt}}
  \end{minipage}\par\nopagebreak\vspace{8pt}
}
\newcounter{popsub}
\newcommand{\courssub}[1]{%
  \stepcounter{popsub}%
  \par\vspace{9pt}\noindent{\popxbold\fontsize{11.5}{13}\selectfont\color{popink}{\color{poporange}\thepopsec.\thepopsub}\hspace{6pt}#1}\par\nopagebreak\vspace{4pt}
}

% ============================================================
% Boîtes pédagogiques — même recette : fond blanc, bordure épaisse colorée,
% ombre dure encre, badge plein. Titre optionnel [#1].
% ============================================================
\tcbset{popbase/.style={breakable,enhanced,colback=white,boxrule=2.3pt,arc=9pt,
  shadow={3pt}{-3pt}{0pt}{popink},left=14pt,right=14pt,top=11pt,bottom=11pt,before skip=11pt,after skip=15pt,
  fontupper=\popsemi\fontsize{10.6}{14.5}\selectfont\color{popink},
  fontlower=\popsemi\fontsize{10.6}{14.5}\selectfont\color{popink}}}
% Coche dessinée (le glyphe ✓ n'existe pas dans les polices du document)
\newcommand{\popcheck}[1]{\tikz[baseline=-0.5ex]\draw[#1,line width=1.6pt,line cap=round,line join=round](-0.13,0.0)--(-0.04,-0.09)--(0.13,0.1);}
\providecommand{\coche}{\popcheck{popgreen}}
\providecommand{\resultat}[1]{\par\smallskip\noindent\fcolorbox{popgreen}{popgreenL}{\parbox{\dimexpr\linewidth-2\fboxsep-2\fboxrule\relax}{\textbf{Résultat.} #1}}\par\smallskip}
\newcommand{\popboxhead}[3]{\poppill{#1}{#2}{white}\ifx\relax#3\relax\else\hspace{6pt}{\popxbold\fontsize{10.6}{12}\selectfont\color{popink}#3}\fi\par\vspace{7pt}}

\newtcolorbox{aretenir}[1][]{popbase,colframe=popgreen,before upper={\popboxhead{À retenir}{popgreen}{#1}}}
% Texte en allemand (document de lecture, dialogue, extrait) : cadre
% d'étude. Un titre optionnel (source, auteur) s'affiche dans l'en-tête.
\newtcolorbox{textebox}[1][]{popbase,colframe=popblue,before upper={\popboxhead{Text}{popblue}{#1}\begin{otherlanguage}{ngerman}},after upper={\end{otherlanguage}}}
% Marques ✓ / ✗ (forme correcte / fautive) souvent utilisées par les modèles
\providecommand{\cross}{\textcolor{poppink}{\ensuremath{\times}}}
\providecommand{\xmark}{\cross}\providecommand{\crossmark}{\cross}\providecommand{\cmark}{\ensuremath{\checkmark}}
% Ligne de réponse / justification à compléter (inventée par les modèles)
\providecommand{\justification}[1]{\par\noindent\textit{Justification~:} \dotfill\par}
% \textipa (package tipa non chargé) : la police ne couvre pas l'API -> simple passage
\providecommand{\textipa}[1]{#1}
% Soulignement (ulem non chargé) : \uline, \ul
\providecommand{\uline}[1]{\underline{#1}}\providecommand{\ul}[1]{\underline{#1}}
% \alert (beamer) inventée par les modèles : texte mis en évidence
\providecommand{\alert}[1]{\textbf{\textcolor{poppink}{#1}}}
% variantes ASCII d'ordinaux écrites en macro
\providecommand{\iere}{\textsuperscript{re}}\providecommand{\ieme}{\textsuperscript{e}}\providecommand{\ere}{\textsuperscript{re}}\providecommand{\eme}{\textsuperscript{e}}\providecommand{\er}{\textsuperscript{er}}\providecommand{\ie}{\textsuperscript{e}}
% Ordinaux français écrits en macro par les modèles : 1\ière, 3\ième
\newcommand{\ière}{\textsuperscript{re}}\newcommand{\ième}{\textsuperscript{e}}
% \exemple (inventée par les modèles) : étiquette « Exemple : »
\providecommand{\exemple}{\textbf{Exemple~:}\ }
% \square utilisable aussi hors mode math (cases à cocher)
\AtBeginDocument{\let\squareMath\square\renewcommand{\square}{\ensuremath{\squareMath}}}
% Mot ou phrase en allemand à l'intérieur d'un paragraphe français
\newcommand{\all}[1]{\ifmmode\text{\textit{#1}}\else\begingroup\selectlanguage{ngerman}\itshape #1\endgroup\fi}
\let\esp\all % alias toléré
\newtcolorbox{exemplebox}[1][]{popbase,colframe=poporange,before upper={\popboxhead{Exemple}{poporange}{#1}}}
\newtcolorbox{definition}[1][]{popbase,colframe=popblue,before upper={\popboxhead{Définition}{popblue}{#1}}}
\newtcolorbox{propriete}[1][]{popbase,colframe=poppurple,before upper={\popboxhead{Propriété}{poppurple}{#1}}}
\newtcolorbox{methode}[1][]{popbase,colframe=popgold,before upper={\popboxhead{Méthode}{popgold}{#1}}}
\newtcolorbox{attention}[1][]{popbase,colframe=poppink,before upper={\popboxhead{Attention}{poppink}{#1}}}
\newtcolorbox{experience}[1][]{popbase,colframe=popblue,colback=popblueL,before upper={\popboxhead{Expérience}{popblue}{#1}}}
% « Le savais-tu ? » : carte violette pleine (contexte, culture, Cameroun)
\newtcolorbox{savaistu}[1][]{popbase,colback=poppurple,colframe=popink,
  fontupper=\popsemi\fontsize{10.4}{14.2}\selectfont\color{white},
  before upper={\poppill{Le savais-tu\,?}{white}{poppurple}\ifx\relax#1\relax\else\hspace{6pt}{\popxbold\color{white}#1}\fi\par\vspace{7pt}}}
% Exercice résolu : énoncé en haut, \tcblower puis la solution sur fond crème
\newcounter{popexo}\newcounter{popexr}
\newtcolorbox{exoresolu}[1][]{popbase,colframe=poporange,colbacklower=popcream,
  segmentation style={popink,line width=1.2pt,dashed},
  before upper={\stepcounter{popexr}\popboxhead{Exercice résolu \thepopexr}{poporange}{#1}},
  before lower={\poppill{Solution}{popink}{popcream}\par\vspace{6pt}}}
% Exercice (non résolu en place) — la correction est dans la partie Corrigés
\newtcolorbox{exercice}[1][]{popbase,colframe=popink,
  before upper={\stepcounter{popexo}\popboxhead{Exercice \thepopexo}{popink}{#1}}}
% Corrigé
\newtcolorbox{corrige}[1][]{popbase,colframe=popgreen,colback=popgreenL,
  before upper={\popboxhead{Corrigé}{popgreen}{#1}}}
% Objectifs / prérequis / bilan
\newtcolorbox{objectifs}{popbase,colframe=popink,colback=poporangeL,
  before upper={\popboxhead{Objectifs de la leçon}{popink}{}}}
\newtcolorbox{prerequis}{popbase,colframe=popink,colback=popblueL,
  before upper={\popboxhead{Prérequis}{popblue}{}}}
\newtcolorbox{bilan}{popbase,colframe=popink,colback=white,boxrule=2.8pt,
  before upper={\popboxhead{Fiche bilan}{poporange}{L'essentiel à connaître pour l'examen}}}
% Figure encadrée avec légende
\newtcolorbox{popfigure}[1][]{enhanced,breakable=false,colback=white,colframe=popline,boxrule=1.4pt,arc=8pt,
  left=10pt,right=10pt,top=10pt,bottom=8pt,before skip=10pt,after skip=12pt,halign=center,
  lower separated=false,halign lower=center,
  fontlower=\popsemi\fontsize{9}{11.5}\selectfont\color{popmuted},
  #1}
\newcommand{\legende}[1]{\par\vspace{6pt}{\popsemi\fontsize{9}{11.5}\selectfont\color{popmuted}#1}}

% ============================================================
% Signature OnBuch+
% ============================================================
\newcommand{\poplogoinline}[1]{{\poplogo\fontsize{#1}{#1}\selectfont\color{popink}OnBuch\color{poporange}+}}
\newcommand{\signatureonbuch}{%
  \begin{tcolorbox}[enhanced,colback=popink,colframe=popink,boxrule=2.2pt,arc=10pt,
      drop shadow=poporange,left=14pt,right=14pt,top=10pt,bottom=10pt,before skip=0pt,after skip=0pt]
    \tikz[baseline=-0.7ex]\node[circle,fill=popgreen,draw=popcream,line width=1.3pt,minimum size=22pt,inner sep=0]{\popcheck{white}};%
    \hspace{9pt}\begin{minipage}[c]{0.62\linewidth}
      {\popxbold\fontsize{12}{13}\selectfont\color{popcream}Signé\ \textbullet\ {\poplogo OnBuch\color{poporange}+}}\par\vspace{2pt}
      {\raggedright\popsemi\fontsize{8}{10.5}\selectfont\color{popline}\MakeUppercase{Équipe pédagogique OnBuch+}\\\MakeUppercase{Conforme au programme officiel MINESEC}\par}
    \end{minipage}\hfill
    {\poplogo\fontsize{22}{22}\selectfont\color{popcream}OnBuch\color{poporange}+}
  \end{tcolorbox}%
}

% ============================================================
% Page de garde
% #1 titre · #2 matière · #3 niveau · #4 module · #5 n° leçon · #6 durée ·
% #7 description / objectifs (texte ou itemize)
% ============================================================
\newcommand{\pagegarde}[7]{%
  \thispagestyle{empty}
  \newgeometry{a4paper,margin=1.7cm,top=1.6cm,bottom=1.4cm}%
  \begin{tikzpicture}[remember picture,overlay]
    \fill[poporange!18] ([xshift=-3.2cm,yshift=-3.6cm]current page.north east) circle (4.6cm);
    \fill[poppurple!14] ([xshift=2.4cm,yshift=7.6cm]current page.south west) circle (3.8cm);
  \end{tikzpicture}%
  {\poplogo\fontsize{30}{30}\selectfont\color{popink}OnBuch\color{poporange}+}\hfill
  \raisebox{0.6ex}{\poppill{Cours signé \textbullet\ Premium}{popink}{popcream}}\par
  \vspace{1pt}\popsquiggle{poppurple}\par
  \vspace{8pt}
  \begin{tcolorbox}[enhanced,colback=poporange,colframe=popink,boxrule=2.8pt,arc=13pt,
      drop shadow=popink,left=18pt,right=18pt,top=12pt,bottom=13pt,after skip=10pt]
    \poppill{#2 \textbullet\ #3}{popink}{popcream}\hspace{5pt}\poppill{#4}{white}{popink}\par\vspace{8pt}
    {\popdisplay\fontsize{14}{14}\selectfont\color{popink}LEÇON #5}\par\vspace{4pt}
    {\raggedright\hyphenpenalty=10000\exhyphenpenalty=10000\popdisplay\fontsize{25}{27}\selectfont\color{popcream}\MakeUppercase{#1}\par}
    \vspace{8pt}
    {\popxbold\fontsize{9.5}{9.5}\selectfont\color{popink}\MakeUppercase{Durée conseillée : #6}}
  \end{tcolorbox}
  \begin{tcolorbox}[enhanced,colback=white,colframe=popink,boxrule=2.2pt,arc=10pt,
      drop shadow=popink,left=16pt,right=16pt,top=10pt,bottom=10pt,after skip=8pt]
    \poppill{Dans cette leçon}{poppurple}{white}\par\vspace{5pt}
    \popsemi\fontsize{9.3}{12.6}\selectfont\color{popink}#7
  \end{tcolorbox}
  \noindent\begin{minipage}[t]{0.487\linewidth}
    \begin{tcolorbox}[enhanced,colback=popgreen,colframe=popink,boxrule=2.2pt,arc=10pt,
        drop shadow=popink,left=11pt,right=11pt,top=10pt,bottom=10pt,height=3.1cm,valign=center]
      \tcbox[colback=white,colframe=popink,boxrule=1.3pt,arc=5pt,boxsep=0pt,nobeforeafter,
        left=3pt,right=3pt,top=3pt,bottom=3pt,valign=center]{\qrcode[height=1.9cm]{https://play.google.com/store/apps/details?id=com.luvvix.onbuch&hl=fr}}%
      \hspace{8pt}\begin{minipage}[c]{0.5\linewidth}
        {\popdisplay\fontsize{11}{12.5}\selectfont\color{white}\MakeUppercase{Télécharge OnBuch}}\par\vspace{4pt}
        {\raggedright\popsemi\fontsize{8.4}{11}\selectfont\color{white}Gratuit \textbullet\ Google Play\\Tuteur IA, annales, concours, résultats\par}
      \end{minipage}
    \end{tcolorbox}
  \end{minipage}\hfill
  \begin{minipage}[t]{0.487\linewidth}
    \begin{tcolorbox}[enhanced,colback=poppurple,colframe=popink,boxrule=2.2pt,arc=10pt,
        drop shadow=popink,left=11pt,right=11pt,top=10pt,bottom=10pt,height=3.1cm,valign=center]
      \tikz[baseline=-0.7ex]\node[circle,fill=white,draw=popink,line width=1.3pt,minimum size=1.6cm,inner sep=0]{\popdisplay\fontsize{24}{24}\selectfont\color{poppurple}+};%
      \hspace{8pt}\begin{minipage}[c]{0.6\linewidth}
        {\popdisplay\fontsize{11}{12.5}\selectfont\color{white}\MakeUppercase{Passe à OnBuch+}}\par\vspace{4pt}
        {\raggedright\popsemi\fontsize{8.4}{11}\selectfont\color{white}Tous les cours signés, Tuteur IA illimité, fiches PDF — onglet OnBuch+ de l'app\par}
      \end{minipage}
    \end{tcolorbox}
  \end{minipage}
  \vspace{8pt}
  \popcontenu
  \vfill
  \signatureonbuch
  \restoregeometry
  \newpage
}

% Rangée « Ce cours contient » (page de garde)
\newcommand{\poptile}[3]{% #1 couleur #2 gros texte #3 légende
  \begin{tcolorbox}[enhanced,colback=#1,colframe=popink,boxrule=2pt,arc=9pt,shadow={2.5pt}{-2.5pt}{0pt}{popink},
      left=6pt,right=6pt,top=9pt,bottom=9pt,halign=center,valign=center,height=1.75cm,width=\linewidth,nobeforeafter]
    {\popdisplay\fontsize{12.5}{13}\selectfont\color{white}\MakeUppercase{#2}}\par\vspace{3pt}
    {\centering\popsemi\fontsize{7.8}{9.5}\selectfont\color{white}#3\par}
  \end{tcolorbox}}
\newcommand{\popcontenu}{%
  \par\noindent{\popxboldls\fontsize{7.5}{7.5}\selectfont\color{popmuted}CE COURS SIGNÉ CONTIENT}\par\vspace{7pt}
  \noindent\begin{minipage}[t]{0.235\linewidth}\poptile{popblue}{Cours}{complet, illustré, pas à pas}\end{minipage}\hfill
  \begin{minipage}[t]{0.235\linewidth}\poptile{poporange}{Méthodes}{et exercices résolus}\end{minipage}\hfill
  \begin{minipage}[t]{0.235\linewidth}\poptile{poppink}{Exercices}{type examen + corrigés}\end{minipage}\hfill
  \begin{minipage}[t]{0.235\linewidth}\poptile{popgreen}{Fiche bilan}{l'essentiel à retenir}\end{minipage}\par}

% Sommaire
\newtcolorbox{sommaire}{enhanced,colback=white,colframe=popink,boxrule=2.2pt,arc=10pt,
  drop shadow=popink,left=18pt,right=18pt,top=13pt,bottom=13pt,before skip=6pt,after skip=20pt,
  before upper={\poppill{Sommaire}{poppurple}{white}\par\vspace{9pt}\popsemi\fontsize{10.6}{14.5}\selectfont\color{popink}}}

% Signature de fin de cours — poussée tout en bas de la dernière page
\newcommand{\findecours}{%
  \par\vfill\nopagebreak
  \begin{center}\popsquiggle{poporange}\end{center}\vspace{4pt}
  \signatureonbuch
}
\AtBeginDocument{\def\llbracket{\mathopen{[\mkern-3.5mu[}}\def\rrbracket{\mathclose{]\mkern-3.5mu]}}} % intervalles d'entiers (glyphe ⟦ absent de la police)
% Glyphes absents de Fira Math : redéfinis à partir de glyphes disponibles
\AtBeginDocument{%
  \def\setminus{\mathbin{\backslash}}\let\smallsetminus\setminus
  \def\vdots{\mathord{\vbox{\baselineskip3.2pt\lineskiplimit0pt\kern4pt\hbox{.}\hbox{.}\hbox{.}}}}%
  \def\ddots{\mathinner{\mkern1mu\raise6.5pt\hbox{.}\mkern2mu\raise3.6pt\hbox{.}\mkern2mu\raise0.7pt\hbox{.}\mkern1mu}}%
  \def\triangle{\mathord{\scalebox{0.9}{$\symup{\Delta}$}}}%
  \def\nabla{\mathord{\raisebox{\depth}{\scalebox{1}[-1]{$\symup{\Delta}$}}}}%
  \def\checkmark{\popcheck{popdark}}%
  \def\star{\mathbin{\ast}}\def\bigstar{\mathord{\ast}}%
  \def\diamond{\mathbin{\scalebox{0.6}{$\Diamond$}}}%
  \def\bigcup{\mathop{\mathchoice{\raisebox{-0.3em}{\scalebox{1.7}{$\cup$}}}{\scalebox{1.25}{$\cup$}}{\cup}{\cup}}\displaylimits}%
  \def\bigcap{\mathop{\mathchoice{\raisebox{-0.3em}{\scalebox{1.7}{$\cap$}}}{\scalebox{1.25}{$\cap$}}{\cap}{\cap}}\displaylimits}%
  \def\lesseqqgtr{\mathrel{\lessgtr}}\def\bigoplus{\mathop{\oplus}\displaylimits}%
}
\providecommand{\ssi}{si et seulement si} % abréviation fréquente
\AtBeginDocument{\providecommand{\wideparen}{\overparen}} % arc de cercle

%% FIGFIX-BEGIN v5 — correctifs globaux des figures (docs/onbuch-generation/tools/figfix)
\usepackage{adjustbox}\usepackage{etoolbox}\tcbuselibrary{raster}
% toute figure TikZ/pgfplots plus large que la ligne est réduite à la largeur disponible
\BeforeBeginEnvironment{tikzpicture}{\begin{adjustbox}{max width=\linewidth}}
\AfterEndEnvironment{tikzpicture}{\end{adjustbox}}
% légendes pgfplots lisibles par-dessus les courbes
\pgfplotsset{every axis/.append style={legend style={fill=white,fill opacity=0.92,text opacity=1,draw=black!15,inner sep=3pt}}}
% étiquettes d'arêtes (syntaxe edge["..."]) avec fond blanc
\tikzset{every edge quotes/.append style={fill=white,inner sep=1.2pt}}
%% FIGFIX-END
\begin{document}

\pagegarde{Les maladies et l'importance de la pratique du sport}{Allemand}{1ère A}{Chapitre 3 : Environnement, santé et bien-être}{ALA12}{2 h}{Cette leçon propose l'étude d'un texte sur les maladies et l'importance de la pratique sportive. Elle permet d'acquérir le vocabulaire de la santé et du sport, de maîtriser les verbes pronominaux et l'expression du conseil avec « sollen » et « müssen ».
\vspace{4pt}
\begin{multicols}{2}\small
\begin{itemize}
\item À la fin de cette leçon, je sais comprendre globalement et en détail un texte allemand sur les maladies et le sport.
\item À la fin de cette leçon, je sais employer correctement le vocabulaire thématique de la santé, de la maladie et du sport (articles et pluriels compris).
\item À la fin de cette leçon, je sais reconnaître et conjuguer les verbes pronominaux au présent.
\item À la fin de cette leçon, je sais distinguer les verbes pronominaux à pronom réfléchi accusatif et datif.
\item À la fin de cette leçon, je sais exprimer le conseil et l'obligation avec « sollen » et « müssen ».
\item À la fin de cette leçon, je sais répondre par écrit et oralement à des questions de compréhension.
\end{itemize}\end{multicols}}

\begin{sommaire}
\begin{multicols}{2}
\begin{enumerate}
\item Texte support : « Sport ist die beste Medizin »
\item Compréhension du texte
\item Lexique thématique : santé, maladie et sport
\item Grammaire : les verbes pronominaux
\item Grammaire : le conseil avec « sollen » et « müssen »
\item Méthode du commentaire et production guidée
\item Activité d'intégration
\item Méthodes et exercices résolus
\item Exercices
\item Corrigés des exercices
\item Fiche bilan
\end{enumerate}
\end{multicols}
\end{sommaire}

\begin{objectifs}
\begin{itemize}
\item À la fin de cette leçon, je sais comprendre globalement et en détail un texte allemand sur les maladies et le sport.
\item À la fin de cette leçon, je sais employer correctement le vocabulaire thématique de la santé, de la maladie et du sport (articles et pluriels compris).
\item À la fin de cette leçon, je sais reconnaître et conjuguer les verbes pronominaux au présent.
\item À la fin de cette leçon, je sais distinguer les verbes pronominaux à pronom réfléchi accusatif et datif.
\item À la fin de cette leçon, je sais exprimer le conseil et l'obligation avec « sollen » et « müssen ».
\item À la fin de cette leçon, je sais répondre par écrit et oralement à des questions de compréhension.
\item À la fin de cette leçon, je sais rédiger un court paragraphe de conseils de santé en allemand.
\item À la fin de cette leçon, je sais commenter un texte selon la méthode du commentaire de texte (introduction, développement, conclusion).
\end{itemize}
\end{objectifs}

\begin{prerequis}
\begin{itemize}
\item La conjugaison des verbes réguliers et irréguliers au présent (Seconde)
\item Les articles définis et indéfinis au nominatif, accusatif et datif (Seconde)
\item Les verbes modaux de base : können, wollen, dürfen (Seconde)
\item Le vocabulaire du corps humain et de la vie quotidienne (Seconde)
\item La place du verbe dans la phrase déclarative et la phrase avec modal (début d'année de Première)
\end{itemize}
\end{prerequis}

\newpage

\coursec{Texte support : « Sport ist die beste Medizin »}

Pendant la saison des pluies à Yaoundé, le paludisme et le rhume frappent de nombreux élèves : certains manquent les cours, d'autres se rendent chez le médecin. En Allemagne, en Autriche et en Suisse, les médecins conseillent souvent la même chose : \all{„Sie sollen mehr Sport treiben !“} (« Vous devriez faire plus de sport ! »). Le sport est d'ailleurs remboursé par certaines caisses d'assurance maladie allemandes. Comment parler de la maladie, de la santé et donner des conseils en allemand ? C'est l'objectif de cette leçon.

\begin{definition}[Le texte support]
Un \cle{texte support} est un document authentique ou adapté qui sert de base à la compréhension écrite et orale et à l'étude de la langue (vocabulaire, grammaire, conjugaison). Il est le point de départ de toute la leçon : on l'observe, on le comprend, puis on en tire les règles de langue.
\end{definition}

\courssub{Lecture globale du texte}

\textbf{Première étape : l'écoute.}\par
Le professeur lit le texte à voix haute, lentement, deux fois. Les élèves écoutent \emph{sans lire}, le livre fermé : l'objectif est de saisir le sens global uniquement par l'oreille. À la deuxième écoute, chacun note sur son cahier les mots qu'il reconnaît.

\textbf{Deuxième étape : la lecture individuelle silencieuse.}\par
Chaque élève lit ensuite le texte seul, en silence, en s'aidant du glossaire. Il est interdit de traduire mot à mot : on cherche d'abord à répondre à trois questions simples : \emph{De qui parle-t-on ? Quel est son problème ? Comment se termine l'histoire ?}

\begin{methode}[Comment aborder un texte allemand]
\begin{enumerate}
\item Lire d'abord le \cle{titre} et le \cle{glossaire} : ils annoncent le sujet.
\item Repérer les \cle{mots transparents} (mots proches du français ou de l'anglais) : ils donnent des indices de sens sans dictionnaire.
\item Lire le texte une première fois pour le \cle{sens global} (qui ? quoi ? où ?).
\item Lire une deuxième fois pour les \cle{détails} : souligner les informations clés.
\item Formuler une hypothèse de sens, puis vérifier avec les questions de compréhension.
\end{enumerate}
\end{methode}

\courssub{Le texte : « Sport ist die beste Medizin »}

\begin{textebox}[Sport ist die beste Medizin]
Der junge Paul aus Yaoundé ist oft krank. Er hat Fieber und Husten und fühlt sich müde. Fast jeden Monat hat er eine Erkältung, und in der Schule kann er sich nicht gut konzentrieren. Seine Mutter sagt : „Du musst zum Arzt gehen !“

Paul geht also zum Arzt in ein Krankenhaus in der Stadt. Der Arzt untersucht ihn lange : er hört das Herz ab, schaut in den Hals und misst die Temperatur. Dann lächelt er und sagt : „Sie sind nicht wirklich krank, Herr Mbarga. Aber Sie sind schwach, weil Sie sich nicht genug bewegen. Sie müssen mehr Sport treiben. Sie sollen jeden Tag eine halbe Stunde laufen oder Fußball spielen. Und Sie sollen früh schlafen gehen und viel Wasser trinken.“

Paul folgt dem Rat des Arztes. Am Morgen läuft er jetzt dreißig Minuten vor der Schule. Er trainiert regelmäßig, schwimmt am Wochenende im Schwimmbad und spielt am Nachmittag mit seinen Freunden Fußball auf dem Sportplatz. Am Abend geht er früh ins Bett und trinkt keinen Kaffee mehr.

Nach drei Monaten fühlt er sich fit und gesund. Er ist selten krank, hat kein Fieber mehr und kann sich in der Schule besser konzentrieren. Seine Noten sind jetzt viel besser, und seine Mutter ist sehr zufrieden. Paul sagt zu seinen Freunden : „Sport ist wirklich die beste Medizin !“
\end{textebox}

\textbf{Glossaire (mots nouveaux) :}

\begin{center}
\begin{tabularx}{\linewidth}{|X|X|}
\hline
\rowcolor{popblueL} \textbf{Deutsch} & \textbf{Français} \\
\hline
die Krankheit, -en & la maladie \\
\hline
sich erkälten (die Erkältung, -en) & s'enrhumer (le rhume) \\
\hline
der Husten (Sg.) & la toux \\
\hline
müde & fatigué \\
\hline
der Arzt, die Ärzte / untersuchen & le médecin / examiner (ausculter) \\
\hline
das Krankenhaus, die Krankenhäuser & l'hôpital \\
\hline
schwach (contraire : stark) & faible (fort) \\
\hline
sich bewegen & bouger, faire de l'exercice \\
\hline
der Rat, die Ratschläge & le conseil \\
\hline
regelmäßig & régulièrement \\
\hline
die Gesundheit (Sg.) & la santé \\
\hline
zufrieden & content, satisfait \\
\hline
\end{tabularx}
\end{center}

\courssub{Les mots transparents : de précieux indices}

Avant même de traduire, certains mots du texte se comprennent presque tout seuls, car ils ressemblent au français ou à l'anglais. Ce sont les \cle{mots transparents}.

\begin{center}
\begin{tabularx}{\linewidth}{|X|X|X|}
\hline
\rowcolor{popblueL} \textbf{Deutsch} & \textbf{Français} & \textbf{Remarque} \\
\hline
die Medizin & la médecine & ressemblance totale \\
\hline
der Sport & le sport & identique \\
\hline
der Doktor & le docteur & ressemblance totale \\
\hline
trainieren & s'entraîner & proche de l'anglais \emph{to train} \\
\hline
das Fieber & la fièvre & attention au genre : \all{das}, pas \all{die} \\
\hline
die Temperatur & la température & ressemblance totale \\
\hline
sich konzentrieren & se concentrer & verbe pronominal \\
\hline
\end{tabularx}
\end{center}

\begin{attention}[Faux amis et pièges de genre]
\begin{itemize}
\item \all{das Fieber} est \emph{neutre} en allemand, alors que « la fièvre » est féminin en français. Toujours apprendre le nom \textbf{avec son article}.
\item \all{der Husten} signifie « la toux », pas « l'haleine ».
\item \all{der Rat} signifie « le conseil », pas « le rat » (le rat se dit \all{die Ratte}).
\item \all{bekommen} signifie « recevoir », pas « devenir » (devenir se dit \all{werden}).
\end{itemize}
\end{attention}

\courssub{Première hypothèse de sens}

À partir du titre \all{„Sport ist die beste Medizin“} (« Le sport est le meilleur médicament ») et des mots transparents (\all{der Sport}, \all{die Medizin}, \all{der Arzt}, \all{das Fieber}, \all{trainieren}), on peut formuler une hypothèse avant même la lecture détaillée :

\begin{exemplebox}[Hypothèse de sens]
Le texte parle probablement d'une personne malade qui consulte un médecin, et le sport joue le rôle d'un remède. Le titre suggère une idée générale : faire du sport, c'est se soigner.
\end{exemplebox}

Cette hypothèse sera vérifiée par la lecture détaillée. C'est une démarche essentielle en compréhension de l'écrit : on \emph{prédit}, puis on \emph{confirme} ou on \emph{corrige}.

\courssub{Lecture détaillée paragraphe par paragraphe}

On relit maintenant le texte en soulignant les informations clés : \textbf{qui ? où ? problème ? solution ? résultat ?}

\textbf{Paragraphe 1 — la présentation du problème.}\par
\all{Der junge Paul aus Yaoundé ist oft krank. Er hat Fieber und Husten und fühlt sich müde.} On apprend \emph{qui} (Paul, un jeune de Yaoundé), \emph{où} (à Yaoundé, au Cameroun) et \emph{quel est le problème} (il est souvent malade : fièvre, toux, fatigue, rhumes fréquents). À souligner : \all{oft krank}, \all{Fieber und Husten}, \all{müde}.

\textbf{Paragraphe 2 — la consultation et le conseil.}\par
Paul va à l'hôpital. Le médecin l'examine et prononce la phrase centrale du texte : \all{„Sie müssen mehr Sport treiben. Sie sollen jeden Tag eine halbe Stunde laufen oder Fußball spielen.“} C'est la \emph{solution} proposée. À souligner : \all{untersucht ihn}, \all{nicht wirklich krank}, \all{müssen}, \all{sollen}. On remarque déjà les deux verbes du conseil, \all{müssen} et \all{sollen}, qui seront étudiés en grammaire.

\textbf{Paragraphe 3 — l'application du conseil.}\par
Paul suit le conseil du médecin : \all{Paul folgt dem Rat des Arztes.} On relève ses nouvelles habitudes : courir le matin, nager le week-end, jouer au football, se coucher tôt. À souligner : \all{trainiert regelmäßig}, \all{schwimmt am Wochenende}, \all{spielt Fußball}.

\textbf{Paragraphe 4 — le résultat.}\par
\all{Nach drei Monaten fühlt er sich fit und gesund.} Après trois mois, Paul est en forme, rarement malade, plus concentré, et ses notes s'améliorent. La phrase finale reprend le titre : \all{„Sport ist wirklich die beste Medizin !“} À souligner : \all{fit und gesund}, \all{selten krank}, \all{kann sich besser konzentrieren}.

\begin{popfigure}
\begin{tikzpicture}[mindmap, grow cyclic, every node/.style={concept}, concept color=popblue!50, text=popdark, root concept/.append style={concept color=poporange!60, font=\bfseries}, level 1/.append style={level distance=3.2cm, sibling angle=90}]
\node[align=center]{der Text\\ „Sport ist die\\ beste Medizin“}
  child[concept color=popgreen!50] { node[align=center]{Personen\\ Paul, seine\\ Mutter, der Arzt} }
  child[concept color=poppink!50] { node[align=center]{Problem\\ oft krank:\\ Fieber, Husten,\\ müde} }
  child[concept color=popgold!60] { node[align=center]{Rat\\ mehr Sport,\\ früh schlafen,\\ viel Wasser} }
  child[concept color=poppurple!40] { node[align=center]{Ergebnis\\ fit und gesund,\\ selten krank,\\ bessere Noten} };
\end{tikzpicture}
\legende{Carte mentale du texte : les quatre informations clés à retenir.}
\end{popfigure}

\courssub{Questions de compréhension globale}

\begin{exoresolu}[Vérifier sa compréhension globale]
\textbf{Énoncé.} Réponds en allemand par des phrases complètes :
\begin{enumerate}
\item Wer ist oft krank ?
\item Was sagt der Arzt zu Paul ?
\item Wie fühlt sich Paul nach drei Monaten ?
\end{enumerate}
\tcblower
\textbf{Solution détaillée.}
\begin{enumerate}
\item \all{Der junge Paul aus Yaoundé ist oft krank.} (« Le jeune Paul de Yaoundé est souvent malade. ») — On reprend la structure de la question : \all{wer} (qui) appelle un sujet au nominatif.
\item \all{Der Arzt sagt zu Paul : „Sie müssen mehr Sport treiben.“} (« Le médecin dit à Paul : “Vous devez faire plus de sport.” ») — Le verbe \all{sagen} se conjugue : \all{er sagt} ; \all{zu Paul} reprend le complément de la question.
\item \all{Nach drei Monaten fühlt er sich fit und gesund.} (« Après trois mois, il se sent en forme et en bonne santé. ») — Attention : \all{sich fühlen} est un verbe pronominal ; le pronom réfléchi \all{sich} se place juste après le verbe conjugué.
\end{enumerate}
\end{exoresolu}

\begin{attention}[Erreurs fréquentes des francophones à la lecture]
\begin{itemize}
\item Ne pas confondre \all{krank} (malade) et \all{der Krankenwagen} (l'ambulance) : la famille de mots est régulière, observe les racines : \all{krank} → \all{die Krankheit} (la maladie) → \all{das Krankenhaus} (l'hôpital).
\item \all{„Sie sollen...“} avec majuscule à \all{Sie} est la forme de politesse « vous » ; le médecin vouvoie Paul. Ne pas traduire par « ils ».
\item En allemand, le verbe conjugué est toujours en \textbf{deuxième position} : \all{Nach drei Monaten fühlt er sich fit.} Le sujet \all{er} passe après le verbe, contrairement au français (« Après trois mois, il se sent... »).
\end{itemize}
\end{attention}

\begin{savaistu}[Le sport remboursé en Allemagne]
En Allemagne, certaines caisses d'assurance maladie (\all{die Krankenkassen}) remboursent les cours de sport de prévention (gymnastique, natation, yoga), car le sport régulier réduit les coûts de santé. Les médecins allemands peuvent même prescrire de l'activité physique comme un médicament : c'est la \all{Bewegungstherapie}, la « thérapie par le mouvement ». Le titre de notre texte est donc une réalité du système de santé allemand !
\end{savaistu}

\begin{exercice}[À toi de jouer]
\begin{enumerate}
\item Retrouve dans le texte deux mots de la famille de \all{krank} et traduis-les.
\item Souligne dans le texte toutes les phrases qui contiennent un conseil du médecin.
\item Recopie la carte mentale et ajoute une branche « Orte » (les lieux) avec les lieux mentionnés dans le texte.
\end{enumerate}
\end{exercice}

\begin{corrige}[Exercice « À toi de jouer »]
\begin{enumerate}
\item \all{das Krankenhaus} (l'hôpital) et \all{krank} dans \all{nicht wirklich krank} (pas vraiment malade). Famille complète à retenir : \all{krank} (malade) → \all{die Krankheit} (la maladie) → \all{das Krankenhaus} (l'hôpital) → \all{der Kranke} (le malade).
\item \all{„Sie müssen mehr Sport treiben. Sie sollen jeden Tag eine halbe Stunde laufen oder Fußball spielen. Und Sie sollen früh schlafen gehen und viel Wasser trinken.“}
\item Orte : \all{Yaoundé}, \all{das Krankenhaus}, \all{die Schule}, \all{das Schwimmbad}, \all{der Sportplatz}.
\end{enumerate}
\end{corrige}

\begin{aretenir}[L'essentiel de cette étape]
\begin{itemize}
\item Un texte se lit en deux temps : sens global d'abord, détails ensuite.
\item Le titre et les mots transparents permettent de formuler une hypothèse de sens avant la lecture.
\item La lecture détaillée repose sur cinq questions : \emph{qui ? où ? problème ? solution ? résultat ?}
\item Le texte « Sport ist die beste Medizin » raconte comment Paul, souvent malade, retrouve la santé grâce au sport régulier conseillé par son médecin.
\end{itemize}
\end{aretenir}

\coursec{Compréhension du texte}

Dans la section précédente, nous avons lu et écouté le texte support \all{„Sport ist die beste Medizin“}. Nous avons découvert l'histoire de Paul, un élève camerounais souvent fatigué, qui consulte un médecin et change ses habitudes. Il ne suffit pas de lire un texte : il faut le \cle{comprendre}, c'est-à-dire en saisir le sens global, retrouver les informations précises et savoir les reformuler. C'est exactement l'exercice demandé au baccalauréat. Cette section vous apprend, étape par étape, à répondre aux questions de compréhension.

\courssub{Rappel du texte support}

Relisons le texte avant de répondre aux questions. Repérez en lisant les mots-clés : \all{krank}, \all{Arzt}, \all{Sport}, \all{gesund}.

\begin{textebox}[„Sport ist die beste Medizin“]
Paul ist sechzehn Jahre alt und wohnt in Yaoundé. Er geht in die erste Klasse eines Gymnasiums. Seit einigen Monaten fühlt er sich oft müde und krank. Er hat manchmal Fieber und Kopfschmerzen, und er kann sich in der Schule nicht gut konzentrieren. Seine Noten werden schlechter, und seine Eltern machen sich Sorgen.

Eines Tages geht Paul mit seiner Mutter zum Arzt. Der Arzt untersucht ihn sorgfältig und stellt viele Fragen: „Schläfst du genug? Isst du gesund? Machst du Sport?“ Paul antwortet: „Ich schlafe wenig, ich esse oft Fastfood, und Sport? Nein, ich habe keine Zeit dafür. Ich sitze den ganzen Tag in der Schule oder vor dem Handy.“

Der Arzt lächelt und sagt: „Paul, du bist nicht wirklich krank, aber dein Körper braucht Bewegung. Du sollst jeden Tag mindestens dreißig Minuten Sport machen. Du musst auch mehr schlafen und besser essen: viel Obst und Gemüse, wenig Fastfood. Sport ist die beste Medizin!“

Paul folgt dem Rat des Arztes. Jetzt spielt er dreimal pro Woche Fußball mit seinen Freunden auf dem Sportplatz der Schule. Am Wochenende macht er einen langen Spaziergang mit seinem Vater. Er schläft jetzt acht Stunden pro Nacht und isst jeden Morgen Früchte.

Nach drei Monaten fühlt sich Paul viel besser. Er hat kein Fieber mehr, er kann sich gut konzentrieren, und seine Noten sind wieder gut. „Der Arzt hatte Recht“, sagt er zu seiner Mutter. „Sport ist wirklich die beste Medizin!“
\end{textebox}

\textbf{Glossaire (mots difficiles) :}

\begin{center}
\begin{tabularx}{\linewidth}{|X|X|}\hline
\rowcolor{popblueL} \textbf{Deutsch} & \textbf{Français} \\ \hline
sich fühlen (fühlte sich, hat sich gefühlt) & se sentir \\ \hline
das Fieber, - & la fièvre \\ \hline
die Kopfschmerzen (Pl.) & les maux de tête \\ \hline
sich konzentrieren & se concentrer \\ \hline
sich Sorgen machen & se faire du souci \\ \hline
untersuchen (hat untersucht) & examiner, auscultater \\ \hline
die Bewegung, -en & le mouvement, l'activité physique \\ \hline
der Rat, die Ratschläge & le conseil \\ \hline
dem Rat folgen & suivre le conseil \\ \hline
der Spaziergang, die Spaziergänge & la promenade \\ \hline
Recht haben (hatte Recht, hat Recht gehabt) & avoir raison \\ \hline
\end{tabularx}
\end{center}

\courssub{Méthode : comment répondre à une question de compréhension}

\begin{methode}[Répondre à une question de compréhension en trois étapes]
\begin{enumerate}
\item \textbf{Repérer le mot-clé de la question.} Dans \all{„Warum geht Paul zum Arzt?“}, les mots-clés sont \all{warum} (pourquoi) et \all{Arzt}.
\item \textbf{Localiser la phrase correspondante dans le texte.} Cherchez le mot-clé ou un synonyme : ici, le deuxième paragraphe (\all{„Eines Tages geht Paul mit seiner Mutter zum Arzt“}) et le premier paragraphe qui explique la cause (\all{„Er fühlt sich oft müde und krank“}).
\item \textbf{Reformuler sans recopier mot à mot.} Construisez une phrase complète avec vos propres mots : \all{„Paul geht zum Arzt, weil er sich oft müde und krank fühlt.“}
\end{enumerate}
\end{methode}

\begin{propriete}[Les questions de compréhension au baccalauréat]
Au Bac, les questions de compréhension portent d'abord sur le \cle{sens global} du texte (2 à 3 questions : le thème, les personnages, le lieu), puis sur des \cle{détails précis} (4 à 6 questions : causes, actions, sentiments, chronologie). Les réponses doivent être des \textbf{phrases complètes}, avec sujet, verbe conjugué et complément — jamais un seul mot ni une simple recopie du texte.
\end{propriete}

\courssub{Compréhension globale du texte}

Avant les détails, saisissons l'ensemble. Trois questions suffisent pour vérifier la compréhension globale.

\begin{exemplebox}[Questions de compréhension globale]
\begin{enumerate}
\item \all{Worum geht es in dem Text?} (De quoi parle le texte ?)\\
\all{Der Text handelt von einem Schüler, Paul, der oft krank ist und dann dank Sport wieder gesund wird.} — « Le texte parle d'un élève, Paul, souvent malade, qui guérit grâce au sport. »
\item \all{Wer sind die Personen?} (Qui sont les personnages ?)\\
\all{Die Personen sind Paul, seine Mutter, sein Vater und der Arzt.} — « Les personnages sont Paul, sa mère, son père et le médecin. »
\item \all{Wo spielt die Handlung?} (Où se passe l'action ?)\\
\all{Die Handlung spielt in Yaoundé, in der Schule, beim Arzt und auf dem Sportplatz.} — « L'action se passe à Yaoundé, à l'école, chez le médecin et sur le terrain de sport. »
\end{enumerate}
\end{exemplebox}

\courssub{Compréhension détaillée : questions en allemand}

Passons maintenant aux informations précises. Voici les quatre questions classiques de cette leçon, avec leurs réponses rédigées en phrases complètes.

\begin{exoresolu}[Questions détaillées sur le texte]
\textbf{Énoncé.} Répondez en allemand par des phrases complètes :
\begin{enumerate}
\item \all{Warum geht Paul zum Arzt?}
\item \all{Was sagt der Arzt?}
\item \all{Was macht Paul jetzt?}
\item \all{Wie fühlt er sich nach drei Monaten?}
\end{enumerate}
\tcblower
\textbf{Solution détaillée.}
\begin{enumerate}
\item Mot-clé : \all{warum} + \all{Arzt}. Le premier paragraphe donne la cause. Réponse : \all{Paul geht zum Arzt, weil er sich oft müde und krank fühlt und weil er manchmal Fieber und Kopfschmerzen hat.} — « Paul va chez le médecin parce qu'il se sent souvent fatigué et malade et qu'il a parfois de la fièvre et des maux de tête. »
\item Mot-clé : \all{der Arzt sagt}. Troisième paragraphe. Réponse : \all{Der Arzt sagt, dass Paul jeden Tag mindestens dreißig Minuten Sport machen soll und dass er mehr schlafen und besser essen muss.} — « Le médecin dit que Paul doit faire au moins trente minutes de sport par jour et qu'il doit aussi dormir davantage et mieux manger. »
\item Mot-clé : \all{jetzt}. Quatrième paragraphe. Réponse : \all{Jetzt spielt Paul dreimal pro Woche Fußball mit seinen Freunden, er macht Spaziergänge mit seinem Vater, er schläft acht Stunden pro Nacht und isst jeden Morgen Früchte.} — « Maintenant, Paul joue au football trois fois par semaine avec ses amis, il fait des promenades avec son père, il dort huit heures par nuit et mange des fruits chaque matin. »
\item Mot-clé : \all{nach drei Monaten}. Dernier paragraphe. Réponse : \all{Nach drei Monaten fühlt er sich viel besser: Er hat kein Fieber mehr, er kann sich gut konzentrieren, und seine Noten sind wieder gut.} — « Après trois mois, il se sent beaucoup mieux : il n'a plus de fièvre, il peut bien se concentrer et ses notes sont de nouveau bonnes. »
\end{enumerate}
\end{exoresolu}

\courssub{Questions en français (pratique du Bac)}

Au baccalauréat, certaines questions sont posées en français et exigent une réponse en français. La règle reste la même : phrase complète, reformulation, fidélité au texte.

\begin{exemplebox}[Questions-réponses en français]
\begin{enumerate}
\item \emph{Pourquoi les parents de Paul s'inquiètent-ils ?}\\
Les parents de Paul s'inquiètent parce que ses notes deviennent mauvaises et qu'il est souvent fatigué et malade.
\item \emph{Quelles questions le médecin pose-t-il à Paul ?}\\
Le médecin lui demande s'il dort assez, s'il mange sainement et s'il fait du sport.
\item \emph{Quel est le mode de vie de Paul avant la visite chez le médecin ?}\\
Avant la visite, Paul dort peu, mange souvent de la restauration rapide et reste assis toute la journée à l'école ou devant son téléphone portable.
\item \emph{Que conclut Paul à la fin du texte ?}\\
Paul conclut que le médecin avait raison et que le sport est vraiment le meilleur des remèdes.
\end{enumerate}
\end{exemplebox}

\begin{attention}[Ne traduisez pas mot à mot !]
Le titre \all{„Sport ist die beste Medizin“} ne se traduit pas littéralement par « le sport est la meilleure médecine », qui sonne artificiel en français. La traduction naturelle est : « \textbf{Le sport est le meilleur des remèdes} ». De même, \all{„der Arzt hatte Recht“} se traduit par « le médecin avait raison », et non « avait droit ». Cherchez toujours l'équivalent français naturel.
\end{attention}

\courssub{Vrai ou faux : justifier par une citation du texte}

L'exercice \all{„Richtig oder falsch?“} est très fréquent. La justification par une \textbf{citation exacte} du texte (entre guillemets allemands) est obligatoire.

\begin{exoresolu}[Richtig oder falsch ?]
\textbf{Énoncé.} Dites si les phrases suivantes sont vraies (\all{richtig}) ou fausses (\all{falsch}) et justifiez par une citation du texte.
\begin{enumerate}
\item \all{Paul ist siebzehn Jahre alt.}
\item \all{Paul macht vor dem Arztbesuch viel Sport.}
\item \all{Der Arzt empfiehlt dreißig Minuten Sport pro Tag.}
\item \all{Nach drei Monaten hat Paul noch Fieber.}
\end{enumerate}
\tcblower
\textbf{Solution détaillée.}
\begin{enumerate}
\item \all{Falsch.} Le texte dit : \all{„Paul ist sechzehn Jahre alt.“} (seize ans, pas dix-sept).
\item \all{Falsch.} Paul déclare lui-même : \all{„Sport? Nein, ich habe keine Zeit dafür.“}
\item \all{Richtig.} Le médecin dit : \all{„Du sollst jeden Tag mindestens dreißig Minuten Sport machen.“}
\item \all{Falsch.} Le dernier paragraphe affirme : \all{„Er hat kein Fieber mehr.“} (il n'a plus de fièvre).
\end{enumerate}
\end{exoresolu}

\courssub{La progression du récit : du malade au sportif}

Comprendre un texte, c'est aussi comprendre sa \cle{structure}. L'histoire de Paul suit une progression logique en cinq étapes, de la maladie à la guérison.

\begin{popfigure}
\begin{center}
\begin{tikzpicture}[
  node distance=0.45cm,
  etape/.style={rectangle, rounded corners=3pt, draw=popblue, fill=popblueL, thick, align=center, minimum width=2.6cm, minimum height=0.9cm, font=\small},
  fleche/.style={-{Stealth[length=2.5mm]}, thick, popdark}
]
\node[etape, fill=poppinkL, draw=poppink, align=center] (krank){\all{krank}\\ Paul est malade};
\node[etape, right=of krank, align=center] (arzt){\all{Arztbesuch}\\ visite chez\\ le médecin};
\node[etape, right=of arzt, align=center] (rat){\all{Rat}\\ le conseil :\\ sport, sommeil};
\node[etape, right=of rat, align=center] (training){\all{Training}\\ football et\\ promenades};
\node[etape, fill=popgreenL, draw=popgreen, right=of training, align=center] (gesund){\all{gesund}\\ Paul est guéri};
\draw[fleche] (krank) -- (arzt);
\draw[fleche] (arzt) -- (rat);
\draw[fleche] (rat) -- (training);
\draw[fleche] (training) -- (gesund);
\end{tikzpicture}
\end{center}
\legende{Organigramme de la progression du récit : \all{krank} $\rightarrow$ \all{Arztbesuch} $\rightarrow$ \all{Rat} $\rightarrow$ \all{Training} $\rightarrow$ \all{gesund}}
\end{popfigure}

Cette progression est un excellent plan pour le résumé : chaque étape correspond à une phrase de votre reformulation.

\courssub{Du texte vers la grammaire : chasse aux verbes}

La compréhension prépare la grammaire. Cherchons dans le texte deux familles de verbes que nous étudierons dans les sections suivantes.

\textbf{1. Les verbes pronominaux.} Le texte en contient trois :

\begin{center}
\begin{tabularx}{\linewidth}{|X|X|X|}\hline
\rowcolor{popblueL} \textbf{Verbe pronominal} & \textbf{Phrase du texte} & \textbf{Traduction} \\ \hline
\all{sich fühlen} & \all{Er fühlt sich oft müde und krank.} & Il se sent souvent fatigué et malade. \\ \hline
\all{sich konzentrieren} & \all{Er kann sich in der Schule nicht gut konzentrieren.} & Il ne peut pas bien se concentrer à l'école. \\ \hline
\all{sich Sorgen machen} & \all{Seine Eltern machen sich Sorgen.} & Ses parents se font du souci. \\ \hline
\end{tabularx}
\end{center}

Observez : le pronom réfléchi \all{sich} change selon la personne (\all{ich fühle mich}, \all{du fühlst dich}, \all{er fühlt sich}) et se place juste après le verbe conjugué. Nous verrons la règle complète dans la section 4.

\textbf{2. Les verbes modaux du conseil.} Le médecin utilise deux modaux :

\begin{center}
\begin{tabularx}{\linewidth}{|X|X|X|}\hline
\rowcolor{popblueL} \textbf{Modal} & \textbf{Phrase du texte} & \textbf{Nuance} \\ \hline
\all{sollen} & \all{Du sollst jeden Tag mindestens dreißig Minuten Sport machen.} & conseil, recommandation (« tu devrais ») \\ \hline
\all{müssen} & \all{Du musst auch mehr schlafen und besser essen.} & nécessité, obligation (« tu dois ») \\ \hline
\end{tabularx}
\end{center}

Observez : le modal est conjugué à la deuxième place de la phrase, et l'infinitif (\all{machen}, \all{schlafen}, \all{essen}) va \textbf{à la fin}. Nous étudierons cette construction dans la section 5.

\courssub{Exemples clés traduits}

\begin{exemplebox}[Deux phrases modèles à retenir]
\begin{itemize}
\item \all{Er fühlt sich müde.} — « Il se sent fatigué. » (verbe pronominal \all{sich fühlen} : le pronom \all{sich} suit le verbe conjugué.)
\item \all{Er folgt dem Rat des Arztes.} — « Il suit le conseil du médecin. » (Remarquez le datif : \all{dem Rat}, \all{des Arztes} au génitif.)
\end{itemize}
\end{exemplebox}

\courssub{Reformulation orale : le résumé guidé en trois phrases}

Dernier exercice de compréhension : résumer le texte oralement en \textbf{trois phrases}, une par grande étape du récit. Utilisez l'organigramme comme plan et les amorces suivantes.

\begin{methode}[Résumer un texte en trois phrases]
\begin{enumerate}
\item \textbf{Phrase 1 — le problème} : \all{Paul ist oft müde und krank, deshalb geht er mit seiner Mutter zum Arzt.}
\item \textbf{Phrase 2 — la solution} : \all{Der Arzt rät ihm, jeden Tag Sport zu machen, mehr zu schlafen und besser zu essen.}
\item \textbf{Phrase 3 — le résultat} : \all{Nach drei Monaten fühlt sich Paul viel besser, und seine Noten sind wieder gut.}
\end{enumerate}
Entraînez-vous à dire ces trois phrases à voix haute, sans regarder le texte, puis avec vos propres mots.
\end{methode}

\begin{exercice}[À vous de jouer]
\begin{enumerate}
\item Répondez en allemand par des phrases complètes : \all{Wo wohnt Paul? Was isst Paul vor dem Arztbesuch? Mit wem macht er am Wochenende einen Spaziergang?}
\item \all{Richtig oder falsch?} Justifiez par une citation : \all{a) Paul schläft jetzt acht Stunden pro Nacht. b) Der Arzt sagt, Paul ist wirklich krank.}
\item Résumez le texte en trois phrases avec vos propres mots.
\end{enumerate}
\end{exercice}

\begin{corrige}[Exercice « À vous de jouer »]
\begin{enumerate}
\item \all{Paul wohnt in Yaoundé. Vor dem Arztbesuch isst er oft Fastfood. Am Wochenende macht er einen Spaziergang mit seinem Vater.}
\item a) \all{Richtig} : \all{„Er schläft jetzt acht Stunden pro Nacht.“} b) \all{Falsch} : \all{„Paul, du bist nicht wirklich krank, aber dein Körper braucht Bewegung.“}
\item Exemple : \all{Paul ist oft krank und geht zum Arzt. Der Arzt empfiehlt Sport, gesundes Essen und mehr Schlaf. Nach drei Monaten ist Paul wieder gesund und hat gute Noten.}
\end{enumerate}
\end{corrige}

\begin{aretenir}[L'essentiel de la section]
\begin{itemize}
\item Compréhension globale d'abord (thème, personnages, lieu), compréhension détaillée ensuite (causes, actions, résultats).
\item Méthode en trois étapes : mot-clé de la question $\rightarrow$ phrase du texte $\rightarrow$ reformulation en phrase complète.
\item Au vrai/faux, la justification par une citation exacte est obligatoire.
\item Le texte suit la progression : \all{krank} $\rightarrow$ \all{Arztbesuch} $\rightarrow$ \all{Rat} $\rightarrow$ \all{Training} $\rightarrow$ \all{gesund}.
\item Le texte contient des verbes pronominaux (\all{sich fühlen}, \all{sich konzentrieren}) et des modaux du conseil (\all{sollen}, \all{müssen}) : ce sont les deux points de grammaire des sections suivantes.
\end{itemize}
\end{aretenir}

\coursec{Lexique thématique : santé, maladie et sport}

Dans la section précédente, nous avons lu et compris le texte \all{„Sport ist die beste Medizin“}. Pour bien parler de ce thème — et pour pouvoir résumer le texte, répondre aux questions et produire nos propres phrases —, il faut maintenant organiser et mémoriser son vocabulaire. C'est l'objet de cette section : ranger les mots par \cle{champs lexicaux}, apprendre chaque nom avec son \cle{article} et son \cle{pluriel}, et découvrir les \cle{familles de mots}.

\begin{definition}[Champ lexical et famille de mots]
Un \cle{champ lexical} regroupe les mots qui appartiennent à un même thème : par exemple, \all{der Arzt}, \all{das Fieber} et \all{das Krankenhaus} appartiennent au champ lexical de la maladie. Une \cle{famille de mots} regroupe les mots construits sur la même racine : \all{krank} (malade) → \all{die Krankheit} (la maladie) → \all{das Krankenhaus} (l'hôpital, litt. « maison des malades »). Apprendre le vocabulaire par champs et par familles est la méthode la plus efficace pour retenir durablement.
\end{definition}

\courssub{Observation : le vocabulaire dans un texte}

Avant d'étudier les listes, lisons ce court texte original. Soulignez mentalement tous les mots qui parlent de la santé, de la maladie ou du sport.

\begin{textebox}[„Paul ist krank“ — texte de vocabulaire]
Paul ist Schüler in Yaoundé. Er ist sechzehn Jahre alt und normalerweise sehr fit. Aber seit drei Tagen ist er krank. Er hat Fieber und Husten, und sein Kopf tut weh. Er ist sehr müde und kann nicht zur Schule gehen.

Seine Mutter sagt: „Du musst zum Arzt gehen!“ Am Morgen gehen sie zusammen ins Krankenhaus. Die Ärztin untersucht Paul. Sie sagt: „Sie haben eine Erkältung, aber keine Grippe. Das ist nicht gefährlich.“ Sie gibt Paul ein Medikament und einen guten Rat: „Trinken Sie viel Wasser, schlafen Sie viel und bleiben Sie drei Tage zu Hause.“

Pauls Freunde besuchen ihn am Nachmittag. Sie sagen: „Gute Besserung!“ Paul lächelt. Nach einer Woche ist er wieder gesund. Jetzt treibt er wieder Sport: Er spielt Fußball mit seiner Mannschaft und trainiert dreimal pro Woche. „Sport ist wichtig für die Gesundheit“, sagt er. „Wer sich bewegt, bleibt fit!“

\textbf{Glossaire :} das Fieber (–) : la fièvre ; der Husten (–) : la toux ; untersuchen : examiner ; die Erkältung (-en) : le rhume ; die Grippe (-n) : la grippe ; der Rat (Ratschläge) : le conseil ; die Mannschaft (-en) : l'équipe ; sich bewegen : bouger, faire de l'exercice.
\end{textebox}

\textbf{Questions de repérage :}
\begin{enumerate}
\item Welche Krankheit hat Paul? (Quelle maladie Paul a-t-il ?)
\item Wer untersucht Paul? (Qui examine Paul ?)
\item Was macht Paul nach einer Woche? (Que fait Paul après une semaine ?)
\end{enumerate}

\courssub{Le champ lexical de la maladie}

\begin{definition}[die Krankheit]
\all{die Krankheit} (pluriel : \all{die Krankheiten}) signifie « la maladie ». Prononciation : « krank-haït ». Le suffixe \all{-heit} forme des noms féminins à partir d'adjectifs : \all{krank} → \all{die Krankheit}.
\end{definition}

\begin{center}
\begin{tabularx}{\linewidth}{|X|X|X|}\hline
\rowcolor{popblueL} Singulier & Pluriel & Traduction \\ \hline
die Krankheit & die Krankheiten & la maladie \\ \hline
das Fieber & (–) & la fièvre \\ \hline
der Husten & (–) & la toux \\ \hline
die Erkältung & die Erkältungen & le rhume \\ \hline
die Grippe & die Grippen & la grippe \\ \hline
der Schmerz & die Schmerzen & la douleur \\ \hline
die Müdigkeit & (–) & la fatigue \\ \hline
\end{tabularx}
\end{center}

\begin{aretenir}[Remarque sur les pluriels]
\all{das Fieber}, \all{der Husten} et \all{die Müdigkeit} s'emploient presque toujours au singulier (comme « la fièvre », « la toux », « la fatigue » en français). Retenez surtout : \all{die Krankheiten} (les maladies) et \all{die Schmerzen} (les douleurs). Notez la construction utile : \all{Kopfschmerzen haben} = avoir mal à la tête ; \all{Bauchschmerzen haben} = avoir mal au ventre. \all{Kopfschmerzen} et \all{Bauchschmerzen} sont des \cle{pluralia tantum} (noms qui n'existent qu'au pluriel).
\end{aretenir}

\begin{exemplebox}[Exemples traduits]
\all{Er hat Fieber.} « Il a de la fièvre. »

\all{Sie hat eine Erkältung und Husten.} « Elle a un rhume et de la toux. »

\all{Ich habe Kopfschmerzen.} « J'ai mal à la tête. »

\all{Die Grippe ist gefährlich.} « La grippe est dangereuse. »
\end{exemplebox}

\courssub{Le champ lexical de la santé et du corps}

\begin{center}
\begin{tabularx}{\linewidth}{|X|X|X|}\hline
\rowcolor{popblueL} Singulier & Pluriel & Traduction \\ \hline
die Gesundheit & (–) & la santé \\ \hline
der Körper & die Körper & le corps \\ \hline
der Kopf & die Köpfe & la tête \\ \hline
der Bauch & die Bäuche & le ventre \\ \hline
\end{tabularx}
\end{center}

Les adjectifs essentiels : \all{gesund} (en bonne santé), \all{krank} (malade), \all{fit} (en forme), \all{müde} (fatigué).

\begin{exemplebox}[Exemples traduits]
\all{Sport ist gut für die Gesundheit.} « Le sport est bon pour la santé. » (Accusatif après \all{für})

\all{Nach dem Training ist er müde, aber fit.} « Après l'entraînement, il est fatigué mais en forme. »

\all{Mein Körper braucht Bewegung.} « Mon corps a besoin de mouvement. » (Accusatif)
\end{exemplebox}

\begin{attention}[« Gesundheit ! »]
En Allemagne, quand quelqu'un éternue, on dit \all{„Gesundheit!“} (« Santé ! »), comme notre « À tes souhaits ! ». Ne confondez pas : \all{die Gesundheit} est un nom abstrait (la santé), il n'a pas de pluriel.
\end{attention}

\courssub{Le champ lexical du médecin et du traitement}

\begin{center}
\begin{tabularx}{\linewidth}{|X|X|X|}\hline
\rowcolor{popblueL} Singulier & Pluriel & Traduction \\ \hline
der Arzt & die Ärzte & le médecin (homme) \\ \hline
die Ärztin & die Ärztinnen & le médecin (femme) \\ \hline
das Krankenhaus & die Krankenhäuser & l'hôpital \\ \hline
die Medizin & (die Medizinen) & la médecine (science) \\ \hline
das Medikament & die Medikamente & le médicament \\ \hline
der Rat & die Ratschläge & le conseil \\ \hline
\end{tabularx}
\end{center}

Verbe associé : \all{untersuchen} (examiner, transitif) : \all{Der Arzt untersucht den Patienten.} « Le médecin examine le patient. » (Accusatif \all{den Patienten})

\begin{attention}[Faux ami : die Medizin / das Medikament]
\all{die Medizin} désigne la \textbf{médecine} (la discipline, la science, la faculté) : \all{Er studiert Medizin.} « Il étudie la médecine. » Pour « un médicament » (substance), on dit \all{das Medikament} ou \all{das Arzneimittel} : \all{Die Ärztin gibt ihm ein Medikament.} « Le médecin lui donne un médicament. » Ne dites donc jamais \all{„eine Medizin nehmen“} pour « prendre un médicament », mais \all{ein Medikament nehmen}.
\end{attention}

\begin{exemplebox}[Exemples traduits]
\all{Sie geht zum Arzt.} « Elle va chez le médecin. » (\all{zum} = \all{zu dem}, Dativ)

\all{Die Ärztin untersucht den Jungen im Krankenhaus.} « Le médecin examine le garçon à l'hôpital. » (\all{im} = \all{in dem}, Dativ lieu)

\all{Der Arzt gibt mir einen guten Rat.} « Le médecin me donne un bon conseil. » (Dativ \all{mir}, Accusatif \all{einen guten Rat})
\end{exemplebox}

\courssub{Le champ lexical du sport}

\begin{center}
\begin{tabularx}{\linewidth}{|X|X|X|}\hline
\rowcolor{popblueL} Singulier & Pluriel & Traduction \\ \hline
der Sport & (–) & le sport \\ \hline
der Fußball & die Fußbälle & le football / le ballon \\ \hline
das Training & die Trainings & l'entraînement \\ \hline
die Mannschaft & die Mannschaften & l'équipe \\ \hline
der Sportler & die Sportler & le sportif \\ \hline
\end{tabularx}
\end{center}

Verbes associés : \all{trainieren} (s'entraîner), \all{laufen} (courir), \all{schwimmen} (nager), \all{spielen} (jouer), \all{sich bewegen} (bouger, faire de l'exercice — verbe pronominal, voir section grammaire).

\begin{exemplebox}[Exemples traduits]
\all{Wir treiben jeden Tag Sport.} « Nous faisons du sport chaque jour. » (\all{treiben} + Accusatif)

\all{Die Mannschaft trainiert dreimal pro Woche.} « L'équipe s'entraîne trois fois par semaine. »

\all{Er läuft jeden Morgen vor der Schule.} « Il court chaque matin avant l'école. » (\all{vor} + Dativ)
\end{exemplebox}

\courssub{Les familles de mots}

Observer les familles permet de deviner le sens de mots nouveaux. Voici les trois familles du texte :

\begin{center}
\begin{tabularx}{\linewidth}{|X|X|X|}\hline
\rowcolor{popblueL} Racine & Famille & Traduction \\ \hline
krank (malade) & die Krankheit, das Krankenhaus & la maladie, l'hôpital \\ \hline
gesund (en bonne santé) & die Gesundheit, gesund bleiben & la santé, rester en bonne santé \\ \hline
der Sport (le sport) & der Sportler, sportlich & le sportif, sportif (adj.) \\ \hline
\end{tabularx}
\end{center}

\begin{propriete}[Les suffixes \all{-heit} et \all{-keit}]
Le suffixe \all{-heit} transforme un adjectif en nom \textbf{féminin} : \all{krank} → \all{die Krankheit} ; \all{gesund} → \all{die Gesundheit}. La variante \all{-keit} s'utilise après les suffixes adjectivaux \all{-ig}, \all{-lich}, \all{-sam}, \all{-bar}, \all{-haft} (ex. : \all{möglich} → \all{die Möglichkeit}, \all{freundlich} → \all{die Freundlichkeit}) et pour \all{müde} → \all{die Müdigkeit}. \textbf{Tous les noms en \all{-heit} et \all{-keit} sont féminins} : c'est un excellent réflexe pour l'article !
\end{propriete}

\courssub{Expressions utiles à mémoriser}

\begin{center}
\begin{tabularx}{\linewidth}{|X|X|}\hline
\rowcolor{popblueL} Deutsch & Français \\ \hline
krank sein & être malade \\ \hline
zum Arzt gehen & aller chez le médecin \\ \hline
Sport treiben & faire du sport \\ \hline
sich fit halten & se maintenir en forme \\ \hline
gesund bleiben & rester en bonne santé \\ \hline
Gute Besserung! & Bon rétablissement ! \\ \hline
\end{tabularx}
\end{center}

\begin{savaistu}[„Gute Besserung!“]
\all{„Gute Besserung!“} est la formule à dire à un malade en Allemagne, l'équivalent de « Bon rétablissement ! ». Les Allemands la disent aussi au téléphone ou l'écrivent sur une carte. Au Cameroun, on rend visite aux malades à l'hôpital ; en Allemagne aussi, et on n'oublie jamais cette petite phrase de politesse !
\end{savaistu}

\courssub{Carte mentale du vocabulaire}

\begin{popfigure}
\begin{tikzpicture}[
  mindmap,
  grow cyclic,
  every node/.style={concept, align=center, font=\small},
  concept color=popblue,
  level 1/.append style={level distance=3.2cm, sibling angle=90},
  root concept/.append style={concept color=popdark, font=\small\bfseries}
]
\node[root concept, align=center]{Gesundheit\\und Sport}
  child[concept color=poporange]{ node {Krankheit} 
    child[grow=180, concept color=poporange!70, level distance=2.4cm]{ node[align=center]{das Fieber\\der Husten\\die Grippe} } }
  child[concept color=popgreen]{ node {Körper}
    child[grow=270, concept color=popgreen!70, level distance=2.4cm]{ node[align=center]{der Kopf\\der Bauch\\fit, müde} } }
  child[concept color=poppurple]{ node {Arzt}
    child[grow=0, concept color=poppurple!70, level distance=2.4cm]{ node[align=center]{das Krankenhaus\\das Medikament\\untersuchen} } }
  child[concept color=poppink]{ node {Sport}
    child[grow=90, concept color=poppink!70, level distance=2.4cm]{ node[align=center]{trainieren\\die Mannschaft\\sich bewegen} } };
\end{tikzpicture}
\legende{Carte mentale : le vocabulaire de la santé et du sport en quatre champs lexicaux.}
\end{popfigure}

\courssub{Exercice résolu et entraînement}

\begin{exoresolu}[Complétez avec le bon mot]
Énoncé : Complétez avec : \all{Fieber}, \all{Arzt}, \all{Mannschaft}, \all{Gesundheit}, \all{Medikament}.

1. Paul hat 39 Grad \dots\ — er ist krank. 2. Die Ärztin gibt ihm ein \dots\ gegen die Erkältung. 3. Sport ist gut für die \dots. 4. Unsere \dots\ spielt heute Fußball. 5. Mein Vater ist \dots\ in Douala.
\tcblower
Solution : 1. \all{Paul hat 39 Grad \textbf{Fieber} — er ist krank.} (Paul a 39 de fièvre.) 2. \all{Die Ärztin gibt ihm ein \textbf{Medikament}.} (et non \all{Medizin} : un médicament = \all{das Medikament}). 3. \all{Sport ist gut für die \textbf{Gesundheit}.} 4. \all{Unsere \textbf{Mannschaft} spielt heute Fußball.} 5. \all{Mein Vater ist \textbf{Arzt} in Douala.} (sans article après \all{sein} pour la profession : \all{Er ist Arzt.})
\end{exoresolu}

\begin{attention}[Pièges des francophones]
\begin{itemize}
\item \all{bekommen} ne signifie pas « devenir » mais « \textbf{recevoir} » : \all{Er bekommt ein Medikament.} « Il reçoit un médicament. » « Devenir » se dit \all{werden} : \all{Er wird Arzt.} « Il devient médecin. »
\item N'oubliez pas la majuscule à tous les noms : \all{das Fieber}, \all{die Gesundheit}, \all{der Sport}.
\item Apprenez toujours le nom \textbf{avec son article} : on dit \all{der Husten} (masculin) mais \all{die Grippe} (féminin) et \all{das Fieber} (neutre).
\end{itemize}
\end{attention}

\begin{exercice}[À vous]
1. Traduisez : a) Elle va chez le médecin. b) Nous faisons du sport chaque jour. c) Il a de la fièvre et de la toux. 2. Formez la famille de mots de \all{krank} et de \all{gesund}. 3. Donnez l'article et le pluriel de : \all{Krankenhaus}, \all{Mannschaft}, \all{Medikament}, \all{Arzt}.
\end{exercice}

\begin{corrige}[Exercice « À vous »]
1. a) \all{Sie geht zum Arzt.} b) \all{Wir treiben jeden Tag Sport.} c) \all{Er hat Fieber und Husten.} 2. \all{krank} → \all{die Krankheit}, \all{das Krankenhaus} ; \all{gesund} → \all{die Gesundheit}, \all{gesund bleiben}. 3. \all{das Krankenhaus, die Krankenhäuser} ; \all{die Mannschaft, die Mannschaften} ; \all{das Medikament, die Medikamente} ; \all{der Arzt, die Ärzte}.
\end{corrige}

Ce vocabulaire est maintenant votre boîte à outils : gardez-la ouverte pour la suite de la leçon, car les verbes pronominaux comme \all{sich fit halten} et \all{sich bewegen} seront au cœur de la prochaine section de grammaire.

\coursec{Grammaire : les verbes pronominaux}

\courssub{De l'observation à la règle}

Dans la section précédente, nous avons appris le vocabulaire de la santé, de la maladie et du sport : \all{die Krankheit}, \all{der Sport}, \all{gesund}, \all{der Arzt}, \all{das Fieber}, \all{trainieren}. Reprenons maintenant deux phrases de notre texte support \all{Sport ist die beste Medizin} et observons-les de près :

\all{Er fühlt sich fit.} « Il se sent en forme. »

\all{Er konzentriert sich besser.} « Il se concentre mieux. »

Que remarque-t-on ? Le verbe (\all{fühlt}, \all{konzentriert}) est accompagné d'un petit mot : \all{sich}. Ce mot renvoie au sujet \all{er} : c'est lui-même qui se sent, lui-même qui se concentre. C'est exactement le mécanisme du français « il \emph{se} sent », « il \emph{se} concentre ». Ce petit mot s'appelle le \cle{pronom réfléchi} (\all{das Reflexivpronomen}), et les verbes qui l'accompagnent sont les \cle{verbes pronominaux} (\all{reflexive Verben}).

\begin{definition}[Le verbe pronominal]
Un \cle{verbe pronominal} (\all{reflexives Verb}) est un verbe qui se conjugue toujours avec un \cle{pronom réfléchi} (\all{sich} à l'infinitif). Ce pronom renvoie au sujet : l'action revient sur celui qui la fait. Dans le dictionnaire, on note ces verbes avec \all{sich} : \all{sich freuen}, \all{sich fühlen}, \all{sich erkälten}.
\end{definition}

\begin{propriete}[Place et accord du pronom réfléchi]
Deux règles d'or :
\begin{enumerate}
\item Le pronom réfléchi \textbf{s'accorde avec le sujet} : \all{ich} $\rightarrow$ \all{mich}, \all{du} $\rightarrow$ \all{dich}, \all{er/sie/es} $\rightarrow$ \all{sich}, etc.
\item Le pronom réfléchi se place \textbf{immédiatement après le verbe conjugué}, dans la phrase affirmative comme dans la question (phrase principale).
\end{enumerate}
\all{Ich freue mich.} « Je me réjouis. » \quad \all{Freust du dich ?} « Te réjouis-tu ? » \quad \all{Sie erholt sich schnell.} « Elle se rétablit vite. »
\end{propriete}

\courssub{Les pronoms réfléchis : Akkusativ et Dativ}

Le pronom réfléchi possède deux séries de formes, selon qu'il est à l'\textbf{accusatif} (cas le plus fréquent) ou au \textbf{datif}. Bonne nouvelle : seules les personnes \all{ich} et \all{du} changent (\all{mich/mir}, \all{dich/dir}) ; toutes les autres formes sont identiques.

\begin{center}
\begin{tabular}{|l|l|l|l|}
\hline
\rowcolor{popblueL} Personne & Akkusativ & Dativ & Français \\
\hline
ich & mich & mir & me \\
\hline
du & dich & dir & te \\
\hline
er / sie / es & sich & sich & se \\
\hline
wir & uns & uns & nous \\
\hline
ihr & euch & euch & vous \\
\hline
sie / Sie & sich & sich & se \\
\hline
\end{tabular}
\end{center}

La figure suivante résume visuellement ces deux séries et met en évidence les deux seules formes qui changent.

\begin{popfigure}
\begin{tikzpicture}
\node[draw=popblue, fill=popblueL, rounded corners, minimum width=5.2cm, minimum height=0.9cm, font=\bfseries] (t1) at (0,0) {Akkusativ (cas général)};
\node[draw=poporange, fill=poporangeL, rounded corners, minimum width=5.2cm, minimum height=0.9cm, font=\bfseries] (t2) at (6.2,0) {Dativ (cas particulier)};
\node[draw=popline, fill=white, rounded corners, minimum width=5.2cm, minimum height=0.7cm] at (0,-0.9) {ich $\rightarrow$ mich};
\node[draw=popline, fill=white, rounded corners, minimum width=5.2cm, minimum height=0.7cm] at (0,-1.7) {du $\rightarrow$ dich};
\node[draw=popline, fill=white, rounded corners, minimum width=5.2cm, minimum height=0.7cm] at (0,-2.5) {er/sie/es $\rightarrow$ sich};
\node[draw=popline, fill=white, rounded corners, minimum width=5.2cm, minimum height=0.7cm] at (0,-3.3) {wir $\rightarrow$ uns};
\node[draw=popline, fill=white, rounded corners, minimum width=5.2cm, minimum height=0.7cm] at (0,-4.1) {ihr $\rightarrow$ euch};
\node[draw=popline, fill=white, rounded corners, minimum width=5.2cm, minimum height=0.7cm] at (0,-4.9) {sie/Sie $\rightarrow$ sich};
\node[draw=poporange, fill=poppinkL, rounded corners, minimum width=5.2cm, minimum height=0.7cm] at (6.2,-0.9) {ich $\rightarrow$ \textbf{mir}};
\node[draw=poporange, fill=poppinkL, rounded corners, minimum width=5.2cm, minimum height=0.7cm] at (6.2,-1.7) {du $\rightarrow$ \textbf{dir}};
\node[draw=popline, fill=white, rounded corners, minimum width=5.2cm, minimum height=0.7cm] at (6.2,-2.5) {er/sie/es $\rightarrow$ sich};
\node[draw=popline, fill=white, rounded corners, minimum width=5.2cm, minimum height=0.7cm] at (6.2,-3.3) {wir $\rightarrow$ uns};
\node[draw=popline, fill=white, rounded corners, minimum width=5.2cm, minimum height=0.7cm] at (6.2,-4.1) {ihr $\rightarrow$ euch};
\node[draw=popline, fill=white, rounded corners, minimum width=5.2cm, minimum height=0.7cm] at (6.2,-4.9) {sie/Sie $\rightarrow$ sich};
\end{tikzpicture}
\legende{Les pronoms réfléchis : seuls \all{mir} et \all{dir} (en rose) diffèrent de l'accusatif.}
\end{popfigure}

\courssub{Conjugaison au présent des verbes du thème}

Voici la conjugaison complète au présent des verbes pronominaux essentiels de notre chapitre. Remarquez que le verbe se conjugue normalement ; seul le pronom réfléchi s'ajoute.

\begin{center}
\begin{tabular}{|l|l|l|}
\hline
\rowcolor{popblueL} sich freuen (se réjouir) & sich fühlen (se sentir) & sich erkälten (s'enrhumer) \\
\hline
ich freue mich & ich fühle mich & ich erkälte mich \\
\hline
du freust dich & du fühlst dich & du erkältest dich \\
\hline
er/sie/es freut sich & er/sie/es fühlt sich & er/sie/es erkältet sich \\
\hline
wir freuen uns & wir fühlen uns & wir erkälten uns \\
\hline
ihr freut euch & ihr fühlt euch & ihr erkältet euch \\
\hline
sie/Sie freuen sich & sie/Sie fühlen sich & sie/Sie erkälten sich \\
\hline
\end{tabular}
\end{center}

\begin{center}
\begin{tabular}{|l|l|l|}
\hline
\rowcolor{popgreenL} sich konzentrieren (se concentrer) & sich ausruhen (se reposer) & sich bewegen (bouger) \\
\hline
ich konzentriere mich & ich ruhe mich aus & ich bewege mich \\
\hline
du konzentrierst dich & du ruhst dich aus & du bewegst dich \\
\hline
er/sie/es konzentriert sich & er/sie/es ruht sich aus & er/sie/es bewegt sich \\
\hline
wir konzentrieren uns & wir ruhen uns aus & wir bewegen uns \\
\hline
ihr konzentriert euch & ihr ruht euch aus & ihr bewegt euch \\
\hline
sie/Sie konzentrieren sich & sie/Sie ruhen sich aus & sie/Sie bewegen sich \\
\hline
\end{tabular}
\end{center}

\begin{attention}[\all{sich ausruhen} : un verbe à particule séparable]
\all{sich ausruhen} est à la fois pronominal \emph{et} à particule séparable : la particule \all{aus-} va en fin de phrase. \all{Ich ruhe mich nach dem Training aus.} « Je me repose après l'entraînement. » Ne dites jamais \all{Ich ausruhe mich}.
\end{attention}

\begin{exemplebox}[Exemples traduits]
\begin{itemize}
\item \all{Ich erkälte mich oft.} « Je m'enrhume souvent. »
\item \all{Du musst dich ausruhen.} « Tu dois te reposer. »
\item \all{Wir bewegen uns viel.} « Nous bougeons beaucoup. »
\item \all{Nach dem Sport fühlt er sich fit.} « Après le sport, il se sent en forme. »
\item \all{Die Schüler konzentrieren sich im Unterricht besser.} « Les élèves se concentrent mieux en classe. »
\item \all{Verletzt du dich oft beim Fußball ?} « Te blesses-tu souvent au football ? »
\end{itemize}
\end{exemplebox}

\courssub{Accusatif ou datif ? Le cas du complément direct}

Règle générale : le pronom réfléchi est à l'\textbf{accusatif}. Mais il passe au \textbf{datif} quand le verbe possède déjà un autre complément d'objet direct (COD). Comparez :

\all{Ich wasche mich.} « Je me lave. » (pas d'autre COD $\rightarrow$ accusatif)

\all{Ich wasche mir die Hände.} « Je me lave les mains. » (COD : \all{die Hände} $\rightarrow$ datif)

\all{Du putzt dir die Zähne.} « Tu te brosses les dents. »

\all{Er bricht sich das Bein.} « Il se casse la jambe. »

\begin{methode}[Choisir entre accusatif et datif]
\begin{enumerate}
\item Repérer le verbe pronominal et son sujet.
\item Chercher dans la phrase un \textbf{autre complément d'objet direct} (souvent une partie du corps : \all{die Hände}, \all{die Zähne}, \all{das Bein}).
\item S'il y a un autre COD $\rightarrow$ le pronom réfléchi est au \textbf{datif} (\all{mir}, \all{dir}, \all{sich}...).
\item Sinon $\rightarrow$ le pronom réfléchi est à l'\textbf{accusatif} (\all{mich}, \all{dich}, \all{sich}...).
\end{enumerate}
\end{methode}

\courssub{Comparaison français / allemand}

Le français et l'allemand fonctionnent souvent de la même manière, mais pas toujours. Certains verbes sont pronominaux dans une langue et pas dans l'autre, et les prépositions associées diffèrent.

\begin{center}
\begin{tabularx}{\linewidth}{|X|X|X|}
\hline
\rowcolor{popblueL} Français & Deutsch & Remarque \\
\hline
se sentir & sich fühlen & parallèle parfait \\
\hline
se reposer & sich ausruhen & particule séparable en plus \\
\hline
se souvenir de & sich erinnern an & préposition \all{an} + Akk, pas \all{von} \\
\hline
s'intéresser à & sich interessieren für & préposition \all{für} + Akk \\
\hline
se réjouir & sich freuen & pronominal dans les deux langues \\
\hline
se dépêcher & sich beeilen & pronominal dans les deux langues \\
\hline
attraper un rhume & sich erkälten & pronominal en allemand seulement \\
\hline
se blesser & sich verletzen & parallèle parfait \\
\hline
\end{tabularx}
\end{center}

\begin{attention}[Verbes pronominaux en allemand mais pas en français]
Certains verbes sont obligatoirement pronominaux en allemand alors que le français n'a pas de « se » : \all{sich erkälten} « attraper un rhume », \all{sich erholen} « se rétablir », \all{sich fit halten} « se maintenir en forme ». Inversement, attention aux prépositions : \all{sich erinnern an} = « se souvenir \emph{de} ». La règle d'or : \textbf{apprenez chaque verbe avec sa construction complète} (pronom + préposition éventuelle + cas).
\end{attention}

\coursec{Grammaire : le conseil avec « sollen » et « müssen »}

\courssub{Conjugaison au présent}

Les verbes modaux \all{müssen} et \all{sollen} sont indispensables pour exprimer le conseil, l'obligation ou la recommandation. Ils sont irréguliers au présent (changement de voyelle pour \all{müssen} aux singulières).

\begin{center}
\begin{tabular}{|l|l|l|}
\hline
\rowcolor{popblueL} müssen (devoir / avoir à) & sollen (devoir / être censé) & Remarque \\
\hline
ich muss & ich soll & pas de \emph{-e} à la 1ère pers. sg. \\
\hline
du musst & du sollst & \\
\hline
er/sie/es muss & er/sie/es soll & pas de \emph{-t} à la 3ème pers. sg. \\
\hline
wir müssen & wir sollen & formes régulières au pluriel \\
\hline
ihr müsst & ihr sollt & \\
\hline
sie/Sie müssen & sie/Sie sollen & \\
\hline
\end{tabular}
\end{center}

\courssub{Emploi : obligation vs conseil}

\begin{propriete}[Nuances entre \all{müssen} et \all{sollen}]\ 
\begin{itemize}
\item \textbf{\all{müssen} + infinitif} : obligation forte, contrainte extérieure, nécessité objective. \all{Ich muss zum Arzt.} « Je dois aller chez le médecin. » (C'est nécessaire, j'ai mal).
\item \textbf{\all{sollen} + infinitif} : conseil, recommandation, attente morale ou médicale, obligation plus faible. \all{Du sollst mehr Wasser trinken.} « Tu devrais boire plus d'eau. » (C'est un bon conseil).
\end{itemize}
\end{propriete}

\begin{attention}[La négociation : \all{nicht müssen} $\neq$ \all{nicht sollen}]
C'est un piège majeur pour les francophones :
\begin{itemize}
\item \all{Du musst nicht kommen.} = « Tu n'as \emph{pas à} venir / Il n'est \emph{pas nécessaire} que tu viennes. » (Pas d'obligation).
\item \all{Du sollst nicht kommen.} = « Tu \emph{ne dois pas} venir / Il est \emph{déconseillé} que tu viennes. » (Interdiction / Déconseil).
\end{itemize}
En français, « ne pas devoir » est ambigu ; en allemand, le choix du modal lève l'ambiguïté.
\end{attention}

\courssub{Place du pronom réfléchi avec les modaux}

Avec un modal (\all{müssen}, \all{sollen}, \all{können}, \all{wollen}...), l'infinitif du verbe principal va en fin de phrase. Le pronom réfléchi se place \textbf{avant l'infinitif}, juste après le modal conjugué (ou après le sujet si le modal est en position 2).

\begin{exemplebox}[Modaux + Verbes pronominaux]
\begin{itemize}
\item \all{Er muss sich ausruhen.} « Il doit se reposer. » (Obligation)
\item \all{Du sollst dich mehr bewegen.} « Tu devrais bouger davantage. » (Conseil)
\item \all{Wir wollen uns fit halten.} « Nous voulons rester en forme. »
\item \all{Ich kann mich nicht konzentrieren.} « Je n'arrive pas à me concentrer. »
\end{itemize}
\end{exemplebox}

\begin{methode}[Construire une phrase avec modal + verbe pronominal]
\begin{enumerate}
\item Conjuguer le modal (\all{müssen}/\all{sollen}) à la 2ème place.
\item Placer le \textbf{sujet} (si ce n'est pas déjà fait).
\item Placer le \textbf{pronom réfléchi} (accordé au sujet) \textbf{immédiatement après le modal} (ou après le sujet).
\item Placer les compléments (temps, lieu, manière).
\item Mettre l'\textbf{infinitif du verbe principal} (avec sa particule séparable le cas échéant) \textbf{en toute fin de phrase}.
\end{enumerate}
Exemple : \all{Kevin / müssen / sich / nach der Schule / ausruhen} $\rightarrow$ \all{Kevin muss sich nach der Schule ausruhen.}
\end{methode}

\courssub{Texte d'application}

\begin{textebox}[Ein gesunder Tag in Yaoundé]
Mama Amina steht jeden Morgen um sechs Uhr auf. Sie bewegt sich viel: Sie geht zu Fuß zum Markt in Mokolo. „Ich fühle mich fit und gesund“, sagt sie und lacht. Ihr Sohn Kevin ist anders. Er bewegt sich wenig \textbf{und} spielt nur am Handy. Letzte Woche hat er sich erkältet und hatte Fieber. Der Arzt sagt: „Du musst dich mehr bewegen und du sollst dich nach der Schule ausruhen. Sport ist wichtig!“ Jetzt spielt Kevin jeden Nachmittag Fußball mit seinen Freunden. Er konzentriert sich in der Schule besser und erkältet sich nicht mehr so oft. Nach dem Training wäscht er sich die Hände und ruht sich eine Stunde aus. „Ich freue mich über den Sport“, sagt er. „Ich erhole mich schnell und fühle mich stark.“ Auch seine Schwester hält sich fit: Sie tanzt gern. Die ganze Familie interessiert sich jetzt für Sport und Gesundheit.
\end{textebox}

\textbf{Glossaire} : \all{aufstehen} (se lever, sep.) ; \all{der Markt, die Märkte} (le marché) ; \all{das Handy, die Handys} (le portable) ; \all{das Fieber} (la fièvre) ; \all{das Training} (l'entraînement) ; \all{stark} (fort) ; \all{die Gesundheit} (la santé) ; \all{sich erholen} (se rétablir) ; \all{sich fit halten} (se maintenir en forme) ; \all{sich interessieren für} (s'intéresser à).

\textbf{Questions de compréhension} :
\begin{enumerate}
\item \all{Wann steht Mama Amina auf ?}
\item \all{Warum geht Kevin jetzt zum Arzt nicht mehr so oft ?}
\item \all{Was macht Kevin nach dem Training ?}
\end{enumerate}

\begin{exoresolu}[Conjuguez et choisissez le bon modal]
Énoncé. 1. Conjuguez au présent : \all{sich erkälten} (ich) ; \all{sich ausruhen} (du) ; \all{sich fühlen} (wir) ; \all{sich konzentrieren} (ihr). \\
2. Traduisez : « Tu dois te reposer. » ; « Tu devrais bouger davantage. » ; « Je n'ai pas à me lever tôt. » (Pas d'obligation). \\
3. Choisissez la forme correcte : \all{Ich wasche (mich / mir) die Hände.} ; \all{Er (muss / soll) sich mehr bewegen.} (Conseil du médecin).
\tcblower
Solution détaillée. 
1. \all{ich erkälte mich} ; \all{du ruhst dich aus} (particule \all{aus} en finale) ; \all{wir fühlen uns} ; \all{ihr konzentriert euch}. 
2. \all{Du musst dich ausruhen.} (obligation forte) ; \all{Du sollst dich mehr bewegen.} (conseil) ; \all{Ich muss nicht früh aufstehen.} (\all{nicht müssen} = pas d'obligation). 
3. \all{Ich wasche mir die Hände} (COD \all{die Hände} $\rightarrow$ datif \all{mir}) ; \all{Er soll sich mehr bewegen.} (Conseil / Recommandation médicale).
\end{exoresolu}

\begin{exercice}[À vous]
1. Conjuguez au présent : \all{sich verletzen} (er), \all{sich fit halten} (sie, pluriel), \all{sich freuen} (ich). 
2. Traduisez en allemand : 
   a) « Nous bougeons beaucoup. » 
   b) « Elle doit se reposer. » (obligation) 
   c) « Je me lave les mains. » 
   d) « Tu ne dois pas fumer. » (Interdiction / Déconseil fort) 
   e) « Tu n'as pas à venir demain. » (Pas d'obligation).
3. Rédigez trois phrases sur votre journée avec \all{sich fühlen}, \all{sich bewegen} et \all{sich ausruhen}.
4. Complétez avec \all{müssen} ou \all{sollen} (conjugués) + infinitif : 
   \all{Du \_\_\_\_\_\_ mehr Wasser trinken.} (Conseil) \quad 
   \all{Ich \_\_\_\_\_\_ jetzt nach Hause gehen.} (Obligation, il est tard) \quad 
   \all{Wir \_\_\_\_\_\_ nicht so viel Süßes essen.} (Déconseil).
\end{exercice}

\begin{corrige}[Exercices : verbes pronominaux et modaux]
1. \all{er verletzt sich} ; \all{sie halten sich fit} ; \all{ich freue mich}. 
2. a) \all{Wir bewegen uns viel.} 
   b) \all{Sie muss sich ausruhen.} 
   c) \all{Ich wasche mir die Hände.} (Datif !) 
   d) \all{Du sollst nicht rauchen.} (\all{sollen} + \all{nicht} = interdiction/déconseil). 
   e) \all{Du musst nicht morgen kommen.} (\all{nicht müssen} = pas d'obligation). 
3. Exemple : \all{Ich fühle mich heute fit. Ich bewege mich viel in der Schule. Am Abend ruhe ich mich eine Stunde aus.} 
4. \all{Du sollst mehr Wasser trinken.} ; \all{Ich muss jetzt nach Hause gehen.} ; \all{Wir sollen nicht so viel Süßes essen.}
\end{corrige}

\begin{aretenir}[L'essentiel : Verbes pronominaux \& Conseils]
\begin{itemize}
\item \textbf{Verbe pronominal} = verbe + pronom réfléchi (\all{sich}) accordé au sujet. Place : après le verbe conjugué (phrase principale) ou après le sujet (subordonnée).
\item \textbf{Akkusativ} (\all{mich, dich, sich...}) : cas général. \textbf{Dativ} (\all{mir, dir, sich...}) : si un autre COD existe (souvent partie du corps : \all{Ich wasche mir die Hände}).
\item \textbf{Verbes à particule séparable} : \all{sich ausruhen} $\rightarrow$ \all{Ich ruhe mich aus}.
\item \textbf{Modaux} : \all{müssen} = obligation/nécessité ; \all{sollen} = conseil/recommandation.
\item \textbf{Négation} : \all{nicht müssen} = « ne pas avoir à » ; \all{nicht sollen} = « ne pas devoir / être déconseillé ».
\item \textbf{Ordre modal + pronominal} : Modal (conjugué) + Pr. réfléchi + Compléments + Infinitif (fin). Ex: \all{Du musst dich ausruhen.}
\item \textbf{Constructions à apprendre} : \all{sich erinnern an} (Akk), \all{sich interessieren für} (Akk), \all{sich freuen über} (Akk), \all{sich verletzen}, \all{sich erkälten}, \all{sich fit halten}.
\end{itemize}
\end{aretenir}

\coursec{Grammaire : le conseil avec « sollen » et « müssen »}

Après avoir étudié les verbes pronominaux qui permettent de décrire des actions que l'on fait sur soi-même (se laver, se reposer, se sentir), nous abordons maintenant un outil indispensable pour la communication quotidienne : \textbf{exprimer le conseil et l'obligation}. Dans notre texte support \emph{« Sport ist die beste Medizin »}, le médecin dit au patient : \all{Sie sollen mehr Sport treiben.} (Vous devriez faire plus de sport) et \all{Sie müssen sich mehr bewegen.} (Vous devez bouger davantage). Ces deux verbes modaux, \cle{sollen} et \cle{müssen}, structurent toute recommandation sanitaire ou règle de vie.

\courssub{Observation et distinction des nuances}

\begin{textebox}[Dialogue chez le médecin : conseils et obligations]
\all{Herr Doktor, ich fühle mich oft müde und habe Rückenschmerzen.}
« Docteur, je me sens souvent fatigué et j'ai mal au dos. »

\all{Sie sollen weniger sitzen und mehr Sport treiben.}
« Vous devriez moins rester assis et faire plus de sport. »

\all{Und Sie müssen jeden Tag mindestens 30 Minuten laufen.}
« Et vous devez courir au moins 30 minutes par jour. »

\all{Aber ich muss doch arbeiten! Ich habe keine Zeit.}
« Mais je dois travailler ! Je n'ai pas le temps. »

\all{Sie müssen nicht Marathon laufen. Spazieren reicht.}
« Vous n'êtes pas obligé de courir un marathon. Marcher suffit. »

\all{Ah, gut. Und ich soll nicht rauchen, oder?}
« Ah, bon. Et je ne devrais pas fumer, n'est-ce pas ? »

\all{Richtig. Sie dürfen nicht rauchen. Das ist schlecht für Ihre Gesundheit.}
« Exact. Vous ne devez pas fumer. C'est mauvais pour votre santé. »
\end{textebox}

\begin{definition}[Vocabulaire du dialogue]
\begin{itemize}
    \item \all{sich fühlen} (verbe pronominal) : se sentir
    \item \all{müde} : fatigué
    \item \all{die Rückenschmerzen} (pl.) : les maux de dos
    \item \all{sitzen} : rester assis
    \item \all{weniger} : moins
    \item \all{mindestens} : au moins
    \item \all{laufen} : courir / marcher
    \item \all{der Marathon, -s} : le marathon
    \item \all{spazieren} : se promener / marcher
    \item \all{reichen} : suffire
    \item \all{rauchen} : fumer
    \item \all{die Gesundheit} (sing.) : la santé
\end{itemize}
\end{definition}

\begin{propriete}[Sollen : le conseil et la recommandation]
Le verbe modal \cle{sollen} exprime un \textbf{conseil}, une \textbf{recommandation} ou une \textbf{attente venant d'un tiers} (médecin, parent, professeur, loi, opinion publique). Il se traduit par \emph{« devoir »} au sens de recommandation (\emph{je devrais, tu devrais}) ou par \emph{« il faut que »}.
\begin{itemize}
    \item \all{Der Arzt sagt: „Du sollst mehr Wasser trinken.“} = Le médecin dit : « Tu devrais boire plus d'eau. »
    \item \all{Man soll täglich Obst essen.} = Il faudrait manger des fruits chaque jour. (Recommandation générale, impersonnelle)
    \item \all{Was soll ich tun?} = Que dois-je faire ? / Que me conseillez-vous de faire ?
\end{itemize}
\end{propriete}

\begin{propriete}[Müssen : l'obligation et la nécessité forte]
Le verbe modal \cle{müssen} exprime une \textbf{obligation} (contrainte extérieure : loi, règle, ordre) ou une \textbf{nécessité} (vitale, logique). Il se traduit par \emph{« devoir »} (\emph{je dois, tu dois}) ou \emph{« être obligé de »}.
\begin{itemize}
    \item \all{Kranke müssen zum Arzt gehen.} = Les malades doivent aller chez le médecin.
    \item \all{Ich muss jetzt meine Medizin nehmen.} = Je dois prendre mon médicament maintenant.
    \item \all{Man muss gesund essen, um fit zu bleiben.} = Il faut manger sain pour rester en forme. (Nécessité logique)
\end{itemize}
\end{propriete}

\begin{center}
\begin{tabularx}{\linewidth}{|X|X|X|}
\hline
\rowcolor{popblueL} \textbf{Contexte} & \textbf{sollen (conseil)} & \textbf{müssen (obligation)} \\
\hline
Médecin au patient & \all{Sie sollen sich ausruhen.} (Reposez-vous.) & \all{Sie müssen das Antibiotikum nehmen.} (Prenez l'antibiotique.) \\
\hline
Règle générale / hygiène & \all{Man soll sich die Hände waschen.} (Il faut se laver les mains.) & \all{Man muss bei Rot anhalten.} (On doit s'arrêter au feu rouge.) \\
\hline
Objectif personnel & \all{Ich soll abnehmen.} (Je devrais maigrir — avis du médecin.) & \all{Ich muss abnehmen, sonst wird mein Herz krank.} (Je dois maigrir, sinon mon cœur tombera malade — nécessité vitale.) \\
\hline
\end{tabularx}
\end{center}

\courssub{Conjugaison au présent (Präsens)}

Les deux verbes sont des \textbf{verbes modaux} : la voyelle du radical change aux trois personnes du singulier, et il n'y a \textbf{pas de terminaison} à la 1\textsuperscript{re} et à la 3\textsuperscript{e} personnes du singulier.

\begin{center}
\begin{tabular}{|l|c|c|c|c|c|c|}
\hline
\rowcolor{popblueL} \textbf{Verbe} & \textbf{ich} & \textbf{du} & \textbf{er/sie/es} & \textbf{wir} & \textbf{ihr} & \textbf{sie/Sie} \\
\hline
\textbf{sollen} & soll & sollst & soll & sollen & sollt & sollen \\
\hline
\textbf{müssen} & muss & musst & muss & müssen & müsst & müssen \\
\hline
\end{tabular}
\end{center}

\begin{attention}[Pièges de conjugaison pour francophones]
\begin{itemize}
    \item \textbf{1\textsuperscript{re} personne du singulier :} pas de \emph{-e} ! \all{ich soll} (et non *ich solle), \all{ich muss} (et non *ich müsse).
    \item \textbf{3\textsuperscript{e} personne du singulier :} pas de \emph{-t} ! \all{er soll}, \all{er muss} (et non *er sollt, *er musst). C'est une particularité des modaux : \all{er muss}, mais \all{er lernt}.
    \item \textbf{Umlaut avec müssen :} le \emph{ü} n'apparaît qu'au pluriel (\all{wir müssen}, \all{ihr müsst}, \all{sie/Sie müssen}). Au singulier, on écrit \emph{u} : \all{ich muss}, \all{du musst}, \all{er muss}. \all{du musst} se prononce « mouste ».
    \item \textbf{ß ou ss ?} \all{müssen} s'écrit avec \emph{ü} et \emph{ss} (jamais \emph{ß} après un \emph{ü} bref).
\end{itemize}
\end{attention}

\courssub{Construction de la phrase : la place du verbe}

Règle d'or des verbes modaux : \textbf{le modal conjugué occupe la 2\textsuperscript{e} position (place du verbe conjugué) et l'infinitif du verbe principal va à la toute fin de la phrase.}

\begin{exemplebox}[Structure : modal + infinitif]
\all{Du \textbf{sollst} mehr Wasser \textbf{trinken}.} \\
(Sujet — modal en 2\textsuperscript{e} position — compléments — infinitif à la fin.)
\par\medskip
\all{Wir müssen jeden Tag eine halbe Stunde trainieren.} \\
« Nous devons nous entraîner une demi-heure chaque jour. »
\par\medskip
\all{Sie sollen sich mehr bewegen.} \\
« Vous devriez bouger davantage. » (Avec un verbe pronominal, le pronom réfléchi \all{sich} suit le modal et l'infinitif reste à la fin.)
\end{exemplebox}

\begin{propriete}[La négation : le piège « nicht müssen »]
La place de \cle{nicht} change radicalement le sens avec \textbf{müssen} :
\begin{itemize}
    \item \all{sollen + nicht} = \textbf{conseil de ne pas faire} (déconseil). \\
    \all{Du sollst nicht rauchen.} = Tu ne devrais pas fumer. (Conseil : évite de fumer.)
    \item \all{müssen + nicht} = \textbf{absence d'obligation} (liberté de ne pas faire). \textbf{Ce n'est pas une interdiction !} \\
    \all{Du musst nicht kommen.} = Tu n'es pas obligé de venir. (Tu peux venir si tu veux, mais ce n'est pas exigé.)
    \item Pour exprimer l'\textbf{interdiction} (il ne faut pas), on utilise \cle{nicht dürfen} : \\
    \all{Du darfst nicht rauchen.} = Il t'est interdit de fumer / Tu ne dois pas fumer.
\end{itemize}
\end{propriete}

\begin{center}
\begin{tabularx}{\linewidth}{|X|X|X|}
\hline
\rowcolor{popblueL} \textbf{Phrase allemande} & \textbf{Traduction française} & \textbf{Valeur} \\
\hline
\all{Du musst nicht arbeiten.} & Tu n'es pas obligé de travailler. & Absence d'obligation (liberté) \\
\hline
\all{Du darfst nicht arbeiten.} & Tu ne dois pas travailler. & Interdiction \\
\hline
\all{Du sollst nicht arbeiten.} & Tu ne devrais pas travailler. & Déconseil (ex. : avis du médecin) \\
\hline
\end{tabularx}
\end{center}

\courssub{Autres moyens d'exprimer le conseil}

Pour varier vos productions écrites (lettres de conseils, articles sur la santé), utilisez aussi ces structures :

\begin{propriete}[L'impératif (Imperativ) : conseil direct]
Utilisé pour des consignes claires, des modes d'emploi ou entre amis proches.
\begin{itemize}
    \item \all{Treib mehr Sport!} (Fais plus de sport ! — forme \emph{du}, sans \emph{-e} obligatoire)
    \item \all{Esst mehr Gemüse!} (Mangez plus de légumes ! — forme \emph{ihr})
    \item \all{Gehen Sie jeden Tag spazieren!} (Allez vous promener chaque jour ! — forme \emph{Sie}, polie)
\end{itemize}
\end{propriete}

\begin{propriete}[« Es ist wichtig, dass... » : conseil impersonnel]
Structure idéale pour les rédactions (articles, discours, lettres). Après \all{dass}, le verbe conjugué va à la fin de la subordonnée.
\begin{itemize}
    \item \all{Es ist wichtig, dass du regelmäßig trainierst.} = Il est important que tu t'entraînes régulièrement.
    \item \all{Es ist wichtig, dass man sich gesund ernährt.} = Il est important de se nourrir sainement.
\end{itemize}
\end{propriete}

\begin{propriete}[« Man sollte... » (Konjunktiv II de sollen) : conseil atténué]
\all{Man sollte mehr schlafen.} = On devrait dormir davantage. (Plus doux, plus poli que \all{Man soll...}.) Formes : \all{ich sollte, du solltest, er sollte, wir sollten, ihr solltet, sie sollten}.
\end{propriete}

\begin{center}
\begin{tabularx}{\linewidth}{|X|X|X|}
\hline
\rowcolor{popblueL} \textbf{Structure} & \textbf{Exemple} & \textbf{Registre / nuance} \\
\hline
\all{sollen + Inf.} & \all{Du sollst Sport treiben.} & Conseil standard (venant d'un tiers) \\
\hline
\all{müssen + Inf.} & \all{Du musst Sport treiben.} & Obligation forte / nécessité \\
\hline
Impératif & \all{Treib Sport!} & Ordre / conseil direct entre amis \\
\hline
\all{Es ist wichtig, dass...} & \all{Es ist wichtig, dass du Sport treibst.} & Argumentation écrite / formel \\
\hline
\all{Man sollte + Inf.} & \all{Man sollte Sport treiben.} & Conseil général atténué / politesse \\
\hline
\end{tabularx}
\end{center}

\courssub{Échelle des modalités : de la permission à l'obligation}

\begin{popfigure}
\begin{tikzpicture}[
    node distance=0.8cm and 1.5cm,
    box/.style={rectangle, rounded corners, draw=popdark, thick, fill=popcream, minimum height=1.2cm, minimum width=3.2cm, align=center, font=\small},
    arrow/.style={-Stealth, thick, popblue},
    labelnode/.style={font=\tiny, text=popmuted, align=center}
]
\node[box, fill=popgreenL, align=center] (duerfen){\textbf{dürfen}\\Permission};
\node[box, fill=popgoldL, right=of duerfen, align=center] (sollen){\textbf{sollen}\\Conseil};
\node[box, fill=poporangeL, right=of sollen, align=center] (muessen){\textbf{müssen}\\Obligation};

\draw[arrow] (duerfen.east) -- (sollen.west) node[midway, above, labelnode] {plus fort};
\draw[arrow] (sollen.east) -- (muessen.west) node[midway, above, labelnode] {plus fort};

\node[labelnode, below=0.3cm of duerfen, align=center]{\all{Du darfst bleiben.}\\(Tu peux rester.)};
\node[labelnode, below=0.3cm of sollen, align=center]{\all{Du sollst bleiben.}\\(Tu devrais rester.)};
\node[labelnode, below=0.3cm of muessen, align=center]{\all{Du musst bleiben.}\\(Tu dois rester.)};

\node[draw=poppink, thick, rounded corners, fill=poppinkL, below=1.6cm of sollen, align=center, font=\small, text width=8cm] (warn) {\textbf{Attention à la négation :}\\
\all{nicht müssen} = pas obligé (liberté)\\
\all{nicht dürfen} = interdit\\
\all{nicht sollen} = déconseillé};
\draw[poppink, dashed, thick] (muessen.south) |- (warn.east);
\draw[poppink, dashed, thick] (duerfen.south) |- (warn.west);
\end{tikzpicture}
\legende{Échelle de force modale : permission (dürfen) < conseil (sollen) < obligation (müssen). Attention aux négations !}
\end{popfigure}

\courssub{Exercice résolu : choisir le bon modal et la bonne forme}

\begin{exoresolu}[Conseils pour un ami malade]
\textbf{Énoncé :} Complétez les phrases avec \all{sollen} ou \all{müssen} (conjugués) + l'infinitif entre parenthèses. Justifiez le choix.
\begin{enumerate}
    \item Der Arzt sagt: Du \_\_\_\_\_\_\_\_ mehr Wasser \_\_\_\_\_\_\_\_ (trinken).
    \item In der Schule \_\_\_\_\_\_\_\_ die Schüler fleißig \_\_\_\_\_\_\_\_ (lernen).
    \item Wir \_\_\_\_\_\_\_\_ nicht so viel \_\_\_\_\_\_\_\_ (fernsehen), das ist ungesund.
    \item Ich \_\_\_\_\_\_\_\_ heute nicht \_\_\_\_\_\_\_\_ (arbeiten), ich habe frei.
    \item \_\_\_\_\_\_\_\_ ich zum Arzt \_\_\_\_\_\_\_\_ (gehen)? Mein Bein tut weh.
\end{enumerate}

\tcblower
\textbf{Solution détaillée :}
\begin{enumerate}
    \item \all{Du \textbf{sollst} mehr Wasser \textbf{trinken}.} \textbf{Justification :} le médecin donne un \textbf{conseil} (\all{sollen}). 2\textsuperscript{e} pers. sing. $\rightarrow$ \all{sollst}. Infinitif \all{trinken} à la fin.
    \item \all{In der Schule \textbf{müssen} die Schüler fleißig \textbf{lernen}.} \textbf{Justification :} règle scolaire = \textbf{obligation} (\all{müssen}). Sujet pluriel \all{die Schüler} $\rightarrow$ \all{müssen}. Attention : la phrase commence par un complément, donc le sujet vient après le verbe (inversion).
    \item \all{Wir \textbf{sollen} nicht so viel \textbf{fernsehen}.} \textbf{Justification :} \all{nicht sollen} = \textbf{déconseil} (c'est malsain). Sujet \all{wir} $\rightarrow$ \all{sollen}. \all{fernsehen} est un verbe à particule séparable, mais à l'infinitif avec un modal il reste en un seul mot à la fin.
    \item \all{Ich \textbf{muss} heute nicht \textbf{arbeiten}.} \textbf{Justification :} \textbf{piège !} \all{Ich muss nicht arbeiten} = \textbf{je ne suis pas obligé de travailler} (j'ai congé). Si c'était une interdiction, ce serait \all{ich darf nicht arbeiten}. Ici, absence d'obligation $\rightarrow$ \all{müssen} + \all{nicht}.
    \item \all{\textbf{Soll} ich zum Arzt \textbf{gehen}?} \textbf{Justification :} demande de \textbf{conseil} (\all{sollen}). 1\textsuperscript{re} pers. sing. $\rightarrow$ \all{soll}. Dans une question sans mot interrogatif, le modal est en 1\textsuperscript{re} position.
\end{enumerate}
\end{exoresolu}

\courssub{Exercices d'entraînement}

\begin{exercice}[Transformation : conseil et obligation]
Transformez les phrases en remplaçant le modal indiqué par l'autre modal (\all{sollen} $\leftrightarrow$ \all{müssen}) et adaptez la négation si nécessaire pour garder un sens logique. Traduisez.
\begin{enumerate}
    \item \all{Du sollst mehr Gemüse essen.} ($\rightarrow$ \all{müssen})
    \item \all{Er muss Tabletten nehmen.} ($\rightarrow$ \all{sollen})
    \item \all{Wir sollen nicht so spät schlafen gehen.} ($\rightarrow$ \all{müssen} + sens logique)
    \item \all{Ich muss nicht aufstehen.} (Gardez \all{müssen} et expliquez la différence de sens avec \all{ich darf nicht aufstehen}.)
\end{enumerate}
\end{exercice}

\begin{exercice}[Production écrite : message de conseils]
Votre ami camerounais \textbf{Paul} (17 ans, Yaoundé) vous écrit : il est stressé par les examens, dort mal, mange des fast-foods et ne fait plus de sport. Répondez-lui par un court message (5 à 6 phrases) en utilisant \textbf{au moins 3 structures différentes} de conseil (\all{sollen}, \all{müssen}, impératif, \all{Es ist wichtig, dass...}, \all{Man sollte...}). Utilisez le vocabulaire de la leçon (\all{der Stress}, \all{schlafen}, \all{sich ernähren}, \all{trainieren}, \all{sich bewegen}).
\end{exercice}

\begin{corrige}[Exercice 1 : Transformation]
\begin{enumerate}
    \item \all{Du musst mehr Gemüse essen.} = Tu dois manger plus de légumes. (Passage du conseil à l'obligation forte.)
    \item \all{Er soll Tabletten nehmen.} = Il devrait prendre des comprimés. (Passage de l'obligation au conseil.)
    \item \all{Wir müssen nicht so spät schlafen gehen.} = Nous ne sommes pas obligés d'aller nous coucher si tard. (Attention : \all{sollen nicht} = déconseil ; \all{müssen nicht} = absence d'obligation. Le sens change ! Pour garder le sens « il ne faut pas », il faudrait \all{Wir dürfen nicht so spät schlafen gehen.})
    \item \all{Ich muss nicht aufstehen.} = Je ne suis pas obligé de me lever (je peux rester au lit). \textbf{vs} \all{Ich darf nicht aufstehen.} = Il m'est interdit de me lever (interdiction médicale, par exemple).
\end{enumerate}
\end{corrige}

\begin{corrige}[Exercice 2 : Message de conseils (proposition)]
\all{Lieber Paul,}
\par\medskip
\all{Du \textbf{sollst} weniger Fastfood \textbf{essen} und dich gesünder \textbf{ernähren}.}
\all{Es \textbf{ist wichtig, dass} du jede Nacht acht Stunden \textbf{schläfst}.}
\all{Du \textbf{musst} dich unbedingt mehr \textbf{bewegen}, zum Beispiel joggen oder Fußball \textbf{spielen}.}
\all{\textbf{Treib} regelmäßig \textbf{Sport}, das hilft gegen den \textbf{Stress}!}
\all{Man \textbf{sollte} auch mal das Handy weglegen und sich \textbf{entspannen}.}
\par\medskip
\all{Viel Erfolg bei den Prüfungen! Dein Freund}
\par\medskip
« Cher Paul, tu devrais manger moins de fast-food et te nourrir plus sainement. Il est important que tu dormes huit heures chaque nuit. Tu dois absolument bouger davantage, par exemple faire du jogging ou jouer au football. Fais du sport régulièrement, cela aide contre le stress ! On devrait aussi poser le portable de temps en temps et se détendre. Bonne chance pour les examens ! Ton ami. »
\end{corrige}

\begin{savaistu}[Le conseil au Baccalauréat]
Dans les épreuves d'expression écrite, les sujets du type \emph{« Schreiben Sie einen Brief an einen kranken Freund »} (Écrivez une lettre à un ami malade) ou \emph{« Geben Sie Tipps für ein gesundes Leben »} (Donnez des conseils pour une vie saine) sont classiques. Les correcteurs attendent une \textbf{variété des structures de conseil} : ne vous contentez pas de \all{sollen} ! Montrez que vous maîtrisez l'impératif (\all{Geh spazieren!}), la structure impersonnelle (\all{Es ist wichtig, dass...}) et le Konjunktiv II (\all{Man sollte...}). Évitez aussi la traduction littérale du français « tu as à faire » : en allemand, on dit simplement \all{du musst} ou \all{du sollst}.
\end{savaistu}

\begin{aretenir}[Résumé : le conseil et l'obligation]
\begin{itemize}
    \item \cle{sollen} = conseil, recommandation (source : un tiers). $\rightarrow$ \emph{je devrais / il faudrait que.}
    \item \cle{müssen} = obligation, nécessité forte (source : règle, loi, nécessité). $\rightarrow$ \emph{je dois / il faut que.}
    \item \textbf{Conjugaison au présent :} \all{ich soll / ich muss}, \all{du sollst / du musst}, \all{er soll / er muss}, \all{wir/sie/Sie sollen / müssen}.
    \item \textbf{Syntaxe :} modal conjugué en 2\textsuperscript{e} position \dots\ infinitif à la fin de la phrase.
    \item \textbf{Négation :} \all{nicht sollen} = déconseil ; \all{nicht müssen} = pas obligé (liberté) ; \all{nicht dürfen} = interdit.
    \item \textbf{Alternatives :} impératif (\all{Treib Sport!}), \all{Es ist wichtig, dass...}, \all{Man sollte...}.
\end{itemize}
\end{aretenir}

\coursec{Méthode du commentaire et production guidée}

Après avoir étudié le vocabulaire de la santé et du sport, les verbes pronominaux et l'expression du conseil avec \all{sollen} et \all{müssen}, nous allons maintenant mettre ces acquis en pratique. Cette section vous apprend la méthode rigoureuse du commentaire de texte en allemand — une compétence clé pour le Baccalauréat — et vous guide vers une production écrite et orale autonome sur le thème « maladies et importance du sport ».

\courssub{La méthode du commentaire de texte en trois étapes}

\begin{methode}[Comment rédiger un commentaire de texte structuré en allemand]
Un bon commentaire commence toujours par identifier le thème et le type du texte, puis résume les idées sans recopier, et se termine par une opinion personnelle argumentée. Suivez ces trois étapes obligatoires :
\begin{enumerate}
    \item \textbf{Einleitung (Introduction) — 2 à 3 phrases} : Présentez le texte. Citez le \cle{titre}, le \cle{type de texte} (article, récit, dialogue, etc.), le \cle{thème} principal et, si possible, la \cle{source} ou l'auteur. Annoncez le plan de votre commentaire.
    \item \textbf{Hauptteil (Développement) — 2 paragraphes} :
    \begin{itemize}
        \item \textit{Inhaltliche Zusammenfassung (Résumé organisé) :} Relevez les idées principales dans l'ordre logique du texte (début, milieu, fin). Utilisez vos propres mots (pas de copier-coller). Liez les idées par des connecteurs logiques (\all{zuerst, dann, danach, schließlich, außerdem}).
        \item \textit{Analyse (Procédés et portée) :} Identifiez le \cle{ton} (humoristique, sérieux, alarmiste), le \cle{point de vue} (subjectif / objectif), les \cle{arguments} utilisés (exemples, chiffres, avis d'expert) et la \cle{structure} du texte (paragraphes, dialogue, énumération).
    \end{itemize}
    \item \textbf{Schluss (Conclusion) — 2 à 3 phrases} : Faites le \cle{bilan} : le texte est-il convaincant ? Donnez votre \cle{opinion personnelle} justifiée (\all{Meiner Meinung nach... / Ich finde, dass... / Meines Erachtens...}) en lien avec votre contexte (Cameroun, votre expérience).
\end{enumerate}
\end{methode}

\begin{propriete}[Connecteurs logiques indispensables pour le commentaire]
\begin{center}
\begin{tabularx}{\linewidth}{|X|X|X|}
\hline
\rowcolor{popblueL} \textbf{Français} & \textbf{Deutsch} & \textbf{Niveau / Remarque} \\
\hline
Tout d'abord / Premièrement & \all{zuerst / erstens} & A1 — Ordre chronologique \\
\hline
Ensuite / Puis & \all{dann / danach / anschließend} & A1 — Suite logique \\
\hline
De plus / En outre & \all{außerdem / zudem / ferner} & A2 — Ajout d'argument \\
\hline
Cependant / Mais & \all{jedoch / aber / allerdings} & A2 — Opposition / Concession \\
\hline
Par conséquent / C'est pourquoi & \all{deshalb / daher / folglich} & A2 — Conséquence \\
\hline
Finalement / En résumé & \all{schließlich / zusammenfassend} & A2 — Conclusion partielle \\
\hline
À mon avis / Selon moi & \all{meiner Meinung nach / meines Erachtens} & A2 — Opinion personnelle \\
\hline
\end{tabularx}
\end{center}
\end{propriete}

\begin{attention}[Pièges fréquents des francophones au commentaire]
\begin{itemize}
    \item \textbf{Recopier le texte :} Le résumé (\all{Zusammenfassung}) doit être en \emph{vos propres mots}. Citer un mot-clé est acceptable, mais pas une phrase entière.
    \item \textbf{Oublier l'introduction formelle :} Sans mention du titre, du type et du thème, le commentaire est « hors sujet » formel.
    \item \textbf{Mélanger résumé et analyse :} Séparez clairement les deux paragraphes du \all{Hauptteil}. Le premier = \emph{Was steht drin?} (contenu), le second = \emph{Wie ist es gemacht?} (forme/portée).
    \item \textbf{Opinion vague :} « C'est intéressant » (\all{Das ist interessant}) ne suffit pas. Justifiez : \all{Das ist wichtig, weil in Kamerun viele Jugendliche zu wenig Sport treiben.}
    \item \textbf{Allemand approximatif :} Utilisez les structures vues en classe : verbes pronominaux (\all{sich erholen, sich bewegen}), \all{sollen / müssen}, vocabulaire précis (\all{die Vorbeugung, die Lebenserwartung}).
\end{itemize}
\end{attention}

\courssub{Organigramme visuel de la structure du commentaire}

\begin{popfigure}
\begin{tikzpicture}[
    node distance=1.2cm and 2.5cm,
    every node/.style={font=\small, align=center},
    box/.style={draw=popblue, thick, rounded corners, fill=popblueL, minimum width=3.2cm, minimum height=1.1cm},
    arrow/.style={-Stealth, thick, popdark}
]
\node[box, align=center] (einl){Einleitung\\Thema, Textart, Quelle\\Gliederung ankündigen};
\node[box, below=of einl] (haupt) {Hauptteil};
\node[box, below left=0.8cm and 1.5cm of haupt, align=center] (inhalt){Inhalt\\Zusammenfassung\\Ideen in eigener Reihenfolge};
\node[box, below right=0.8cm and 1.5cm of haupt, align=center] (analyse){Analyse\\Ton, Argumente,\\Struktur, Sprachmittel};
\node[box, below=1.5cm of haupt, align=center] (schluss){Schluss\\Bilanz + Eigene Meinung\\begründet};

\draw[arrow] (einl) -- (haupt);
\draw[arrow] (haupt) -- (inhalt);
\draw[arrow] (haupt) -- (analyse);
\draw[arrow] (inhalt) -- (schluss);
\draw[arrow] (analyse) -- (schluss);
\end{tikzpicture}
\legende{Structure type d'un commentaire de texte (Einleitung → Hauptteil → Schluss)}
\end{popfigure}

\courssub{Commentaire modèle sur « Sport ist die beste Medizin »}

Le texte support de la leçon (section 1) était un article court de style journalistique. Voici un commentaire modèle rédigé au niveau attendu (A2/B1), d'environ 9 lignes, que vous pouvez analyser en classe.

\begin{textebox}[Commentaire modèle — Niveau A2/B1]
\textbf{Einleitung:} Der Text „Sport ist die beste Medizin“ ist ein kurzer Zeitungsartikel, der das Thema Gesundheit und Prävention behandelt. Er zeigt am Beispiel von Paul, wie wichtig regelmäßige Bewegung für die Genesung und das Wohlbefinden ist.

\textbf{Hauptteil — Inhalt:} Zuerst wird Paul vorgestellt: Er fühlt sich müde, hat Rückenschmerzen und oft Fieber. Der Arzt rät ihm, sich mehr zu bewegen und sich gesund zu ernähren. Paul folgt dem Rat: Er fängt an, joggen zu gehen und sich besser zu ernähren. Schließlich fühlt er sich viel besser und hat keine Rückenschmerzen mehr.

\textbf{Hauptteil — Analyse:} Der Text ist sachlich und appellativ geschrieben. Er nutzt ein konkretes Beispiel (Fallbeispiel), um das abstrakte Thema „Prävention“ verständlich zu machen. Die direkte Rede des Arztes („Sie sollen...“) wirkt überzeugend. Die Struktur ist chronologisch: Problem – Rat – Lösung.

\textbf{Schluss:} Meiner Meinung nach ist der Text sehr nützlich, auch für uns in Kamerun. Viele Jugendliche in Yaoundé oder Douala sitzen lange in der Schule oder vor Bildschirmen. Wir müssen uns mehr bewegen, um Krankheiten vorzubeugen. Sport ist wirklich die beste Medizin – und er kostet nichts.
\end{textebox}

\begin{exoresolu}[Analyse du commentaire modèle]
\textbf{Énoncé :} Identifiez dans le commentaire modèle les éléments suivants : 1) Les 3 composantes de l'introduction. 2) Les connecteurs logiques utilisés dans le résumé. 3) Les deux arguments de l'analyse. 4) La structure de la conclusion.

\tcblower
\textbf{Solution détaillée :}
\begin{enumerate}
    \item \textbf{Einleitung :} Titre (\all{„Sport ist die beste Medizin“}), Type (\all{Zeitungsartikel}), Thème (\all{Gesundheit und Prävention}), Annonce du plan (implicite : exemple de Paul $\rightarrow$ analyse).
    \item \textbf{Connecteurs (Zusammenfassung) :} \all{Zuerst} (Paul présenté), \all{Der Arzt rät...} (conseil), \all{Paul folgt dem Rat} (action), \all{Schließlich} (résultat final). Structure chronologique claire.
    \item \textbf{Analyse :} Argument 1 = \cle{Fallbeispiel} (exemple concret de Paul) pour rendre le thème accessible. Argument 2 = \cle{direkte Rede} (discours direct du médecin) + ton \cle{appellativ} (incitatif) pour convaincre le lecteur.
    \item \textbf{Schluss :} \cle{Bilan} (\all{sehr nützlich}) + \cle{Transfert au contexte camerounais} (\all{Yaoundé, Douala, Jugendliche, Bildschirme}) + \cle{Opinion justifiée} (\all{Sport kostet nichts, Vorbeugung}).
\end{enumerate}
\end{exoresolu}

\courssub{Production écrite guidée : « Ratschläge für einen kranken Freund »}

Vous allez maintenant rédiger un court message (SMS, e-mail ou mot) à un ami malade. L'objectif est de réinvestir \textbf{le vocabulaire thématique}, \textbf{les verbes pronominaux} et \textbf{les modaux de conseil (sollen / müssen)}.

\begin{methode}[Consignes de rédaction — 5 à 8 phrases]
\begin{enumerate}
    \item \textbf{Contexte :} Votre ami(e) \cle{Paul / Marie} est malade (grippe, fatigue, mal de dos, fièvre). Vous lui écrivez.
    \item \textbf{Contraintes obligatoires :}
    \begin{itemize}
        \item Utilisez \textbf{au moins 2 verbes pronominaux} différents (ex. \all{sich ausruhen, sich erholen, sich hinlegen, sich besser fühlen, sich bewegen}).
        \item Utilisez \textbf{au moins une fois \all{sollen} et une fois \all{müssen}} (conjugués à la 2e personne du singulier : \all{du sollst / du musst}).
        \item Employez \textbf{5 mots du vocabulaire de la leçon} (ex. \all{das Fieber, die Erkältung, der Arzt, das Medikament, die Gesundheit, trainieren, gesund}).
    \end{itemize}
    \item \textbf{Structure type :} Empathie $\rightarrow$ Conseil concret (sollen/müssen + infinitif) $\rightarrow$ Encouragement.
\end{enumerate}
\end{methode}

\begin{exemplebox}[Phrases modèles pour la production — À adapter et combiner]
\begin{itemize}
    \item \all{Liebe Marie, ich höre, dass du krank bist. Gute Besserung!} — « Chère Marie, j'apprends que tu es malade. Bon rétablissement ! »
    \item \all{Du sollst dich jetzt gut ausruhen und viel schlafen.} — « Tu devrais bien te reposer et beaucoup dormir. » (\cle{verbe pronominal \all{sich ausruhen}} + \cle{\all{sollen}})
    \item \all{Du musst viel Wasser oder Tee trinken, weil du Fieber hast.} — « Tu dois boire beaucoup d'eau ou du thé car tu as de la fièvre. » (\cle{\all{müssen}} + vocabulaire \all{Fieber})
    \item \all{Du sollst dich nicht zu viel bewegen, aber ein kleiner Spaziergang ist gut.} — « Tu ne devrais pas trop bouger, mais une petite promenade est bonne. » (\cle{verbe pronominal \all{sich bewegen}} + \cle{\all{sollen}} + négation)
    \item \all{Wenn es nicht besser wird, musst du zum Arzt gehen.} — « Si ça ne va pas mieux, tu dois aller chez le médecin. » (\cle{\all{müssen}} + \all{zum Arzt gehen})
    \item \all{Ich hoffe, du fühlst dich bald wieder gesund und fit.} — « J'espère que tu te sentiras bientôt en bonne santé et en forme. » (\cle{verbe pronominal \all{sich fühlen}} + adjectifs)
\end{itemize}
\end{exemplebox}

\begin{exoresolu}[Production écrite type — Corrigé commenté]
\textbf{Énoncé :} Rédigez le message pour Paul (grippe, fièvre, fatigue) en respectant les contraintes.

\tcblower
\textbf{Proposition de rédaction (7 phrases) :}
\all{Lieber Paul, ich habe gehört, dass du eine starke Erkältung und Fieber hast. Du sollst dich jetzt sofort ins Bett legen und dich gut ausruhen. Du musst viel Tee trinken und die Medizin vom Arzt nehmen. Du sollst dich nicht stressen und nicht arbeiten. Wenn das Fieber nicht sinkt, musst du wieder zum Arzt gehen. Ich wünsche dir gute Besserung und hoffe, du fühlst dich bald wieder gesund. Dein Freund, Jean.}

\textbf{Vérification des contraintes :}
\begin{itemize}
    \item Verbes pronominaux (3) : \all{sich hinlegen, sich ausruhen, sich stressen, sich fühlen} (\textbf{validé}).
    \item \all{sollen} (2x) : \all{Du sollst dich... hinlegen}, \all{Du sollst dich nicht stressen} (\textbf{validé}).
    \item \all{müssen} (2x) : \all{Du musst viel Tee trinken}, \all{musst du wieder zum Arzt gehen} (\textbf{validé}).
    \item Vocabulaire thématique (6) : \all{Erkältung, Fieber, Bett, Medizin, Arzt, gesund, Gute Besserung} (\textbf{validé}).
    \item Cohérence : ton amical (\all{du}), progression logique (symptômes $\rightarrow$ repos $\rightarrow$ boisson/médocs $\rightarrow$ alerte $\rightarrow$ vœux).
\end{itemize}
\end{exoresolu}

\begin{attention}[Erreurs à éviter dans votre production]
\begin{itemize}
    \item \textbf{Oubli du pronom réfléchi :} \all{*Du sollst ausruhen} $\rightarrow$ \textbf{Faux}. Correct : \all{Du sollst \textbf{dich} ausruhen}.
    \item \textbf{Mauvaise place du verbe à l'infinitif :} \all{*Du sollst dich ausruhen gut} $\rightarrow$ \textbf{Faux}. L'infinitif (\all{ausruhen}) va à la \textbf{fin de la phrase} (ou de la proposition).
    \item \textbf{Confusion \all{sollen} / \all{müssen} :} \all{sollen} = conseil / recommandation (tu \emph{devrais}), \all{müssen} = obligation / nécessité forte (tu \emph{dois}). Ne les échangez pas au hasard.
    \item \textbf{Ordre des mots en subordinate (si clause \all{wenn}) :} \all{Wenn das Fieber nicht sinkt, musst du...} (verbe en 2e position dans la principale). Pas : \all{*...du musst}.
    \item \textbf{Majuscules aux noms :} \all{das Fieber, der Arzt, die Medizin, die Gesundheit} — \textbf{toujours majuscule}.
\end{itemize}
\end{attention}

\courssub{Production orale guidée : Dialogue médecin / patient (Rollenspiel)}

L'expression orale au Bac évalue votre capacité à interagir dans une situation authentique. Voici un cadre structuré pour un jeu de rôle en binôme.

\begin{experience}[Jeu de rôle : « Beim Hausarzt » — Durée : 3–4 minutes par binôme]
\textbf{Rôles :}
\begin{itemize}
    \item \textbf{Partner A (Der Patient / Die Patientin) :} Vous avez des symptômes depuis 2–3 jours. Décrivez-les précisément (vocabulaire : \all{Fieber, Husten, Halsschmerzen, Kopfschmerzen, Gliederschmerzen, Übelkeit, Schwindel, Ausschlag}). Dites ce que vous avez déjà fait (médicaments, repos). Posez une question sur la durée de la maladie.
    \item \textbf{Partner B (Der Arzt / Die Ärztin) :} Vous utilisez le \textbf{vouvoiement (\all{Sie})} ! Accueillez le patient, posez des questions ciblées (\all{Seit wann? Wo genau tut es weh? Haben Sie Fieber gemessen?}). Donnez \textbf{3 conseils clairs} avec \all{Sie sollen...} et \all{Sie müssen...}. Prescrivez un arrêt de travail si nécessaire (\all{Krankschreibung}). Terminez par \all{Gute Besserung!}
\end{itemize}

\textbf{Préparation (5 min) :} Chaque élève note 5–6 mots-clés (pas de phrases toutes faites) sur une fiche.
\textbf{Jeu (3–4 min) :} Jouez sans lire. L'enseignant circule et note : prononciation, correction grammaticale (surtout \all{Sie}-formes, verbes pronominaux \all{sich fühlen, sich hinlegen}), richesse lexicale, interaction (réagir à l'autre).
\textbf{Débriefing (2 min) :} Échange de rôles. Correction collective de 2–3 erreurs fréquentes entendues.
\end{experience}

\begin{propriete}[Le vouvoiement médical obligatoire — Formes \all{Sie} au présent]
Dans un contexte médical formel, le médecin s'adresse au patient par \all{Sie} (3e personne du pluriel comme forme de politesse). Le patient répond aussi par \all{Sie}.
\begin{center}
\begin{tabular}{|l|l|l|}
\hline
\rowcolor{popblueL} \textbf{Verbe (Infinitif)} & \textbf{Forme \all{Sie} (Médecin $\rightarrow$ Patient)} & \textbf{Traduction} \\
\hline
\all{fühlen} (se sentir) & \all{Wie fühlen Sie sich?} & Comment vous sentez-vous ? \\
\hline
\all{haben} (avoir) & \all{Haben Sie Fieber?} & Avez-vous de la fièvre ? \\
\hline
\all{nehmen} (prendre) & \all{Nehmen Sie dieses Medikament.} & Prenez ce médicament. \\
\hline
\all{sollen} (devoir/conseil) & \all{Sie sollen sich ausruhen.} & Vous devriez vous reposer. \\
\hline
\all{müssen} (devoir/obligation) & \all{Sie müssen zum Spezialisten gehen.} & Vous devez aller chez le spécialiste. \\
\hline
\all{bleiben} (rester) & \all{Sie sollen zu Hause bleiben.} & Vous devriez rester à la maison. \\
\hline
\all{kommen} (venir/revenir) & \all{Kommen Sie in drei Tagen wieder.} & Revenez dans trois jours. \\
\hline
\end{tabular}
\end{center}
\end{propriete}

\begin{textebox}[Modèle de dialogue guidé (squelette à compléter oralement)]
\textbf{Arzt:} Guten Tag, Frau / Herr \dots Was fehlt Ihnen?
\textbf{Patient:} Guten Tag, Herr Doktor. Ich fühle mich seit \dots Tagen nicht gut. Ich habe \dots (symptômes) und \dots
\textbf{Arzt:} Haben Sie Fieber gemessen?
\textbf{Patient:} Ja, \dots Grad. / Nein, aber ich friere / schwitze.
\textbf{Arzt:} Ich untersuche Sie kurz. \dots (Diagnose simulée). Sie haben eine \dots
\textbf{Arzt:} \textbf{Sie sollen} sich \dots (verbe pronominal) und \dots \textbf{Sie müssen} \dots (action forte). Nehmen Sie \dots (médicament).
\textbf{Patient:} Wie lange dauert die Krankheit?
\textbf{Arzt:} Normalerweise \dots Tage. \textbf{Kommen Sie} wieder, wenn es nicht besser wird. Hier ist die Krankschreibung für \dots Tage.
\textbf{Patient:} Vielen Dank, Herr Doktor. Auf Wiederhören.
\textbf{Arzt:} Gute Besserung! Auf Wiedersehen.
\end{textebox}

\courssub{Grille d'auto-évaluation et critères du Baccalauréat}

Avant de rendre votre production écrite ou de passer l'oral, utilisez cette grille pour vous auto-évaluer. Elle reprend les critères officiels d'évaluation du MINESEC pour l'allemand LV2 en Première.

\begin{aretenir}[Grille d'auto-évaluation — Production écrite/orale « Maladies et Sport »]
\begin{center}
\begin{tabularx}{\linewidth}{|X|c|c|c|}
\hline
\rowcolor{popblueL} \textbf{Critère (2 pts chacun = 8 pts total)} & \textbf{Non acquis (0)} & \textbf{En cours (1)} & \textbf{Acquis (2)} \\
\hline
\textbf{1. Vocabulaire thématique} (maladies, symptômes, sport, corps, conseils) & $\le$ 3 mots utilisés, répétitions, anglicismes & 4–6 mots corrects, variés, orthographe approximative & $\ge$ 7 mots précis, articles/pluriels corrects, champ lexical riche \\
\hline
\textbf{2. Verbes pronominaux} (forme, place, sens) & Absents ou tous faux (pronom manquant, place) & 1–2 verbes utilisés, 1 erreur de forme/place & $\ge$ 2 verbes différents, formes parfaites (pronom + infinitif final) \\
\hline
\textbf{3. \all{sollen} / \all{müssen}} (conjugaison, emploi, nuance) & Absents ou confondus, mauvaise conjugaison & Présents tous les deux, 1 erreur conjugaison/emploi & Tous deux corrects, conjugués (\all{du/ Sie}), nuance conseil/obligation respectée \\
\hline
\textbf{4. Cohérence \& Communication} (structure, connecteurs, ton, interaction) & Phrases disjointes, pas de lien, ton inapproprié & Structure visible, quelques connecteurs, ton globalement ok & Structure claire (intro/dev/concl ou dialogue logique), connecteurs variés, ton naturel (Sie/du), communication réussie \\
\hline
\end{tabularx}
\end{center}
\textbf{Score $\ge$ 6/8 :} Objectif atteint. \textbf{Score 4–5 :} Réviser les points faibles. \textbf{Score $\le$ 3 :} Revoir la leçon + refaire l'exercice.
\end{aretenir}

\begin{propriete}[Critères officiels du Bac — Production écrite (Extrait programme MINESEC)]
Au Baccalauréat, la production écrite est évaluée sur la pertinence du contenu, l'exactitude grammaticale et la richesse du vocabulaire : réutiliser les structures de la leçon est la meilleure stratégie.
\begin{itemize}
    \item \textbf{Pertinence (4 pts) :} Réponse à la consigne, respect du format (lettre, mail, dialogue, article), idées développées.
    \item \textbf{Correction grammaticale (6 pts) :} Conjugaison (présent, parfait, modaux), cas (accusatif/datif après verbes/prépositions), ordre des mots (verbe en 2e position, infinitif final), orthographe (majuscules noms, Umlaute, ß).
    \item \textbf{Richesse lexicale (4 pts) :} Variété du vocabulaire (pas de répétition de \all{gut/schlecht/krank}), utilisation du vocabulaire spécifique du thème, expressions idiomatiques simples (\all{Gute Besserung, sich erholen, auf sich aufpassen}).
    \item \textbf{Cohérence / Cohésion (4 pts) :} Connecteurs logiques, paragraphation, progression thématique, registre de langue adapté (formel/informel).
\end{itemize}
\textbf{Total : 18 pts} (souvent converti sur 20).
\end{propriete}

\courssub{Exercices de consolidation finale}

\begin{exercice}[Entraînement commentaire — Texte court]
Lisez le texte suivant et rédigez un \textbf{mini-commentaire} (introduction 2 phrases, résumé 3 phrases, conclusion 2 phrases) en respectant la méthode. Utilisez au moins 3 connecteurs du tableau de la \textbf{Propriété 2}.

\begin{textebox}[Texte d'entraînement : „Bewegung im Alltag“]
In Deutschland bewegen sich immer weniger Menschen genug. Laut einer Studie der WHO treiben 44\% der Frauen und 40\% der Männer zu wenig Sport. Das führt zu Herz-Kreislauf-Erkrankungen, Rückenschmerzen und Übergewicht. Experten raten: Man soll sich mindestens 150 Minuten pro Woche moderat bewegen. Das kann Radfahren, schnelles Gehen oder Schwimmen sein. Auch kleine Änderungen helfen: Treppe statt Aufzug, Fahrrad statt Auto für kurze Wege. Wer sich regelmäßig bewegt, lebt länger und gesünder. Fängt man heute an, profitiert man morgen.
\end{textebox}

\textbf{Glossaire :} \all{die Studie} (l'étude), \all{die WHO} (OMS), \all{die Herz-Kreislauf-Erkrankung} (maladie cardiovasculaire), \all{das Übergewicht} (surpoids), \all{mäßig} (modéré), \all{die Treppe} (l'escalier), \all{der Aufzug} (l'ascenseur), \all{profitieren} (en bénéficier).
\end{exercice}

\begin{exercice}[Thème : Français $\rightarrow$ Allemand — Structures de la leçon]
Traduisez en allemand en utilisant \textbf{obligatoirement} la structure indiquée entre parenthèses.
\begin{enumerate}
    \item Tu devrais te reposer maintenant. (\all{du sollen + sich ausruhen})
    \item Il doit prendre le médicament trois fois par jour. (\all{er müssen + das Medikament nehmen})
    \item Nous devrions faire plus de sport. (\all{wir sollen + Sport treiben})
    \item Elle ne doit pas aller à l'école aujourd'hui. (\all{sie nicht müssen + zur Schule gehen})
    \item Je me sens beaucoup mieux aujourd'hui. (\all{ich mich fühlen + viel besser})
    \item Le médecin m'a dit : « Vous devez boire beaucoup. » (\all{Sie müssen + viel trinken} — style direct)
    \item Mes amis conseillent : « Tu ne devrais pas fumer. » (\all{du nicht sollen + rauchen})
    \item Allongez-vous (formel) et reposez-vous ! (\all{Sie sich hinlegen + sich ausruhen} — impératif formel = infinitif + \all{Sie})
\end{enumerate}
\end{exercice}

\begin{corrige}[Exercice 2 — Thème]
\begin{enumerate}
    \item \all{Du sollst dich jetzt ausruhen.}
    \item \all{Er muss das Medikament dreimal am Tag nehmen.} (ordre : temps - manière - lieu - infinitif final)
    \item \all{Wir sollen mehr Sport treiben.}
    \item \all{Sie muss heute nicht zur Schule gehen.} (\all{nicht} ne se place pas avant \all{müssen} pour « pas obligation » ; \all{zur} = \all{zu der})
    \item \all{Ich fühle mich heute viel besser.}
    \item \all{Der Arzt sagte: „Sie müssen viel trinken.“}
    \item \all{Meine Freunde sagten: „Du sollst nicht rauchen.“}
    \item \all{Legen Sie sich hin und ruhen Sie sich aus!} (Impératif formel : infinitif + \all{Sie} ; verbes pronominaux : pronom \all{sich} après le verbe).
\end{enumerate}
\end{corrige}

\begin{savaistu}[Le sport au Cameroun et en Allemagne — Contexte culturel]
Au Cameroun, le football est roi (\all{der Königssport}), mais la \cle{pratique sportive quotidienne} (\all{regelmäßige Bewegung}) reste faible chez les jeunes scolarisés (manque d'infrastructures, emploi du temps chargé). En Allemagne, le \cle{sport de masse} (\all{Breitensport}) est institutionnalisé : \all{Turnvereine} (clubs de gymnastique) dans chaque village, \all{Schulsport} obligatoire 2–3h/semaine, \all{Betriebssport} (sport en entreprise). L'assurance maladie (\all{Krankenkasse}) rembourse même des cours de prévention (\all{Präventionskurse : Yoga, Rückenschule, Nordic Walking}). Le slogan \all{„Sport ist die beste Medizin“} n'est pas une métaphore : c'est une politique de santé publique. Au Cameroun, des initiatives comme la \cle{Journée nationale du sport} ou les \cle{Jeux de la Francophonie} (Yaoundé 2027) visent à changer les mentalités. \all{Bewegung ist Leben!} — Le mouvement, c'est la vie !
\end{savaistu}

\newpage

\coursec{Activité d'intégration}

\courssub{Objectifs de l'activité}
Cette activité d'intégration vise à mobiliser l'ensemble des savoirs et savoir-faire travaillés lors de la leçon \emph{« Les maladies et l'importance de la pratique du sport »} dans une tâche finale complexe, authentique et motivante. Elle s'inscrit dans l'approche actionnelle préconisée par le programme MINESEC : l'élève devient un acteur social qui utilise la langue pour \textbf{agir} dans une situation de communication réelle.

\begin{aretenir}[Compétences ciblées]
\begin{itemize}
    \item \textbf{Compétence langagière :} Réinvestir le lexique de la santé (\cle{die Krankheit, die Gesundheit, das Fieber, der Schmerz, die Erkältung, die Verletzung}), du sport (\cle{der Sport, das Training, die Bewegung, fit, trainieren}) et du corps (\cle{der Körper, der Kopf, der Bauch, der Rücken, das Bein, der Arm}) avec articles et pluriels corrects.
    \item \textbf{Compétence grammaticale :} Maîtriser la conjugaison des \textbf{verbes pronominaux} au présent de l'indicatif (\all{sich bewegen, sich ausruhen, sich ernähren, sich waschen, sich verletzen, sich fühlen}) et l'ordre des mots (place du pronom réfléchi). Utiliser correctement les verbes modaux \textbf{\all{sollen}} (conseil, recommandation) et \textbf{\all{müssen}} (obligation, nécessité forte) pour formuler des conseils de santé.
    \item \textbf{Compétence discursive :} Produire un texte court, cohérent et structuré (affiche : titre, slogans, conseils structurés) et le présenter à l'oral de manière fluide (2 minutes), en soignant la prononciation (voyelles longues et courtes, \all{ß}, \all{ü}, \all{ö}) et l'intonation.
    \item \textbf{Compétence stratégique :} Travailler en équipe (négociation du contenu, répartition des rôles : rédacteur, illustrateur, présentateur, secrétaire), gérer le temps, s'auto-évaluer et évaluer les pairs (vote).
    \item \textbf{Compétence culturelle :} S'approprier le concept germanophone de \cle{Prävention} (prévention) et de \cle{Gesundheitstag} (journée de la santé), très ancré dans le système scolaire allemand, autrichien et suisse, et le transposer au contexte camerounais (lycée de Yaoundé ou de Douala, cantine, bendskins, chaleur, paludisme, nutrition).
\end{itemize}
\end{aretenir}

\courssub{Situation}
\emph{Contexte :} Le \textbf{Lycée Bilingue de Nkolbisson (Yaoundé)} organise son premier \all{Gesundheitstag} (Journée de la Santé) en partenariat avec le \all{Goethe-Institut Kamerun}. L'administration a lancé un concours inter-classes de Première : chaque groupe d'élèves doit concevoir une \textbf{affiche promotionnelle bilingue (allemand/français)} pour sensibiliser les camarades aux bonnes habitudes d'hygiène de vie, de nutrition et d'activité physique.

\begin{textebox}[Document déclencheur : Appel à projets du Proviseur (extrait)]
\textbf{Avis de concours : „Gesundheitstag 2025 – Dein Körper, deine Verantwortung!“}

Liebe Schülerinnen und Schüler der Klasse 1A,

der Schulvorstand und der Deutschclub laden euch herzlich ein, unsere erste „Gesundheitstag“-Kampagne zu gestalten. Thema: „Gesund leben in Yaoundé – Sport, Ernährung und Hygiene“.

\textbf{Eure Mission:}
Erstellt ein Plakat (DIN A2) auf Deutsch mit französischer Übersetzung, das eure Mitschüler motiviert, gesünder zu leben. Das Plakat muss enthalten:
\begin{enumerate}
    \item Einen \textbf{auffälligen Titel} (Hauptslogan).
    \item \textbf{Fünf konkrete Ratschläge} (Tipps) mit \all{sollen} oder \all{müssen}.
    \item \textbf{Drei verschiedene reflexive Verben} (z. B. \all{sich bewegen}, \all{sich ausruhen}, \all{sich gesund ernähren}).
    \item Eine \textbf{Illustration} (Zeichnung, Collage, Schema) passend zum Thema.
    \item Die \textbf{Namen der Gruppenmitglieder} und die Klasse.
\end{enumerate}

\textbf{Präsentation:} Jede Gruppe stellt ihr Plakat in zwei Minuten vor (auf Deutsch). Die beste Kampagne (Inhalt, Sprache, Kreativität, Präsentation) gewinnt einen Preis (Bücher, Sportzubehör, Besuch im Goethe-Institut).

\textbf{Abgabetermin:} in zwei Wochen. Viel Erfolg!

\emph{Der Schulleiter / Le Proviseur}
\end{textebox}

\begin{savaistu}[Contexte culturel : le « Gesundheitstag » en pays germanophones]
En Allemagne, en Autriche et en Suisse, le \cle{Gesundheitstag} (masculin : \all{der Gesundheitstag}) est une institution dans les écoles (\cle{die Schule, -n}), les entreprises (\cle{der Betrieb, -e}) et les communes (\cle{die Gemeinde, -n}). Il vise la \cle{Prävention} : éviter l'apparition des maladies (\cle{verhindern}) par l'éducation à la santé (\cle{die Gesundheitserziehung}). Thèmes classiques : \cle{der Bewegungsmangel} (manque d'exercice), \cle{die Ernährung} (nutrition), \cle{der Stress} (le stress). Au Cameroun, les défis sont proches (sédentarité, alimentation trop grasse ou trop sucrée, paludisme, fièvre typhoïde), mais le vocabulaire allemand permet d'en parler avec une structure logique. Proverbe allemand à retenir : \all{Vorbeugen ist besser als heilen.} « Mieux vaut prévenir que guérir. »
\end{savaistu}

\courssub{Consignes}
\emph{Travail en groupes de 4 élèves. Durée totale estimée : 2 h (préparation 1 h 15 + présentations 45 min).}

\begin{enumerate}
    \item \textbf{Analyse du document déclencheur (10 min) :} Lisez l'avis du Proviseur. Identifiez les \textbf{contraintes obligatoires} (nombre de conseils, verbes pronominaux, titre, illustration, durée orale). Soulignez le vocabulaire inconnu et demandez l'explication au professeur ou utilisez le dictionnaire.
    \item \textbf{Remue-méninges lexical et grammatical (15 min) :} Sur une feuille de brouillon, chaque groupe dresse trois listes :
    \begin{itemize}
        \item \textit{Liste A :} 10 mots ou expressions du thème « Santé / Maladie / Sport » avec article et pluriel (ex. \all{der Rat, die Ratschläge} ; \all{die Krankheit, die Krankheiten} ; \all{das Training, die Trainings}).
        \item \textit{Liste B :} 5 verbes pronominaux utiles pour la santé, conjugués à la 2\up{e} personne du pluriel : \all{ihr bewegt euch, ihr ruht euch aus, ihr ernährt euch gesund, ihr wascht euch, ihr fühlt euch fit}.
        \item \textit{Liste C :} 5 structures de conseil alternant \all{sollen} (recommandation) et \all{müssen} (obligation vitale ou hygiène).
    \end{itemize}
    \item \textbf{Rédaction de l'affiche – Étape 1 : titre et structure (15 min) :} Choisissez un \textbf{titre accrocheur} (impératif, question, jeu de mots). Exemples : \all{„Beweg dich, Kamerun!“} ; \all{„Dein Körper dankt dir!“} ; \all{„Fit für die Zukunft – Gesundheitstag 2025“}. Structurez l'affiche en 5 zones (un conseil par zone).
    \item \textbf{Rédaction de l'affiche – Étape 2 : les 5 conseils (25 min) :} Rédigez les 5 phrases-conseils en allemand.
    \begin{itemize}
        \item \textbf{Contrainte 1 :} au moins 2 conseils avec \all{ihr sollt} (ex. \all{Ihr sollt jeden Tag 30 Minuten Sport treiben.}).
        \item \textbf{Contrainte 2 :} au moins 2 conseils avec \all{ihr müsst} (ex. \all{Ihr müsst vor dem Essen die Hände waschen.} / \all{Ihr müsst genug Wasser trinken.}).
        \item \textbf{Contrainte 3 :} intégrez \textbf{exactement 3 verbes pronominaux différents} dans l'affiche (ex. \all{ihr bewegt euch}, \all{ihr ruht euch aus}, \all{ihr ernährt euch gesund}). Cochez-les sur votre brouillon.
        \item \textbf{Contrainte 4 :} ajoutez la traduction française \emph{sous} chaque phrase allemande (petits caractères).
    \end{itemize}
    \item \textbf{Finalisation visuelle (15 min) :} Réalisez l'illustration (dessin, schéma « Assiette équilibrée », silhouette qui court, pictogrammes de lavage des mains). Soignez la lisibilité (couleurs, écriture, espacement). Collez les textes.
    \item \textbf{Préparation de l'oral (15 min) :} Répartissez les rôles : \textbf{Présentateur 1} (titre + contexte), \textbf{Présentateur 2} (conseils 1 à 3 + verbes pronominaux), \textbf{Présentateur 3} (conseils 4 et 5 + \all{müssen}/\all{sollen}), \textbf{Présentateur 4} (conclusion + appel à l'action). Répétez à voix haute : fluidité, regard vers le public, pas de lecture mot à mot. Chronométrez-vous (2 min maximum).
    \item \textbf{Présentations et vote (45 min) :} Chaque groupe présente. La classe note sur une grille d'évaluation (contenu /5, grammaire et lexique /5, oral /5, créativité /5, respect des consignes /5). Le professeur note aussi. Dépouillement et remise des prix.
\end{enumerate}

\courssub{Ressources à mobiliser}
Voici un mémo synthétique à garder sous les yeux pendant l'activité. Il ne remplace pas le cours complet, mais rappelle les points de vigilance pour la production.

\begin{propriete}[Rappel : les verbes pronominaux (reflexive Verben) au présent]
\textbf{Structure :} Sujet + verbe conjugué (2\up{e} position) + \textbf{pronom réfléchi à l'accusatif} (\all{mich, dich, sich, uns, euch, sich}) + compléments. Le pronom se place \textbf{immédiatement après le verbe conjugué} (ou après le sujet en cas d'inversion).
\begin{center}
\begin{tabular}{|l|l|l|l|l|l|}
\hline
\rowcolor{popblueL} \textbf{Infinitif} & \textbf{ich} & \textbf{du} & \textbf{er/sie/es} & \textbf{wir} & \textbf{ihr / sie} \\ \hline
\all{sich bewegen} & bewege mich & bewegst dich & bewegt sich & bewegen uns & bewegt euch / bewegen sich \\ \hline
\all{sich ausruhen} & ruhe mich aus & ruhst dich aus & ruht sich aus & ruhen uns aus & ruht euch aus / ruhen sich aus \\ \hline
\all{sich ernähren} & ernähre mich & ernährst dich & ernährt sich & ernähren uns & ernährt euch / ernähren sich \\ \hline
\all{sich waschen} & wasche mich & wäschst dich & wäscht sich & waschen uns & wascht euch / waschen sich \\ \hline
\all{sich fühlen} & fühle mich & fühlst dich & fühlt sich & fühlen uns & fühlt euch / fühlen sich \\ \hline
\all{sich verletzen} & verletze mich & verletzt dich & verletzt sich & verletzen uns & verletzt euch / verletzen sich \\ \hline
\end{tabular}
\end{center}
\textbf{Sens et exemples en contexte santé :}
\begin{itemize}
    \item \all{sich bewegen} : bouger, faire de l'exercice. \all{Ihr bewegt euch zu wenig!} « Vous ne bougez pas assez ! »
    \item \all{sich ausruhen} : se reposer (verbe à particule séparable). \all{Nach dem Sport ruht ihr euch aus.} « Après le sport, vous vous reposez. »
    \item \all{sich ernähren} : se nourrir. \all{Ihr ernährt euch gesund.} « Vous vous nourrissez sainement. »
    \item \all{sich waschen} : se laver. \all{Ihr wascht euch die Hände.} « Vous vous lavez les mains. »
    \item \all{sich fühlen} : se sentir. \all{Wie fühlt ihr euch heute? – Wir fühlen uns fit.} « Comment vous sentez-vous aujourd'hui ? – Nous nous sentons en forme. »
    \item \all{sich verletzen} : se blesser. \all{Beim Fußball verletzt man sich leicht.} « Au football, on se blesse facilement. »
\end{itemize}
\end{propriete}

\begin{attention}[Pièges fréquents pour les francophones]
\begin{itemize}
    \item \textbf{Oubli du pronom :} *\all{Ich bewege} est incorrect ; il faut \all{Ich bewege \textbf{mich}}. En français « je bouge » est intransitif ; en allemand, \all{sich bewegen} est obligatoirement pronominal.
    \item \textbf{Place du pronom :} Sujet + verbe + \textbf{pronom} + compléments. Ex. \all{Ihr bewegt euch heute viel.} (et non *\all{Ihr euch bewegt...}, structure possible seulement dans une subordonnée : \all{..., weil ihr euch viel bewegt.}).
    \item \textbf{Particule séparable :} \all{sich ausruhen} $\rightarrow$ \all{ich ruhe mich \textbf{aus}}. La particule \all{aus} va à la \textbf{fin de la phrase principale}.
    \item \textbf{Accusatif ou datif :} \all{sich waschen} (pronom à l'accusatif : \all{mich, dich, sich...}) mais \all{sich die Hände waschen} (pronom au \textbf{datif} : \all{mir, dir, sich...} + accusatif \all{die Hände}) : \all{Ich wasche \textbf{mir} die Hände.} Pour l'affiche, utilisez la forme simple \all{sich waschen} ou, si la structure est maîtrisée, \all{sich die Hände waschen}.
    \item \textbf{Conjugaison de \all{waschen} :} verbe fort à vocalisme changeant : \all{du \textbf{wäschst}, er \textbf{wäscht}} (et non *\all{du waschst}).
\end{itemize}
\end{attention}

\begin{propriete}[Rappel : le conseil – \all{sollen} ou \all{müssen}]
\begin{center}
\begin{tabularx}{\linewidth}{|X|X|X|}
\hline
\rowcolor{poporangeL} \textbf{Verbe modal} & \textbf{Nuance / degré} & \textbf{Exemple pour l'affiche} \\ \hline
\all{sollen} (\all{ihr sollt}) & \textbf{Recommandation, avis médical} : « il est conseillé de... ». Marque une norme sociale ou experte. Moins contraignant que \all{müssen}. & \all{Ihr sollt täglich Obst und Gemüse essen.} « Vous devriez manger des fruits et des légumes chaque jour. » \\ \hline
\all{müssen} (\all{ihr müsst}) & \textbf{Obligation forte, nécessité vitale, règle d'hygiène stricte} : « il faut / vous devez ». & \all{Ihr müsst bei Fieber zum Arzt gehen.} « Vous devez aller chez le médecin en cas de fièvre. » \all{Ihr müsst sauberes Wasser trinken.} « Vous devez boire de l'eau propre. » \\ \hline
\end{tabularx}
\end{center}
\textbf{Ordre des mots avec un modal :} Sujet + modal conjugué (2\up{e} position) + compléments + \textbf{infinitif en fin de phrase}. Ex. \all{Ihr sollt jeden Tag Sport treiben.}
\end{propriete}

\begin{methode}[Choisir entre \all{sollen} et \all{müssen} : arbre de décision]
\begin{enumerate}
    \item Est-ce une \textbf{règle d'hygiène vitale} (lavage des mains, eau potable, fièvre $\rightarrow$ médecin) ou une \textbf{règle scolaire} ? $\rightarrow$ \textbf{\all{müssen}}.
    \item Est-ce une \textbf{recommandation pour le bien-être à long terme} (sport, sommeil, légumes, pause écran) ? $\rightarrow$ \textbf{\all{sollen}}.
    \item \textbf{Astuce :} alternez pour varier le ton de l'affiche. Trop de \all{müssen} donne un ton autoritaire ; trop de \all{sollen} rend tout optionnel. Visez l'équilibre (2 \all{müssen} et 3 \all{sollen}, ou l'inverse).
\end{enumerate}
\end{methode}

\begin{definition}[Lexique utile pour l'affiche (mots-outils)]
\begin{itemize}
    \item \all{der Gesundheitstag, -e} (journée de la santé) ; \all{das Plakat, -e} (affiche) ; \all{der Tipp, -s} / \all{der Ratschlag, die Ratschläge} (conseil) ; \all{die Regel, -n} (règle).
    \item \all{gesund / ungesund} (sain / malsain) ; \all{fit / unfit} (en forme / pas en forme) ; \all{aktiv / inaktiv} (actif / inactif) ; \all{ausgewogen} (équilibré, pour l'alimentation).
    \item \all{der Bewegungsmangel} (manque d'exercice, sédentarité) ; \all{die Zivilisationskrankheit, -en} (maladie de civilisation : diabète, obésité, hypertension) ; \all{die Prävention} (prévention).
    \item \all{täglich} (quotidiennement) ; \all{regelmäßig} (régulièrement) ; \all{genug / ausreichend} (assez / suffisamment) ; \all{mindestens} (au moins) ; \all{vermeiden} (éviter) ; \all{verhindern} (prévenir, empêcher).
    \item \all{das Trinkwasser} (eau potable) ; \all{sich die Hände waschen} (se laver les mains) ; \all{der Schlaf} (sommeil, sans pluriel) ; \all{der Bildschirm, -e} (écran).
\end{itemize}
\end{definition}

\courssub{Proposition de production et corrigé commenté}
\emph{Voici une production modèle (niveau A2 solide) respectant toutes les contraintes. Elle sert de référence pour l'enseignant et de support de correction collective après les présentations.}

\begin{exoresolu}[Production modèle : affiche « Fit fürs Leben – Gesundheitstag am Lycée Nkolbisson »]
\textbf{Énoncé :} Concevoir l'affiche complète (texte allemand + traduction française + commentaire linguistique) pour le groupe « Die Gesundheitsbotschafter » (Les Ambassadeurs de la Santé).

\tcblower

\textbf{1. TITRE PRINCIPAL (grand format) :}
\all{„Fit fürs Leben – Dein Körper, deine Wahl!“}
« En forme pour la vie – Ton corps, ton choix ! »
\emph{Commentaire :} \all{fürs} = contraction de \all{für das}. \all{Dein Körper, deine Wahl} : parallélisme avec possessifs (\all{dein/deine}), impact fort.

\textbf{2. CONSEIL 1 (hygiène vitale – \cle{müssen}) :}
\all{„Ihr müsst euch vor jedem Essen gründlich die Hände waschen.“}
« Vous devez vous laver soigneusement les mains avant chaque repas. »
\emph{Analyse :} \all{müssen} (obligation d'hygiène). \all{vor jedem Essen} (préposition \all{vor} + datif : \all{jedem Essen}). Structure pronominale au datif : \all{euch} (datif) + \all{die Hände} (accusatif). \all{gründlich} (adverbe de manière). Infinitif \all{waschen} en fin de phrase.

\textbf{3. CONSEIL 2 (nutrition – \cle{sollen} + verbe pronominal n°1) :}
\all{„Ihr sollt euch täglich ausgewogen ernähren: viel Gemüse, wenig Fett.“}
« Vous devriez vous nourrir chaque jour de façon équilibrée : beaucoup de légumes, peu de gras. »
\emph{Analyse :} \all{sollen} (recommandation). \textbf{Verbe pronominal n°1 :} \all{sich ernähren} $\rightarrow$ \all{ihr sollt \textbf{euch} ... ernähren} : le pronom \all{euch} suit directement le verbe conjugué. \all{täglich} (adverbe de temps), \all{ausgewogen} (adverbe de manière).

\textbf{4. CONSEIL 3 (activité physique – \cle{sollen} + verbe pronominal n°2) :}
\all{„Ihr sollt euch jeden Tag mindestens 30 Minuten bewegen – kein Bendskin für kurze Strecken!“}
« Vous devriez bouger au moins 30 minutes par jour – pas de bendskin pour les courtes distances ! »
\emph{Analyse :} ancrage camerounais (\all{das Bendskin} : moto-taxi). \textbf{Verbe pronominal n°2 :} \all{sich bewegen} $\rightarrow$ \all{ihr sollt \textbf{euch} ... bewegen}. \all{mindestens 30 Minuten} (complément de mesure). \all{kein} (négation devant un nom).

\textbf{5. CONSEIL 4 (sommeil et récupération – \cle{müssen} + verbe pronominal n°3) :}
\all{„Ihr müsst nachts acht Stunden schlafen und euch gut ausruhen.“}
« Vous devez dormir huit heures par nuit et bien vous reposer. »
\emph{Analyse :} \all{müssen} (besoin physiologique vital). \all{nachts} (adverbe de temps). \textbf{Verbe pronominal n°3 :} \all{sich ausruhen} (séparable) $\rightarrow$ \all{und \textbf{euch} gut \textbf{ausruhen}} : à l'infinitif avec modal, la particule reste soudée (\all{ausruhen}) en fin de phrase ; le pronom \all{euch} précède les compléments.

\textbf{6. CONSEIL 5 (hydratation et écrans – \cle{sollen}) :}
\all{„Ihr sollt zwei Liter Wasser trinken und das Handy vor dem Schlafen weglegen.“}
« Vous devriez boire deux litres d'eau et poser le portable avant de dormir. »
\emph{Analyse :} \all{sollen} (conseil de prévention). \all{zwei Liter Wasser} (accusatif). \all{vor dem Schlafen} (préposition \all{vor} + datif du substantif verbal \all{das Schlafen}). \all{weglegen} (verbe à particule séparable, à l'infinitif en fin de phrase).

\textbf{7. BAS DE L'AFFICHE (signatures) :}
\all{Gesundheitstag 2025 | Lycée Bilingue Nkolbisson | Goethe-Institut Kamerun}
\all{Gruppe: Die Gesundheitsbotschafter (Amadou, Fatou, Pierre, Grace) | Klasse: 1A}
\end{exoresolu}

\begin{aretenir}[Commentaire global des choix linguistiques (pour la correction collective)]
\begin{itemize}
    \item \textbf{Respect des contraintes :} 5 conseils (2 \all{müssen} : hygiène et sommeil ; 3 \all{sollen} : nutrition, sport, écrans). 3 verbes pronominaux distincts (\all{sich ernähren, sich bewegen, sich ausruhen}) correctement employés à la 2\up{e} personne du pluriel avec le pronom \all{euch} bien placé.
    \item \textbf{Ordre des mots :} respect strict du schéma : Sujet (\all{Ihr}) – modal conjugué (\all{müsst}/\all{sollt}) – pronom réfléchi (\all{euch}) – compléments (temps, manière, lieu) – infinitif en fin de phrase. Ex. \all{Ihr (S) müsst (V) euch (pronom) gut (manière) ausruhen (infinitif).}
    \item \textbf{Orthographe et majuscules :} tous les noms (\all{Tag, Essen, Hände, Wasser, Handy, Schlaf}) prennent la majuscule. Le \all{ß} s'écrit après voyelle longue ou diphtongue : \all{müssen, gründlich, außerdem}. Attention aux trémas : \all{gesund, fühlen, Gemüse}.
    \item \textbf{Registre :} \all{ihr} (tutoiement pluriel) est approprié pour s'adresser aux camarades. Ton direct et engageant, non moralisateur grâce à l'alternance \all{müssen}/\all{sollen}.
    \item \textbf{Ancrage local :} \all{Bendskin}, \all{Lycée Nkolbisson}, \all{Goethe-Institut Kamerun} ancrent la tâche dans la réalité vécue des élèves.
\end{itemize}
\end{aretenir}

\courssub{Bilan de l'activité}
Cette activité du \cle{Gesundheitstag} constitue le \textbf{point d'ancrage final} de la séquence. Elle permet de vérifier l'atteinte des objectifs terminaux du chapitre 3 du programme de Première A.

\begin{experience}[Grille d'auto-évaluation élève (à distribuer après le vote)]
\emph{Entourez la mention correspondante : \textbf{Acquis} / \textbf{En cours} / \textbf{À revoir}.}
\begin{center}
\begin{tabularx}{\linewidth}{|X|l|}
\hline
\rowcolor{popgreenL} \textbf{Critère de compétence} & \textbf{Mon niveau} \\ \hline
\textbf{Lexique :} j'utilise au moins 10 mots du thème (maladie, sport, corps) avec article et pluriel corrects. & Acquis / En cours / À revoir \\ \hline
\textbf{Grammaire (verbes pronominaux) :} je conjugue 3 verbes pronominaux au présent (\all{ihr}) sans erreur de place du pronom (\all{euch}). & Acquis / En cours / À revoir \\ \hline
\textbf{Grammaire (conseil) :} je distingue \all{sollen} (conseil) et \all{müssen} (obligation) et je les conjugue correctement (\all{ihr sollt / ihr müsst}). & Acquis / En cours / À revoir \\ \hline
\textbf{Écrit (affiche) :} mon texte est lisible, structuré (titre, 5 conseils, traduction), et l'orthographe allemande est respectée (majuscules, ß, trémas). & Acquis / En cours / À revoir \\ \hline
\textbf{Oral (présentation) :} je parle 30 secondes sans lire, avec une prononciation claire, le regard vers le public. & Acquis / En cours / À revoir \\ \hline
\textbf{Travail d'équipe :} j'ai participé activement (idées, dessin, rédaction, oral), respecté les autres et tenu les délais. & Acquis / En cours / À revoir \\ \hline
\end{tabularx}
\end{center}
\end{experience}

\begin{methode}[Prolongements et différenciation pédagogique]
\begin{itemize}
    \item \textbf{Approfondissement (niveau A2+) :} transformer l'affiche en \textbf{spot audio ou vidéo de 60 secondes} (enregistrement sur téléphone). Ajouter des connecteurs logiques (\all{zuerst, außerdem, schließlich} : « d'abord, de plus, enfin »).
    \item \textbf{Soutien (niveau A1/A1+) :} fournir une \textbf{affiche à trous} (texte à compléter) avec une banque de mots (verbes pronominaux à l'infinitif, \all{sollen}/\all{müssen} à conjuguer). Réduire à 3 conseils au lieu de 5. Présentation en binôme.
    \item \textbf{Interdisciplinarité (SVT / EPS) :} collaborer avec le professeur de SVT (maladies non transmissibles, nutrition) et d'EPS (endurance, étirements). L'affiche devient support d'exposé en SVT (en français) avec glossaire allemand.
    \item \textbf{Numérique :} créer l'affiche sur un outil de conception en ligne (intégration d'un QR code vers une vidéo d'exercices). Exposition virtuelle sur l'environnement numérique de travail de l'établissement.
\end{itemize}
\end{methode}

\begin{savaistu}[Le saviez-vous ? Santé scolaire : Cameroun et Allemagne]
Au Cameroun, la médecine scolaire existe mais est souvent sous-dotée (une infirmerie, parfois un seul infirmier pour un effectif très important). En Allemagne, le service de santé scolaire (\cle{der Schularzt, die Schulärzte} : médecin scolaire) est bien organisé : bilans de santé, vaccinations, dépistages (vue, audition) et séances d'éducation à la santé. Le \cle{Gesundheitstag} est souvent le point d'orgue annuel de ce service. Au Cameroun, les clubs « Santé » ou les « Jeunes Sauveteurs » (Croix-Rouge) jouent ce rôle de prévention par les pairs, très efficace. Votre affiche s'inscrit dans cette dynamique : « par les jeunes, pour les jeunes ».
\end{savaistu}

\emph{Fin de l'activité d'intégration. Le professeur collecte les grilles d'auto-évaluation, note les productions (écrit + oral) selon la grille officielle de l'évaluation des compétences et remplit la fiche de suivi de progression du chapitre 3.}

\newpage

\coursec{Méthodes et exercices résolus}

\coursec{Leçon ALA12 : Les maladies et l'importance de la pratique du sport}

\courssub{Texte support : « Sport ist die beste Medizin »}

\begin{textebox}[Texte de compréhension écrit/oral — Niveau A1+/A2 — Contexte Cameroun / Allemagne]
Viele Menschen in Kamerun und in Deutschland sind oft krank. Sie bewegen sich nicht genug und essen ungesund. Der Arzt sagt: „Sport ist die beste Medizin.“ Wer regelmäßig trainiert, stärkt sein Herz und fühlt sich besser. Man muss nicht Profi sein: dreißig Minuten Sport pro Tag reichen. Auch Fußballspielen mit Freunden oder ein Spaziergang sind gut für die Gesundheit.
\end{textebox}

\courssub{Savoir-faire 1 : Comprendre un texte sur la santé et le sport}

\begin{methode}[Lire et comprendre un texte de santé en allemand]
\begin{enumerate}
\item \textbf{Première lecture globale.} Lis le titre et le texte une première fois sans dictionnaire. Repère le thème (ici : \all{die Krankheit}, \all{der Sport}, \all{die Gesundheit}) et demande-toi : qui parle ? de quoi ? pourquoi ?
\item \textbf{Repérage des mots-clés.} Souligne les mots du champ lexical de la santé (\all{der Arzt, das Fieber, krank, gesund, das Herz}) et du sport (\all{trainieren, der Sport, die Bewegung, Fußballspielen, der Spaziergang}). Beaucoup se devinent grâce au français : \all{die Medizin} = la médecine, \all{der Patient} = le patient, \all{regelmäßig} = régulièrement.
\item \textbf{Lecture détaillée.} Pour chaque question, localise la phrase du texte qui contient la réponse. Recopie la structure de la phrase et remplace uniquement les éléments demandés.
\item \textbf{Rédaction de la réponse.} Réponds toujours par une phrase complète, en reprenant les mots de la question : \all{Warum ist Sport wichtig?} $\rightarrow$ \all{Sport ist wichtig, weil...} Attention : après \all{weil}, le verbe conjugué va à la fin !
\item \textbf{Vérification.} Relis : verbe en 2\textsuperscript{e} position dans la phrase principale, majuscule à tous les noms, accord du verbe avec le sujet.
\end{enumerate}
\end{methode}

\begin{exoresolu}[Compréhension de texte : « Sport ist die beste Medizin »]
\textbf{Énoncé.} Texte : \all{Viele Menschen in Kamerun und in Deutschland sind oft krank. Sie bewegen sich nicht genug und essen ungesund. Der Arzt sagt: „Sport ist die beste Medizin." Wer regelmäßig trainiert, stärkt sein Herz und fühlt sich besser. Man muss nicht Profi sein: dreißig Minuten Sport pro Tag reichen. Auch Fußballspielen mit Freunden oder ein Spaziergang sind gut für die Gesundheit.}\\
Réponds en allemand par des phrases complètes :
\begin{enumerate}
\item \all{Warum sind viele Menschen oft krank?}
\item \all{Was sagt der Arzt?}
\item \all{Wie viel Sport pro Tag reicht?}
\item \all{Was stärkt das Herz?}
\end{enumerate}
\tcblower
\textbf{Raisonnement.} Pour chaque question, on repère la phrase correspondante du texte, puis on rédige une phrase complète en respectant la place du verbe.
\begin{enumerate}
\item La cause est donnée par \all{Sie bewegen sich nicht genug und essen ungesund}. On répond avec \all{weil} (verbe à la fin ; \all{sich bewegen} est pronominal, le pronom réfléchi reste devant le verbe à la fin).\\
\textbf{Réponse :} \all{Viele Menschen sind oft krank, weil sie sich nicht genug bewegen und ungesund essen.} « Beaucoup de gens sont souvent malades parce qu'ils ne bougent pas assez et mangent malsain. »
\item Citation directe du texte transformée en discours indirect avec \all{dass}.\\
\textbf{Réponse :} \all{Der Arzt sagt, dass Sport die beste Medizin ist.} « Le médecin dit que le sport est la meilleure médecine. » (Après \all{dass}, verbe \all{ist} à la fin.)
\item Le texte dit : \all{dreißig Minuten Sport pro Tag reichen}. Sujet pluriel \rightarrow verbe au pluriel.\\
\textbf{Réponse :} \all{Dreißig Minuten Sport pro Tag reichen.} « Trente minutes de sport par jour suffisent. »
\item Le texte dit : \all{Wer regelmäßig trainiert, stärkt sein Herz}.\\
\textbf{Réponse :} \all{Regelmäßiges Training stärkt das Herz.} ou \all{Wer regelmäßig trainiert, stärkt sein Herz.} « Un entraînement régulier renforce le cœur. »
\end{enumerate}
\textbf{Conclusion.} Chaque réponse reprend les mots de la question, contient un verbe conjugué correctement placé et des noms avec majuscule.
\end{exoresolu}

\begin{exoresolu}[Résumé du texte (Résumé / Zusammenfassung)]
\textbf{Énoncé.} Résume le texte en 3 phrases en allemand en utilisant tes propres mots (vocabulaire du texte).
\tcblower
\textbf{Modèle de résumé :}\\
\all{Viele Menschen in Kamerun und Deutschland sind krank, weil sie sich zu wenig bewegen. Der Arzt empfiehlt Sport als beste Medizin. Schon dreißig Minuten Bewegung am Tag, wie Fußball oder Spazierengehen, helfen der Gesundheit.}\\
\textbf{Traduction :} « Beaucoup de gens au Cameroun et en Allemagne sont malades parce qu'ils bougent trop peu. Le médecin recommande le sport comme meilleure médecine. Déjà trente minutes de mouvement par jour, comme le football ou la marche, aident la santé. »
\end{exoresolu}

\courssub{Savoir-faire 2 : Employer le vocabulaire de la santé et du sport}

\begin{propriete}[Vocabulaire thématique : Noms, articles, pluriels et expressions clés]
\begin{center}
\begin{tabularx}{\linewidth}{|X|X|X|X|}
\hline \rowcolor{popblueL} \textbf{Deutsch (Nom + Article)} & \textbf{Pluriel} & \textbf{Français} & \textbf{Exemple / Expression} \\ \hline
\all{die Krankheit} & \all{die Krankheiten} & la maladie & \all{Er hat eine schwere Krankheit.} \\ \hline
\all{der Arzt} & \all{die Ärzte} & le médecin & \all{Ich gehe zum Arzt.} \\ \hline
\all{das Fieber} & \all{—} & la fièvre & \all{Er hat hohes Fieber.} \\ \hline
\all{der Sport} & \all{—} & le sport & \all{Sport treiben.} \\ \hline
\all{die Gesundheit} & \all{—} & la santé & \all{für die Gesundheit.} \\ \hline
\all{die Bewegung} & \all{die Bewegungen} & le mouvement / l'exercice & \all{Sich genug bewegen.} \\ \hline
\all{das Training} & \all{die Trainings} & l'entraînement & \all{Regelmäßig trainieren.} \\ \hline
\all{der Patient} & \all{die Patienten} & le patient & \all{Der Patient wartet.} \\ \hline
\all{das Medikament} & \all{die Medikamente} & le médicament & \all{Das Medikament nehmen.} \\ \hline
\all{die Medizin} & \all{die Medizinen} & la médecine / le médicament & \all{Die Medizin hilft.} \\ \hline
\end{tabularx}
\end{center}
\textbf{Verbes clés :} \all{krank sein} (être malade), \all{gesund sein/bleiben} (être/rester en bonne santé), \all{trainieren} (s'entraîner), \all{Sport treiben} (faire du sport), \all{zum Arzt gehen} (aller chez le médecin), \all{sich erkälten} (s'enrhumer), \all{sich ausruhen} (se reposer), \all{sich fühlen} (se sentir), \all{stärken} (renforcer).
\end{propriete}

\begin{methode}[Utiliser le lexique thématique avec article et pluriel]
\begin{enumerate}
\item \textbf{Apprends chaque nom avec son article et son pluriel} : \all{die Krankheit, -en}, \all{der Arzt, die Ärzte} (Umlaut !), \all{das Fieber} (pas de pluriel usuel), \all{der Sport}, \all{die Gesundheit}.
\item \textbf{Retiens les verbes associés et leur régime} : \all{Fieber haben} (sans article), \all{Sport treiben}, \all{zum Arzt gehen} (\all{zu} + datif \rightarrow \all{zum}).
\item \textbf{Exploite les familles de mots} : \all{krank} (adj.) $\rightarrow$ \all{die Krankheit} (nom), \all{der Kranke} (nominalisé, le malade) ; \all{gesund} $\rightarrow$ \all{die Gesundheit}, \all{gesund leben/essen}.
\item \textbf{Choisis le bon cas} : après \all{für} (préposition accusative) : \all{für die Gesundheit} ; après \all{mit} (préposition datif) : \all{mit Freunden} ; après \all{zu} (préposition datif) : \all{zum Arzt}.
\item \textbf{Réutilise le lexique dans tes propres phrases} (contexte camerounais) : \all{In Yaoundé spielen viele Jugendliche Fußball. Das ist gut für die Gesundheit.}
\end{enumerate}
\end{methode}

\begin{exoresolu}[Thème : phrases à traduire avec le vocabulaire de la santé]
\textbf{Énoncé.} Traduis en allemand :
\begin{enumerate}
\item Mon frère est malade, il a de la fièvre.
\item Le sport est bon pour la santé.
\item Nous allons chez le médecin à Yaoundé.
\item Le patient prend le médicament.
\item Elle veut rester en bonne santé.
\end{enumerate}
\tcblower
\textbf{Raisonnement et solutions.}
\begin{enumerate}
\item « Mon frère » = \all{mein Bruder} (sujet, nominatif). « Être malade » = \all{krank sein}. « Avoir de la fièvre » = \all{Fieber haben} (accusatif sans article).\\
\all{Mein Bruder ist krank, er hat Fieber.}
\item « Être bon pour » = \all{gut sein für} + accusatif ; \all{die Gesundheit} féminin, article \all{die} inchangé à l'accusatif.\\
\all{Sport ist gut für die Gesundheit.}
\item « Aller chez le médecin » = \all{zum Arzt gehen} (\all{zum} = \all{zu dem}, datif masculin). Verbe \all{gehen} en 2\textsuperscript{e} position. Lieu \all{in Yaoundé} avant le but.\\
\all{Wir gehen in Yaoundé zum Arzt.}
\item « Le patient » = \all{der Patient} (sujet) $\rightarrow$ \all{der Patient} (masc.) mais ici objet \rightarrow \all{den Patient} ? Non, \all{nehmen} prend l'accusatif : \all{den Patienten} (n-déclinaison !).\\
\all{Der Patient nimmt das Medikament.} (Sujet \all{Der Patient}, objet \all{das Medikament}).
\item « Rester en bonne santé » = \all{gesund bleiben}.\\
\all{Sie will gesund bleiben.}
\end{enumerate}
\textbf{Conclusion.} Pièges : \cle{datif après zu (zum Arzt)}, \cle{accusatif après für}, \cle{majuscule aux noms}, \cle{n-déclinaison (der Patient, den Patienten)}, \cle{Fieber haben sans article}.
\end{exoresolu}

\courssub{Savoir-faire 3 : Conjuguer les verbes pronominaux au présent}

\begin{propriete}[Conjugaison des verbes pronominaux au présent — Types courants]
\begin{center}
\begin{tabular}{|l|l|l|l|}
\hline \rowcolor{popblueL} \textbf{Personne} & \textbf{\all{sich bewegen} (bouger)} & \textbf{\all{sich fühlen} (se sentir)} & \textbf{\all{sich erkälten} (s'enrhumer)} \\ \hline
ich & \all{ich bewege mich} & \all{ich fühle mich} & \all{ich erkälte mich} \\ \hline
du & \all{du bewegst dich} & \all{du fühlst dich} & \all{du erkältest dich} \\ \hline
er/sie/es & \all{er bewegt sich} & \all{er fühlt sich} & \all{er erkältet sich} \\ \hline
wir & \all{wir bewegen uns} & \all{wir fühlen uns} & \all{wir erkälten uns} \\ \hline
ihr & \all{ihr bewegt euch} & \all{ihr fühlt euch} & \all{ihr erkältet euch} \\ \hline
sie/Sie & \all{sie bewegen sich} & \all{sie fühlen sich} & \all{sie erkälten sich} \\ \hline
\end{tabular}
\end{center}
\textbf{Verbes à particule séparable :} \all{sich ausruhen} (se reposer), \all{sich anziehen} (s'habiller), \all{sich beeilen} (se dépêcher).\\
\all{ich ruhe mich aus, du ruhst dich aus, er ruht sich aus, wir ruhen uns aus, ihr ruht euch aus, sie ruhen sich aus.} (Particule \all{aus} en fin de phrase).
\end{propriete}

\begin{methode}[Maîtriser les verbes pronominaux]
\begin{enumerate}
\item \textbf{Identifie le verbe pronominal} : il s'emploie avec un pronom réfléchi (\all{sich}) à l'infinitif : \all{sich bewegen}, \all{sich fühlen}, \all{sich erkälten}, \all{sich ausruhen}, \all{sich waschen}.
\item \textbf{Accorde le pronom réfléchi avec le sujet} : \all{ich mich, du dich, er/sie/es sich, wir uns, ihr euch, sie/Sie sich}.
\item \textbf{Place le pronom juste après le verbe conjugué} (position 3 en phrase déclarative) : \all{Ich fühle mich heute besser.} Jamais avant le verbe.
\item \textbf{Accusatif ou Datif ?} \textbf{Règle :} Si le verbe a déjà un complément d'objet direct (accusatif), le pronom réfléchi passe au \textbf{datif} : \all{Ich wasche mir die Hände.} (Je me lave les mains \rightarrow \all{mir} datif, \all{die Hände} accusatif). Sinon, pronom à l'\textbf{accusatif} : \all{Ich bewege mich.}
\item \textbf{Verbe à particule séparable} : la particule va à la fin de la phrase : \all{Wir ruhen uns eine Stunde aus.}
\item \textbf{Dans la subordonnée (avec \all{weil, dass, wenn}...)} : le pronom reste collé au verbe, tous deux à la fin : \all{..., weil ich mich nicht gut fühle.} (Présent) / \all{..., weil ich mich erkältet habe.} (Perfekt).
\end{enumerate}
\end{methode}

\begin{popfigure}
\begin{tikzpicture}[node distance=0.7cm, every node/.style={font=\small, align=center}]
\node (S) {Sujet};
\node (V) [right=of S, align=center]{Verbe conjugué\\(2\textsuperscript{e} pos.)};
\node (P) [right=of V, align=center]{Pronom réfléchi\\(acc./dat.)};
\node (C) [right=of P] {Compléments};
\node (Part) [right=of C, align=center]{Particule\\(si séparable)};
\draw[->, popblue, thick] (S) -- (V);
\draw[->, popblue, thick] (V) -- (P);
\draw[->, popblue, thick] (P) -- (C);
\draw[->, popblue, thick] (C) -- (Part);
\node[below=0.4cm of V, popmuted, text width=2.5cm] {Ex: \all{Ich | bewege | mich | nicht genug.}};
\node[below=0.4cm of Part, popmuted, text width=2.5cm] {Ex: \all{Wir | ruhen | uns | eine Stunde | aus.}};
\end{tikzpicture}
\legende{Structure de la phrase déclarative avec verbe pronominal (et particule séparable).}
\end{popfigure}

\begin{exoresolu}[Conjugaison : verbes pronominaux au présent]
\textbf{Énoncé.} Conjugue au présent et complète :
\begin{enumerate}
\item Ich (sich fühlen) heute nicht gut, ich habe Kopfschmerzen.
\item Nach dem Training (sich ausruhen) wir eine Stunde.
\item (sich erkälten) du oft im Regen?
\item Meine Schwester (sich waschen) die Haare. (Attention au cas !)
\end{enumerate}
\tcblower
\textbf{Raisonnement.} On conjugue d'abord le verbe selon la personne, puis on place le pronom réfléchi correspondant juste après le verbe. On vérifie la particule séparable et le cas (accusatif/datif).
\begin{enumerate}
\item Sujet \all{ich} $\rightarrow$ \all{fühle} + pronom accusatif \all{mich}.\\
\all{Ich fühle mich heute nicht gut, ich habe Kopfschmerzen.} « Je ne me sens pas bien aujourd'hui, j'ai mal à la tête. »
\item Sujet \all{wir} $\rightarrow$ \all{ruhen} + pronom \all{uns} + particule \all{aus} à la fin. Complément de temps \all{eine Stunde} avant la particule.\\
\all{Nach dem Training ruhen wir uns eine Stunde aus.} « Après l'entraînement, nous nous reposons une heure. »
\item Question sans mot interrogatif $\rightarrow$ verbe en 1\textsuperscript{re} position. Sujet \all{du} $\rightarrow$ \all{erkältest} + \all{dich}.\\
\all{Erkältest du dich oft im Regen?} « Est-ce que tu t'enrhumes souvent sous la pluie ? »
\item \all{sich waschen} + objet \all{die Haare} (accusatif) $\rightarrow$ pronom réfléchi au \textbf{DATIF} : \all{mir, dir, sich, uns, euch, sich}. Sujet \all{Meine Schwester} (3\textsuperscript{e} pers. sg) $\rightarrow$ \all{wäscht} + \all{sich} (datif = sich).\\
\all{Meine Schwester wäscht sich die Haare.} « Ma sœur se lave les cheveux. »
\end{enumerate}
\textbf{Conclusion.} Double piège : le pronom réfléchi (\all{mich, uns, dich, sich}) et la particule séparable (\all{aus} à la fin). Triple piège : le cas du pronom (datif si objet accusatif présent). Vérifie toujours les trois.
\end{exoresolu}

\courssub{Savoir-faire 4 : Donner un conseil avec « sollen » et « müssen »}

\begin{propriete}[Conjugaison des modaux \all{sollen} (conseil) et \all{müssen} (obligation) au présent]
\begin{center}
\begin{tabular}{|l|l|l|}
\hline \rowcolor{poporangeL} \textbf{Personne} & \textbf{\all{sollen}} & \textbf{\textbf{müssen}} \\ \hline
ich & \all{soll} & \all{muss} \\ \hline
du & \all{sollst} & \all{musst} \\ \hline
er/sie/es & \all{soll} & \all{muss} \\ \hline
wir & \all{sollen} & \all{müssen} \\ \hline
ihr & \all{sollt} & \all{müsst} \\ \hline
sie/Sie & \all{sollen} & \all{müssen} \\ \hline
\end{tabular}
\end{center}
\textbf{Particularités :} \all{sollen} : pas de modification de voyelle (régulier). \all{müssen} : Umlaut \all{ü} seulement aux pluriels (\all{wir/sie müssen, ihr müsst}), pas aux singuliers (\all{ich muss, du musst, er muss}).
\end{propriete}

\begin{methode}[Exprimer le conseil et l'obligation]
\begin{enumerate}
\item \textbf{Choisis le bon modal} : \all{sollen} = conseil, recommandation (avis extérieur, ex: médecin) ; \all{müssen} = obligation forte, nécessité intérieure ou contrainte extérieure.
\item \textbf{Conjugue le modal en 2\textsuperscript{e} position} (phrase principale).
\item \textbf{Mets l'infinitif du verbe lexical à la toute fin de la phrase} : \all{Du sollst mehr Sport treiben.}
\item \textbf{Avec un verbe pronominal} : l'infinitif garde son pronom réfléchi accordé avec le sujet : \all{Du musst dich mehr bewegen.}
\item \textbf{Nuance (Konjunktiv II - aperçu)} : pour un conseil poli/atténué : \all{sollten} (\all{Du solltest lieber zum Arzt gehen.}). \all{müssten} = obligation hypothétique.
\end{enumerate}
\end{methode}

\begin{popfigure}
\begin{tikzpicture}[node distance=0.7cm, every node/.style={font=\small, align=center}]
\node (S) {Sujet};
\node (M) [right=of S, align=center]{Modal conjugué\\(2\textsuperscript{e} pos.)};
\node (C) [right=of M] {Compléments};
\node (Pro) [right=of C, align=center]{Pronom réfl.\\(si verbe pron.)};
\node (Inf) [right=of Pro, align=center]{Infinitif\\(fin de phrase)};
\draw[->, poporange, thick] (S) -- (M);
\draw[->, poporange, thick] (M) -- (C);
\draw[->, poporange, thick] (C) -- (Pro);
\draw[->, poporange, thick] (Pro) -- (Inf);
\node[below=0.4cm of M, popmuted, text width=2.5cm] {Ex: \all{Du | sollst | mehr Wasser | trinken.}};
\node[below=0.4cm of Inf, popmuted, text width=3cm] {Ex pron.: \all{Du | musst | dich | mehr | bewegen.}};
\end{tikzpicture}
\legende{Structure de la phrase avec modal : l'infinitif (et son pronom réfléchi) à la fin.}
\end{popfigure}

\begin{exoresolu}[Transformation : donner un conseil]
\textbf{Énoncé.} Transforme chaque phrase en conseil avec \all{sollen} ou \all{müssen}, selon le sens indiqué.
\begin{enumerate}
\item \all{Peter trinkt viel Wasser.} (conseil du médecin)
\item \all{Ihr bewegt euch mehr.} (obligation)
\item \all{Maria geht zum Arzt.} (conseil)
\item \all{Wir essen gesund.} (obligation / nécessité)
\item \all{Du ruhst dich aus.} (conseil)
\end{enumerate}
\tcblower
\textbf{Raisonnement.} Le modal prend la 2\textsuperscript{e} position et se conjugue selon le sujet ; le verbe lexical passe à l'infinitif à la fin. Si verbe pronominal, le pronom s'accorde avec le sujet et précède l'infinitif.
\begin{enumerate}
\item Sujet \all{Peter} (il) $\rightarrow$ \all{soll}. Infinitif \all{trinken} à la fin.\\
\all{Peter soll viel Wasser trinken.} « Pierre devrait boire beaucoup d'eau. »
\item Sujet \all{ihr} $\rightarrow$ \all{müszt}. Verbe pronominal \all{sich bewegen} : pronom \all{euch} + infinitif \all{bewegen}.\\
\all{Ihr müsst euch mehr bewegen.} « Vous devez bouger davantage. »
\item Sujet \all{Maria} (elle) $\rightarrow$ \all{soll}. Infinitif \all{gehen} à la fin.\\
\all{Maria soll zum Arzt gehen.} « Maria devrait aller chez le médecin. »
\item Sujet \all{wir} $\rightarrow$ \all{müssen}. Infinitif \all{essen}.\\
\all{Wir müssen gesund essen.} « Nous devons manger sainement. »
\item Sujet \all{du} $\rightarrow$ \all{soll}. Verbe pronominal séparable \all{sich ausruhen} : pronom \all{dich} + infinitif \all{ausruhen}.\\
\all{Du sollst dich ausruhen.} « Tu devrais te reposer. »
\end{enumerate}
\textbf{Conclusion.} Structure immuable : \cle{Sujet + Modal conjugué (2\textsuperscript{e} pos.) + Compléments + Pronom réfléchi (si nécessaire) + Infinitif (fin)}.
\end{exoresolu}

\courssub{Savoir-faire 5 : Produire un court paragraphe sur la santé}

\begin{methode}[Rédiger un court texte guidé (5 à 8 phrases)]
\begin{enumerate}
\item \textbf{Prépare un mini-plan} : 1. Situation / Personne (Qui ?). 2. Problème (Maladie, manque de sport). 3. Cause (\all{weil} + verbe à la fin). 4. Conseil avec \all{sollen}. 5. Conseil avec \all{müssen}. 6. Conclusion positive (\all{deshalb / dann} + V2).
\item \textbf{Utilise les connecteurs simples} : \all{und, aber, denn} (verbe en 2\textsuperscript{e} position) ; \all{weil, dass} (verbe à la fin) ; \all{deshalb / dann} (verbe juste après : \all{Deshalb treibe ich Sport.}).
\item \textbf{Varie les structures imposées} : au moins une phrase avec verbe pronominal, une avec \all{sollen}, une avec \all{müssen}, vocabulaire thème santé/sport.
\item \textbf{Reste simple et correct} : mieux vaut des phrases courtes sans faute qu'une phrase longue mal construite.
\item \textbf{Relis avec la grille de relecture} : Verbe en 2\textsuperscript{e} position ? Infinitif à la fin avec modal ? Majuscules aux noms ? Pronom réfléchi accordé ? Cas (acc/dat) après prépositions (\all{für, mit, zu}) ?
\end{enumerate}
\end{methode}

\begin{exoresolu}[Rédaction : « Mein Freund ist oft krank »]
\textbf{Énoncé.} Ton ami Paul est souvent malade. Écris un court texte (6 à 7 phrases) en allemand : décris son problème, la cause, donne-lui des conseils avec \all{sollen} et \all{müssen}, conclus.
\tcblower
\textbf{Raisonnement.} On suit le mini-plan. On vérifie la place du verbe dans chaque phrase.\\
\textbf{Production modèle :}\\
\all{Mein Freund Paul ist oft krank. Er bewegt sich nicht genug und treibt keinen Sport. Er isst auch oft ungesund. Deshalb hat er oft Fieber und Kopfschmerzen. Er soll jeden Tag dreißig Minuten Sport machen. Er muss auch gesund essen und viel Wasser trinken. Ich sage ihm: „Sport ist die beste Medizin." Dann fühlt er sich bald besser.}\\
« Mon ami Paul est souvent malade. Il ne bouge pas assez et ne fait pas de sport. Il mange aussi souvent malsain. C'est pourquoi il a souvent de la fièvre et mal à la tête. Il devrait faire trente minutes de sport chaque jour. Il doit aussi manger sainement et boire beaucoup d'eau. Je lui dis : „Le sport est la meilleure médecine.“ Alors il se sentira bientôt mieux. »\\
\textbf{Justifications grammaticales :}
\begin{itemize}
\item \all{bewegt sich nicht genug} (verbe pronominal, pronom après verbe, V2).
\item \all{treibt keinen Sport} (accusatif masculin négatif \all{keinen}).
\item \all{Deshalb hat er...} (connecteur \all{deshalb} $\rightarrow$ V2 \all{hat}).
\item \all{soll ... machen} / \all{muss ... trinken} (modal V2, infinitif fin).
\item \all{sage ihm} (verbe \all{sagen} + datif \all{ihm}).
\item \all{Dann fühlt er sich...} (connecteur \all{dann} $\rightarrow$ V2 \all{fühlt}, pronom \all{sich} après).
\end{itemize}
\end{exoresolu}

\begin{savaistu}[Santé et sport : regards croisés Cameroun — Allemagne]
\textbf{Allemagne (DACH) :} Le système de santé (\all{das Gesundheitssystem}) est basé sur l'assurance maladie obligatoire (\all{die Krankenkasse}). Le médecin de famille (\all{der Hausarzt}) est la première porte d'entrée. Le sport est très institutionnalisé : clubs (\all{der Verein}), salles de sport (\all{das Fitnessstudio}), pistes cyclables. La prévention (\all{die Vorsorge}) est remboursée (cures thermales \all{die Kur}).
\textbf{Cameroun :} Accès aux soins plus difficile (coût, distance). Médecine traditionnelle (\all{die traditionelle Medizin}) très présente. Sport roi : le football (\all{Fußball}) sur terrains vagues ou stades (Omnisport Yaoundé, Japoma Douala). Marche à pied quotidienne fréquente par nécessité. Montée du diabète/hypertension (\all{Diabetes, Bluthochdruck}) due à l'urbanisation et l'alimentation grasse/sucrée.
\textbf{Point commun :} L'OMS recommande 150 min/semaine d'activité modérée. \all{Bewegung ist Leben!} — \all{Bewegung ist die beste Medizin!}
\end{savaistu}

\begin{experience}[Jeu de rôle : « Beim Arzt » (Chez le médecin)]
\textbf{Consigne :} Par binôme. A = Patient (malade), B = Médecin.
\textbf{Cartes de rôle :}
\begin{itemize}
\item \textbf{Patient :} Tu as \all{Fieber}, \all{Kopfschmerzen}, \all{Halsschmerzen} depuis 3 jours. Tu ne fais pas de sport. Tu manges \all{ungesund} (Fast Food). Tu demandes un conseil (\all{Was soll ich tun?}).
\item \textbf{Médecin :} Tu examines. Tu donnes diagnostic : \all{Sie haben eine Erkältung / Grippe}. Tu donnes ordonnance : \all{Medikamente nehmen}. Tu donnes conseils mode de vie avec \all{sollen/müssen} : \all{Sie sollen sich ausruhen. Sie müssen viel Wasser trinken. Sie sollen Sport treiben.}
\end{itemize}
\textbf{Expressions utiles :} \all{Was fehlt Ihnen?} (Qu'avez-vous ?), \all{Ich fühle mich schlecht.} (Je me sens mal.), \all{Seit wann?} (Depuis quand ?), \all{Gute Besserung!} (Bon rétablissement !).
\end{experience}

\begin{attention}[Les 5 erreurs qui coûtent des points au Bac]
\begin{enumerate}
\item \textbf{Oublier le pronom réfléchi ou mal l'accorder.} Faux : \all{Ich fühle heute besser.} / \all{Wir bewegen mich.} Correct : \all{Ich fühle mich heute besser.} / \all{Wir bewegen uns.} Le pronom s'accorde \textbf{toujours} avec le sujet (\all{ich$\rightarrow$mich, du$\rightarrow$dich, er$\rightarrow$sich, wir$\rightarrow$uns, ihr$\rightarrow$euch, sie$\rightarrow$sich}).
\item \textbf{Infinitif mal placé avec les modaux.} Faux : \all{Du sollst trinken viel Wasser.} Correct : \all{Du sollst viel Wasser trinken.} L'infinitif va à la \textbf{toute dernière place}, après tous les compléments.
\item \textbf{Verbe mal placé après « weil / dass / wenn ».} Faux : \all{..., weil er ist krank.} Correct : \all{..., weil er krank ist.} Après les conjonctions de subordination, le verbe conjugué va à la \textbf{fin}.
\item \textbf{Oublier la majuscule aux noms.} Faux : \all{der sport ist gut für die gesundheit.} Correct : \all{Der Sport ist gut für die Gesundheit.} \textbf{Tous} les noms prennent la majuscule, même en milieu de phrase.
\item \textbf{Calquer le français (Faux amis / Structures).}
\begin{itemize}
\item « Avoir de la fièvre » $\neq$ \all{haben das Fieber} $\rightarrow$ \cle{\all{Fieber haben}} (sans article).
\item « Aller chez le médecin » $\neq$ \all{gehen zu der Arzt} $\rightarrow$ \cle{\all{zum Arzt gehen}} (\all{zu} + datif \all{dem} $\rightarrow$ \all{zum}).
\item \all{die Krankheit} = la maladie (pas « la crainte» $\rightarrow$ \all{die Angst}).
\item \all{das Gift} = le poison (pas « le cadeau » $\rightarrow$ \all{das Geschenk}).
\item \all{der Sport} = le sport (pas « le soutien » $\rightarrow$ \all{die Unterstützung}).
\end{itemize}
\end{enumerate}
\end{attention}

\newpage

\coursec{Exercices}

\coursec{Les maladies et l'importance de la pratique du sport}

\courssub{Texte support : „Sport ist die beste Medizin“}

\begin{textebox}[Article original : Santé et mouvement au quotidien]
Bewegung hält gesund – das wissen alle. Doch in Yaoundé wie in Berlin sitzen viele Jugendliche stundenlang vor dem Handy. Die Folge: Rückenschmerzen, Übergewicht und häufige Erkältungen. Dr. Meyer, ein Hausarzt in Köln, erklärt: „Wer sich nicht bewegt, schwächt sein Immunsystem. Sport ist die beste Medizin, ohne Nebenwirkungen.“ Er rät: „Man soll sich 30 Minuten am Tag bewegen: joggen, schwimmen oder einfach schnell gehen.“ Wichtig ist auch die Erholung: „Man muss sich nach dem Training ausruhen und genug schlafen.“ In Kamerun starten Schulen Programme wie „Sport pour la Santé“. Beim Bendskin-Tanzen oder Fußballspielen auf dem Quartier-Platz bewegen sich die Jugendlichen mit Freude. Wer sich regelmäßig bewegt, fühlt sich besser und lernt leichter. Also: Handy weg, Sportschuhe an!
\end{textebox}

\begin{exemplebox}[Glossaire du texte]
\begin{itemize}
    \item \all{die Bewegung, -en} : le mouvement, l'exercice physique.
    \item \all{das Immunsystem, -e} : le système immunitaire.
    \item \all{die Nebenwirkung, -en} : l'effet secondaire.
    \item \all{raten (rät, riet, geraten)} : conseiller, recommander.
    \item \all{sich bewegen} : bouger, faire de l'exercice.
    \item \all{sich ausruhen} : se reposer.
    \item \all{sich fühlen} : se sentir (état de santé).
    \item \all{der Quartier-Platz, -e} : la place du quartier (terrain vague).
    \item \all{regelmäßig} : régulièrement.
\end{itemize}
\end{exemplebox}

\courssub{Compréhension du texte}

\begin{exercice}[Compréhension globale et détaillée]
Répondez en allemand par des phrases complètes.
\begin{enumerate}
    \item Was sind die Folgen von Bewegungsmangel nach dem Text? (Nennen Sie drei.)
    \item Was empfiehlt Dr. Meyer als „beste Medizin“?
    \item Wie viel Bewegung pro Tag rät der Arzt an?
    \item Nennen Sie zwei Beispiele für Sportarten aus dem Text.
    \item Was ist nach dem Training genauso wichtig wie die Bewegung selbst?
    \item Wie heißt das Programm für mehr Gesundheit in kamerunischen Schulen?
    \item Warum bewegen sich Jugendliche in Yaoundé gerne? (Nennen Sie zwei Aktivitäten.)
\end{enumerate}
\end{exercice}

\courssub{Lexique thématique : Santé, maladie et sport}

\begin{propriete}[Familles de mots et derivation]
En allemand, les mots se construisent souvent par \textbf{composition} (Nomen + Nomen) ou \textbf{derivation} (préfixes/suffixes). Apprenez le nom \textbf{avec son article et son pluriel}, puis ses dérivés.
\begin{center}
\begin{tabularx}{\linewidth}{|X|X|X|X|}
\hline
\rowcolor{popblueL} \textbf{Nom (Singulier)} & \textbf{Article} & \textbf{Pluriel} & \textbf{Dérivés / Composés} \\ \hline
\all{Krankheit} & \all{die} & \all{die Krankheiten} & \all{krank (adj.), erkranken (v.), das Krankenzimmer} \\ \hline
\all{Sport} & \all{der} & \textendash\ (kein Plural) & \all{sportlich (adj.), der Sportplatz, die Sportart, treiben Sport} \\ \hline
\all{Arzt} & \all{der} & \all{die Ärzte} & \all{ärztlich (adj.), die Arztpraxis, der Zahnarzt, die Ärztin} \\ \hline
\all{Fieber} & \all{das} & \textendash\ (kein Plural) & \all{fiebern (v.), das Fieberthermometer, fieberhaft (adj.)} \\ \hline
\all{Verletzung} & \all{die} & \all{die Verletzungen} & \all{sich verletzen (v.), verletzt (adj.), die Sportverletzung} \\ \hline
\all{Training} & \all{das} & \all{die Trainings} & \all{trainieren (v.), der Trainingsplan, das Trainingsspiel} \\ \hline
\all{Ruhe} & \all{die} & \textendash\ (kein Plural) & \all{ruhen (v.), sich ausruhen (v.), der Ruhetag, ruhig (adj.)} \\ \hline
\all{Medizin} & \all{die} & \all{die Medizinen} & \all{medizinisch (adj.), das Medikament, die Medizinstudentin} \\ \hline
\all{Ernährung} & \all{die} & \textendash & \all{sich ernähren (v.), gesund (adj.), das Nahrungsmittel} \\ \hline
\all{Schlaf} & \all{der} & \all{die Schlafe} & \all{schlafen (v.), der Schlafmangel, schläfrig (adj.)} \\ \hline
\end{tabularx}
\end{center}
\end{propriete}

\begin{attention}[Faux amis et pièges fréquents]
\begin{itemize}
    \item \all{der Arzt} (médecin) $\neq$ \all{der Art} (le genre, l'espèce).
    \item \all{das Gift} = le poison ! Pour « cadeau », dire \all{das Geschenk}.
    \item \all{sich erkälten} = attraper un rhume (refroidissement). \all{erkalten} (sans \textit{sich}) = refroidir (température).
    \item \all{trainieren} = s'entraîner (sport). Pour « former quelqu'un », on utilise \all{ausbilden}.
    \item Majuscules : \textbf{Tous les noms} prennent une majuscule (\all{der Sport}, \all{die Gesundheit}, \all{das Fieber}).
\end{itemize}
\end{attention}

\begin{exercice}[Vocabulaire : Articles, pluriels et familles de mots]
Complétez le tableau. Attention aux Umlauts au pluriel (\all{Arzt $\rightarrow$ Ärzte}) et aux genres.
\begin{center}
\begin{tabularx}{\linewidth}{|X|X|X|X|}
\hline
\rowcolor{popblueL} \textbf{Nom (Singulier)} & \textbf{Article} & \textbf{Pluriel} & \textbf{Dérivé / Composé} \\ \hline
Krankheit & & & \\ \hline
Sport & & \textendash & \\ \hline
Arzt & & & \\ \hline
Fieber & & \textendash & \\ \hline
Verletzung & & & \\ \hline
Training & & & \\ \hline
Ruhe & & \textendash & \\ \hline
Medizin & & & \\ \hline
Ernährung & & \textendash & \\ \hline
Schlaf & & & \\ \hline
\end{tabularx}
\end{center}
\end{exercice}

\courssub{Grammaire : Les verbes pronominaux (Reflexivverben)}

\begin{propriete}[Formation et emploi]
Un \textbf{verbe pronominal} ( \all{Reflexivverb} ) s'accompagne d'un \textbf{pronom réfléchi} (\all{sich} à l'infinitif) qui renvoie au sujet.
\begin{itemize}
    \item \textbf{Accusatif (le plus fréquent) :} Action que le sujet fait \emph{sur lui-même}.
    \all{sich waschen, sich beeilen, sich verletzen, sich erkälten, sich fühlen, sich bewegen, sich erholen, sich anziehen}.
    \item \textbf{Datif :} Action faite \emph{pour soi-même} ou sur une \emph{partie du corps} (le corps devient l'objet direct/accusatif).
    \all{sich die Hände waschen, sich die Zähne putzen, sich etwas kaufen, sich vorstellen}.
\end{itemize}
\textbf{Conjugaison au Présent (ex. \textit{sich fühlen} -- Akk., \textit{sich die Hände waschen} -- Dat.) :}
\begin{center}
\begin{tabular}{|l|l|l|}
\hline
\rowcolor{popblueL} \textbf{Personne} & \textbf{Akkusativ (sich fühlen)} & \textbf{Datif (sich die Hände waschen)} \\ \hline
ich & ich fühle \textbf{mich} & ich wasche \textbf{mir} die Hände \\ \hline
du & du fühlst \textbf{dich} & du wäschst \textbf{dir} die Hände \\ \hline
er/sie/es & er fühlt \textbf{sich} & er wäscht \textbf{sich} die Hände \\ \hline
wir & wir fühlen \textbf{uns} & wir waschen \textbf{uns} die Hände \\ \hline
ihr & ihr fühlt \textbf{euch} & ihr wascht \textbf{euch} die Hände \\ \hline
sie/Sie & sie fühlen \textbf{sich} & sie waschen \textbf{sich} die Hände \\ \hline
\end{tabular}
\end{center}
\textbf{Place du pronom :} Derrière le verbe conjugué (phrase principale), derrière le sujet (subordonnée), avant l'infinitif/groupe verbal (\all{Er muss sich ausruhen}, \all{Er will sich waschen}).
\end{propriete}

\begin{attention}[Verbes à sens changeant / Pièges pour francophones]
\begin{itemize}
    \item \all{sich erinnern an (+ Akk)} = se souvenir de. \textbf{Pas} \all{an sich erinnern}.
    \item \all{sich interessieren für (+ Akk)} = s'intéresser à.
    \item \all{sich freuen auf (+ Akk)} = se réjouir de (futur) ; \all{sich freuen über (+ Akk)} = se réjouir de (passé/présent).
    \item \textbf{Particule séparable :} \all{sich ausruhen, sich anziehen, sich erkälten} $\rightarrow$ \all{Er ruht sich aus.} (Particule \all{aus} à la fin).
    \item \textbf{Corps :} \all{Ich wasche mir die Hände} (Datif \all{mir} !). \textit{Français : "Je me lave les mains" $\rightarrow$ Allemand : "Je me (dat) lave les mains (akk)".}
\end{itemize}
\end{attention}

\begin{exemplebox}[Verbes pronominaux courants du thème (Niveau A1/A2)]
\begin{center}
\begin{tabularx}{\linewidth}{|l|X|l|}
\hline
\rowcolor{popgreenL} \textbf{Verbe (Infinitif)} & \textbf{Sens} & \textbf{Cas} \\ \hline
\all{sich fühlen} & se sentir (santé) & Akk \\ \hline
\all{sich erkälten} & attraper un rhume & Akk \\ \hline
\all{sich verletzen} & se blesser & Akk \\ \hline
\all{sich bewegen} & bouger, faire du sport & Akk \\ \hline
\all{sich ausruhen} & se reposer & Akk \\ \hline
\all{sich erholen} & se remettre, récupérer & Akk \\ \hline
\all{sich beeilen} & se dépêcher & Akk \\ \hline
\all{sich waschen / \dots die Hände waschen} & se laver / se laver les mains & Akk / Dat \\ \hline
\all{sich ernähren} & s'alimenter & Akk \\ \hline
\all{sich impfen lassen} & se faire vacciner & Akk \\ \hline
\end{tabularx}
\end{center}
\end{exemplebox}

\courssub{Grammaire : Le conseil avec „sollen“ und „müssen“}

\begin{propriete}[Les verbes modaux de conseil et d'obligation]
\begin{center}
\begin{tabularx}{\linewidth}{|X|X|X|}
\hline
\rowcolor{poporangeL} \textbf{Verbe} & \textbf{Valeur / Nuance} & \textbf{Exemple} \\ \hline
\all{sollen} & \textbf{Conseil, recommandation} (avis extérieur, médecin, règle morale). \textit{« Tu devrais... »} & \all{Du sollst mehr Sport machen.} \\ \hline
\all{müssen} & \textbf{Obligation, nécessité forte} (règle, loi, urgence, contrainte logique). \textit{« Tu dois... »} & \all{Du musst zum Arzt gehen.} \\ \hline
\all{nicht müssen} & \textbf{Absence de nécessité} (pas l'obligation de ne pas faire !). \textit{« Tu n'as pas besoin de... »} & \all{Du musst keine Tabletten nehmen.} \\ \hline
\all{nicht sollen} & \textbf{Déconseil} (il est préférable de ne pas). \textit{« Tu ne devrais pas... »} & \all{Du sollst nicht rauchen.} \\ \hline
\end{tabularx}
\end{center}
\textbf{Structure de phrase :} Sujet + \textbf{Modalverbe conjugué (2ème place)} + Compléments + \textbf{Infinitif à la fin}.
\all{Du sollst [heute] [mehr Wasser] trinken.}
\all{Wir müssen [uns] [jetzt] ausruhen.} (Pronom réfléchi \textbf{avant} l'infinitif).
\end{propriete}

\begin{methode}[Traduire un conseil (Français $\rightarrow$ Allemand)]
\begin{enumerate}
    \item Identifier la nuance : \textbf{Conseil/Avis} $\rightarrow$ \all{sollen} ; \textbf{Obligation/Urgent} $\rightarrow$ \all{müssen}.
    \item Conjuguer le modal au sujet (je $\rightarrow$ \all{soll/muss}, tu $\rightarrow$ \all{sollst/musst}, il $\rightarrow$ \all{soll/muss}, nous $\rightarrow$ \all{sollen/müssen}, vous $\rightarrow$ \all{sollt/müsst}, ils/Vous(poli) $\rightarrow$ \all{sollen/müssen}).
    \item Placer le sujet en position 1, le modal en position 2.
    \item Si verbe pronominal : mettre le pronom réfléchi \textbf{juste avant l'infinitif} (\all{sich ausruhen $\rightarrow$ ... uns ausruhen}).
    \item Vérifier l'ordre : \textbf{Sujet - Modal - (Compléments) - Pronom Réfléchi - Infinitif}.
\end{enumerate}
\end{methode}

\begin{exemplebox}[Exemples contrastés]
\all{Du sollst mehr Gemüse essen.} « Tu devrais manger plus de légumes. » (Conseil diététique)
\all{Du musst mehr Gemüse essen.} « Tu dois manger plus de légumes. » (Ordre du médecin / Règle stricte)
\all{Er soll sich nicht so sehr anstrengen.} « Il ne devrait pas trop forcer. » (Déconseil)
\all{Er muss sich nicht anstrengen.} « Il n'a pas besoin de forcer. » (Absence de nécessité)
\end{exemplebox}

\courssub{Conjugaison : Verbes pronominaux au Présent}

\begin{exercice}[Conjugaison guidée]
Conjuguez les verbes pronominaux entre parenthèses au \textbf{Präsens}. Placez le pronom réfléchi au bon cas (accusatif ou datif).
\begin{enumerate}
    \item Ich \_\_\_\_\_\_\_\_\_\_\_\_ heute nicht gut. (sich fühlen -- Akk)
    \item Warum \_\_\_\_\_\_\_\_\_\_\_\_ du \_\_\_\_\_\_\_\_\_\_\_\_? (sich beeilen -- Akk)
    \item Der Spieler \_\_\_\_\_\_\_\_\_\_\_\_ am Knie. (sich verletzen -- Akk)
    \item Wir \_\_\_\_\_\_\_\_\_\_\_\_ nach dem Training eine Stunde. (sich ausruhen -- Akk, séparable)
    \item Meine Schwester \_\_\_\_\_\_\_\_\_\_\_\_ schnell \_\_\_\_\_\_\_\_\_\_\_\_. (sich erkälten -- Akk)
    \item \_\_\_\_\_\_\_\_\_\_\_\_ Sie \_\_\_\_\_\_\_\_\_\_\_\_ gut? (sich erholen -- Akk, formel \textit{Sie})
    \item Die Kinder \_\_\_\_\_\_\_\_\_\_\_\_ im Garten. (sich bewegen -- Akk)
    \item Er \_\_\_\_\_\_\_\_\_\_\_\_ die Hände vor dem Essen. (sich waschen -- \textbf{Datif} car partie du corps)
\end{enumerate}
\end{exercice}

\courssub{Entraînement : Traduction, transformation et texte à trous}

\begin{exercice}[Traduction : Conseils avec \textit{sollen} et \textit{müssen}]
Traduisez en allemand. Respectez l'ordre des mots (verbe modal 2ème place, infinitif à la fin, pronom réfléchi avant l'infinitif).
\begin{enumerate}
    \item Tu devrais faire plus de sport. (\textit{mehr Sport machen})
    \item Il doit aller chez le médecin. (\textit{zum Arzt gehen})
    \item Nous devons nous reposer maintenant. (\textit{sich ausruhen})
    \item Vous ne devriez pas (pluriel poli) travailler si vous avez de la fièvre. (\textit{arbeiten}, \textit{Fieber haben})
    \item Je dois me laver les mains. (\textit{sich die Hände waschen} -- attention au cas !)
    \item Elle devrait boire plus d'eau. (\textit{mehr Wasser trinken})
\end{enumerate}
\end{exercice}

\begin{exercice}[Transformation : Impératif $\rightarrow$ Conseil avec \textit{sollen}/\textit{müssen}]
Transformez les impératifs en phrases de conseil. Choisissez \all{sollen} (recommandation) ou \all{müssen} (obligation/nécessité forte). Gardez le même sujet.
\begin{enumerate}
    \item Mach Sport! $\rightarrow$ Du \_\_\_\_\_\_\_\_\_\_\_\_ Sport machen.
    \item Geh zum Arzt! $\rightarrow$ Er \_\_\_\_\_\_\_\_\_\_\_\_ zum Arzt gehen.
    \item Ruht euch aus! $\rightarrow$ Ihr \_\_\_\_\_\_\_\_\_\_\_\_ euch ausruhen.
    \item Wasch dir die Hände! $\rightarrow$ Wir \_\_\_\_\_\_\_\_\_\_\_\_ uns die Hände waschen.
    \item Trinke mehr Wasser! $\rightarrow$ Sie \_\_\_\_\_\_\_\_\_\_\_\_ mehr Wasser trinken.
\end{enumerate}
\end{exercice}

\begin{exercice}[Texte à trous : Une visite chez le médecin]
Complétez le dialogue avec les formes correctes (Présent, Perfekt pour le passé, verbes modaux). Attention aux verbes à particule séparable et aux pronoms réfléchis (cas !).
\begin{textebox}[Dialog bei Dr. Müller]
Arzt: Guten Tag, Herr Nkoto. Was fehlt Ihnen?
Patient: Ich \_\_\_\_\_\_\_\_\_\_\_\_ (fühlen) mich sehr schlapp und ich \_\_\_\_\_\_\_\_\_\_\_\_ (haben) seit gestern Fieber.
Arzt: Haben Sie sich \_\_\_\_\_\_\_\_\_\_\_\_ (verletzen -- Perfekt) beim Sport?
Patient: Nein, aber ich \_\_\_\_\_\_\_\_\_\_\_\_ (trainieren) seit drei Wochen nicht mehr. Ich \_\_\_\_\_\_\_\_\_\_\_\_ (haben) viel Stress im Büro.
Arzt: Das ist das Problem. Sie \_\_\_\_\_\_\_\_\_\_\_\_ (müssen) sich mehr \_\_\_\_\_\_\_\_\_\_\_\_ (bewegen). Und Sie \_\_\_\_\_\_\_\_\_\_\_\_ (sollen) Stress \_\_\_\_\_\_\_\_\_\_\_\_ (vermeiden).
Patient: Soll ich Medikamente nehmen?
Arzt: Nein, Sie \_\_\_\_\_\_\_\_\_\_\_\_ (brauchen) keine Tabletten. Sie \_\_\_\_\_\_\_\_\_\_\_\_ (müssen) sich nur \_\_\_\_\_\_\_\_\_\_\_\_ (erholen) und regelmäßig \_\_\_\_\_\_\_\_\_\_\_\_ (laufen) oder \_\_\_\_\_\_\_\_\_\_\_\_ (schwimmen).
\end{textebox}
\end{exercice}

\begin{exercice}[Appariement : Mots composés et expressions]
Reliez la colonne A à la colonne B. Écrivez les paires (ex. 1--b) et formez une phrase complète en allemand pour chaque paire.
\begin{center}
\begin{tabularx}{\linewidth}{|X|X|}
\hline
\rowcolor{poporangeL} \textbf{A} & \textbf{B} \\ \hline
1. die Sportart & a) eine Krankheit, die durch Viren entsteht \\ \hline
2. der Hausarzt & b) z. B. Fußball, Schwimmen, Joggen \\ \hline
3. die Grippe (Influenza) & c) der Arzt für die Grundversorgung \\ \hline
4. das Rezept & d) man macht es, um fit zu werden \\ \hline
5. das Training & e) der Arzt schreibt es für die Apotheke \\ \hline
6. sich impfen lassen & f) Schutz vor Krankheiten (Prävention) \\ \hline
\end{tabularx}
\end{center}
\end{exercice}

\courssub{Préparation au Baccalauréat}

\begin{exercice}[Compréhension de texte type Baccalauréat]
Lisez le texte suivant attentivement, puis répondez aux questions en allemand (1--5) et en français (6--8).

\begin{textebox}[Jugend und Gesundheit in Deutschland]
Laut einer aktuellen Studie der Bundeszentrale für gesundheitliche Aufklärung (BZgA) treiben nur noch 28 Prozent der Jugendlichen in Deutschland ausreichend Sport. Die Weltgesundheitsorganisation (WHO) empfiehlt mindestens 60 Minuten Bewegung pro Tag. Doch viele Schüler sitzen stundenlang vor dem Smartphone oder der Spielkonsole. Die Folge: Übergewicht, Rückenschmerzen und häufige Erkältungen nehmen zu. „Kinder müssen sich wieder mehr bewegen“, warnt Prof. Dr. Petra Wagner, Sportmedizinerin an der Universität Köln. „Wer sich in der Jugend nicht fit hält, hat als Erwachsener oft Herz-Kreislauf-Probleme.“

Schulen und Vereine versuchen gegenzusteuern. In Nordrhein-Westfalen gibt es das Programm „Fit für die Schule“, das tägliche Bewegungspausen vorsieht. Auch Eltern sind gefragt: „Eltern sollen Vorbilder sein“, sagt Wagner. „Wenn die Eltern joggen oder Rad fahren, machen die Kinder oft mit.“ Doch der Alltag ist stressig. Viele Familien haben keine Zeit für gemeinsame Aktivitäten. Experten fordern daher mehr Schwimmunterricht und sichere Radwege, damit sich Kinder selbstständig bewegen können. Gesundheit fängt nicht in der Arztpraxis an, sondern im Alltag.
\end{textebox}

\textbf{Fragen auf Deutsch (Antworten Sie in ganzen Sätzen):}
\begin{enumerate}
    \item Wie viel Prozent der Jugendlichen treiben laut Studie ausreichend Sport?
    \item Was empfiehlt die WHO für Jugendliche?
    \item Nennen Sie drei negative Folgen von Bewegungsmangel, die im Text stehen.
    \item Was schlägt Prof. Wagner den Eltern vor?
    \item Welche zwei Maßnahmen fordern Experten für die Infrastruktur?
\end{enumerate}

\textbf{Questions en français :}
\begin{enumerate}
    \setcounter{enumi}{5}
    \item Expliquez l'expression „Gesundheit fängt nicht in der Arztpraxis an, sondern im Alltag“ en vous appuyant sur le texte.
    \item Comparez la situation décrite en Allemagne avec celle que vous observez au Cameroun (écoles, quartiers, habitudes des jeunes). Donnez deux similitudes et deux différences.
    \item Trouvez dans le texte deux verbes pronominaux au présent et deux verbes modaux (\textit{müssen, sollen, können, wollen, dürfen, mögen}). Citez-les avec leur sujet.
\end{enumerate}
\end{exercice}

\begin{exercice}[Thème et Version]

\textbf{A. Thème (Français $\rightarrow$ Allemand)}
Traduisez en allemand. Soignez l'ordre des mots (place du verbe, place du pronom réfléchi) et l'orthographe (majuscules, Umlaute, ß).
\begin{enumerate}
    \item Mon ami Paul est souvent malade. Il a attrapé un rhume hier.
    \item Il doit aller chez le médecin et il doit se reposer.
    \item Le médecin lui dit : « Tu devrais faire du sport chaque semaine. »
    \item Paul répond : « Je n'ai pas de temps, je dois beaucoup apprendre. »
    \item Mais sans sport, il ne se sentira pas mieux.
\end{enumerate}

\textbf{B. Version (Allemand $\rightarrow$ Français)}
Traduisez en français le paragraphe suivant.
\begin{textebox}[Version]
Lisa fühlt sich heute Morgen nicht gut. Sie hat Fieber und Halsschmerzen. Sie muss im Bett bleiben und sich ausruhen. Ihre Mutter ruft den Arzt an. Der Arzt kommt und sagt: „Lisa soll viel Tee trinken und sich nicht anstrengen. Sie darf morgen nicht in die Schule gehen.“ Lisa ist traurig, aber sie weiß: Gesundheit ist das Wichtigste. Sie will sich bald erholen, um wieder Fußball zu spielen.
\end{textebox}
\end{exercice}

\begin{exercice}[Production écrite : Lettre de conseils (60--80 Wörter)]
\textbf{Consigne :} Votre ami(e) allemand(e) Alex / Anna vous écrit qu'il/elle est souvent fatigué(e) et malade ces derniers temps. Il/Elle ne fait pas de sport et mange mal. Écrivez-lui une lettre (ou un e-mail) pour lui donner des conseils concrets pour améliorer sa santé.

\textbf{Contraintes :}
\begin{itemize}
    \item Longueur : \textbf{60 à 80 mots} (comptez-les).
    \item Utilisez au moins \textbf{3 verbes pronominaux} différents (\textit{sich fühlen, sich ausruhen, sich bewegen, sich ernähren, sich verletzen, sich erkälten...}).
    \item Utilisez \textbf{sollen} et \textbf{müssen} au moins une fois chacun.
    \item Structure de lettre : Salutation, corps du texte (conseils), formule de politesse, signature.
    \item Vocabulaire du thème : \textit{das Gemüse, das Obst, der Schlaf, der Bildschirm, der Verein, joggen, schwimmen}.
\end{itemize}

\begin{textebox}[Modèle de début]
Lieber Alex / Liebe Anna,

ich habe gehört, dass du dich oft \_\_\_\_\_\_\_\_\_\_\_\_.
\end{textebox}
\end{exercice}

\courssub{Corrigés des exercices}

\begin{popfigure}
\begin{tikzpicture}[
    node distance=1.2cm and 1.5cm,
    box/.style={draw=popblue, thick, rounded corners, fill=popblueL, text width=3.2cm, align=center, font=\small},
    arrow/.style={-Stealth, thick, popdark}
]
\node[box, align=center] (subj){Sujet\\ (Ich, Du, Er...)};
\node[box, right=of subj, align=center] (modal){Verbe Modal\\ \textbf{sollen / müssen}\\ (Position 2)};
\node[box, right=of modal, align=center] (comp){Compléments\\ (Zeit, Ort, Objet)};
\node[box, right=of comp, align=center] (refl){Pronom Réfléchi\\ \textbf{mich, dich, sich,\\ uns, euch, sich}\\ (si verbe pron.)};
\node[box, right=of refl, align=center] (inf){Infinitif\\ \textbf{à la fin}\\ (machen, gehen,\\ sich ausruhen)};
\draw[arrow] (subj) -- (modal);
\draw[arrow] (modal) -- (comp);
\draw[arrow] (comp) -- (refl);
\draw[arrow] (refl) -- (inf);
\node[below=0.3cm of modal, font=\tiny\itshape, text=popmuted] {conjugué};
\node[below=0.3cm of inf, font=\tiny\itshape, text=popmuted] {non conjugué};
\end{tikzpicture}
\legende{Structure de la phrase avec verbes modaux + verbes pronominaux}
\end{popfigure}

\begin{savaistu}[Le saviez-vous ? Santé et culture]
En Allemagne, le \textbf{« Hausarzt »} (médecin de famille) est la première porte d'entrée du système de santé (Prinzip der Hausarztzentrierten Versorgung). On ne va pas directement chez le spécialiste sans « Überweisung » (lettre de recommandation). Au Cameroun, on consulte souvent directement un spécialiste ou un médecin de district. Le mot \textbf{« Rezept »} (ordonnance) vient du latin \textit{recipe} (« prends ! ») — c'était l'en-tête des ordonnances anciennes. En pharmacie (\textbf{Apotheke}), les médicaments sont souvent chers, mais la caisse d'assurance maladie (\textbf{Krankenkasse}) rembourse une grande partie. Le sport scolaire (\textbf{Schulsport}) est obligatoire jusqu'au Bac ; les notes comptent pour la moyenne !
\end{savaistu}

\newpage

\coursec{Corrigés des exercices}

\begin{corrige}[Exercice 1 — Compréhension du texte support]
\begin{enumerate}
    \item Die Folgen von Bewegungsmangel sind: \textbf{Rückenschmerzen, Übergewicht und häufige Erkältungen}. (Reprise directe du texte : „Die Folge: Rückenschmerzen, Übergewicht und häufige Erkältungen.“)
    \item Dr. Meyer empfiehlt \textbf{Sport als die beste Medizin} (ohne Nebenwirkungen).
    \item Er rät zu \textbf{30 Minuten Bewegung am Tag} (joggen, schwimmen oder einfach schnell gehen).
    \item Zwei Beispiele für Sportarten: \textbf{Joggen und Schwimmen} (également acceptés : schnelles Gehen, Bendskin-Tanzen, Fußballspielen).
    \item Nach dem Training ist \textbf{die Erholung} genauso wichtig: Man muss sich ausruhen und genug schlafen.
    \item Das Programm heißt \textbf{„Sport pour la Santé“}.
    \item Die Jugendlichen bewegen sich gerne \textbf{beim Bendskin-Tanzen und beim Fußballspielen auf dem Quartier-Platz} — mit Freude.
\end{enumerate}
\textbf{Points de vigilance :} répondre par des phrases complètes avec sujet + verbe en position 2 ; ne pas recopier des fragments sans verbe ; \all{raten} se conjugue \all{er rät} (verbe fort à voyelle changeante \all{a $\rightarrow$ ä}).
\end{corrige}

\begin{corrige}[Exercice 2 — Vocabulaire : articles, pluriels et familles de mots]
\begin{center}
\begin{tabularx}{\linewidth}{|X|l|X|X|}
\hline
\rowcolor{popblueL} \textbf{Nom} & \textbf{Article} & \textbf{Pluriel} & \textbf{Dérivé / Composé} \\ \hline
Krankheit & die & die Krankheiten & krank / erkranken / das Krankenzimmer \\ \hline
Sport & der & \textendash & sportlich / der Sportplatz / die Sportart \\ \hline
Arzt & der & die Ärzte & ärztlich / die Arztpraxis / die Ärztin \\ \hline
Fieber & das & \textendash & fiebern / das Fieberthermometer \\ \hline
Verletzung & die & die Verletzungen & sich verletzen / die Sportverletzung \\ \hline
Training & das & die Trainings & trainieren / der Trainingsplan \\ \hline
Ruhe & die & \textendash & sich ausruhen / der Ruhetag / ruhig \\ \hline
Medizin & die & \textendash & medizinisch / das Medikament \\ \hline
Ernährung & die & \textendash & sich ernähren / das Nahrungsmittel \\ \hline
Schlaf & der & \textendash & schlafen / der Schlafmangel / schläfrig \\ \hline
\end{tabularx}
\end{center}
\textbf{Points de vigilance :} \all{der Arzt $\rightarrow$ die Ärzte} (Umlaut obligatoire) ; \all{der Sport}, \all{das Fieber}, \all{die Ruhe}, \all{die Ernährung} et \all{der Schlaf} sont des noms \emph{non comptables} : ils n'ont pas de pluriel usuel (on ne dit pas \all{die Schlafe} !) ; \all{die Medizin} au sens de « science médicale » n'a pas de pluriel (au sens de « médicament », on préfère \all{das Medikament, die Medikamente}) ; attention au genre de \all{das Training} (anglicisme, toujours neutre).
\end{corrige}

\begin{corrige}[Exercice 3 — Conjugaison des verbes pronominaux au Présent]
\begin{enumerate}
    \item Ich \textbf{fühle mich} heute nicht gut. (Akkusativ : \all{mich})
    \item Warum \textbf{beeilst du dich}? (Akkusativ : \all{dich} ; le pronom suit le sujet dans la question)
    \item Der Spieler \textbf{verletzt sich} am Knie. (3e personne : \all{sich})
    \item Wir \textbf{ruhen uns} nach dem Training eine Stunde \textbf{aus}. (Verbe à particule séparable : \all{aus} renvoyée à la fin)
    \item Meine Schwester \textbf{erkältet sich} schnell. (Akkusativ : \all{sich})
    \item \textbf{Erholen Sie sich} gut? (Forme de politesse \all{Sie} $\rightarrow$ \all{sich} ; inversion sujet-verbe à la question)
    \item Die Kinder \textbf{bewegen sich} im Garten. (3e personne du pluriel : \all{sich})
    \item Er \textbf{wäscht sich} die Hände vor dem Essen. (Ici \all{sich} est au \textbf{Datif} car \all{die Hände} est l'objet direct ; à la 3e personne, datif et accusatif ont la même forme \all{sich}.)
\end{enumerate}
\textbf{Points de vigilance :} le pronom réfléchi se place juste après le verbe conjugué ; avec un verbe séparable, la particule va en fin de phrase : \all{Wir ruhen uns aus.}
\end{corrige}

\begin{corrige}[Exercice 4 — Traduction : conseils avec \textit{sollen} et \textit{müssen}]
\begin{enumerate}
    \item Du \textbf{sollst} mehr Sport \textbf{machen}. (Conseil $\rightarrow$ \all{sollen} ; infinitif en fin de phrase.)
    \item Er \textbf{muss} zum Arzt \textbf{gehen}. (Obligation $\rightarrow$ \all{müssen}.)
    \item Wir \textbf{müssen uns} jetzt \textbf{ausruhen}. (Le pronom réfléchi \all{uns} se place avant l'infinitif ; la particule \all{aus} reste collée à l'infinitif.)
    \item Sie \textbf{sollen nicht arbeiten}, wenn Sie Fieber haben. (Déconseil $\rightarrow$ \all{nicht sollen} ; subordonnée avec \all{wenn} : verbe conjugué \all{haben} à la fin.)
    \item Ich \textbf{muss mir} die Hände \textbf{waschen}. (\textbf{Datif} \all{mir} car \all{die Hände} est l'objet direct — piège classique du francophone qui écrit \all{mich}.)
    \item Sie \textbf{soll} mehr Wasser \textbf{trinken}. (Conseil $\rightarrow$ \all{sollen}, 3e personne : \all{soll}.)
\end{enumerate}
\textbf{Points de vigilance :} ordre rigide : Sujet -- modal conjugué -- compléments -- (pronom réfléchi) -- infinitif. Ne jamais conjuguer l'infinitif final.
\end{corrige}

\begin{corrige}[Exercice 5 — Transformation : Impératif $\rightarrow$ \textit{sollen} / \textit{müssen}]
\begin{enumerate}
    \item Du \textbf{sollst} Sport machen. (Conseil amical $\rightarrow$ \all{sollen}.)
    \item Er \textbf{muss} zum Arzt gehen. (Nécessité médicale forte $\rightarrow$ \all{müssen}.)
    \item Ihr \textbf{sollt} euch ausruhen. (Conseil $\rightarrow$ \all{sollen} ; pronom \all{euch} avant l'infinitif.)
    \item Wir \textbf{müssen} uns die Hände waschen. (Règle d'hygiène $\rightarrow$ \all{müssen} ; \all{uns} datif.)
    \item Sie \textbf{soll} mehr Wasser trinken. (Conseil $\rightarrow$ \all{sollen}.)
\end{enumerate}
\textbf{Points de vigilance :} le choix \all{sollen}/\all{müssen} dépend de la nuance (avis ou obligation) ; vérifier la conjugaison : \all{du sollst}, \all{ihr sollt}, \all{er/sie soll}.
\end{corrige}

\begin{corrige}[Exercice 6 — Texte à trous : Dialog bei Dr. Müller]
\begin{textebox}[Dialog bei Dr. Müller — Lösung]
Arzt: Guten Tag, Herr Nkoto. Was fehlt Ihnen?

Patient: Ich \textbf{fühle} mich sehr schlapp und ich \textbf{habe} seit gestern Fieber.

Arzt: Haben Sie sich \textbf{verletzt} beim Sport?

Patient: Nein, aber ich \textbf{trainiere} seit drei Wochen nicht mehr. Ich \textbf{habe} viel Stress im Büro.

Arzt: Das ist das Problem. Sie \textbf{müssen} sich mehr \textbf{bewegen}. Und Sie \textbf{sollen} Stress \textbf{vermeiden}.

Patient: Soll ich Medikamente nehmen?

Arzt: Nein, Sie \textbf{brauchen} keine Tabletten. Sie \textbf{müssen} sich nur \textbf{erholen} und regelmäßig \textbf{laufen} oder \textbf{schwimmen}.
\end{textebox}
\textbf{Justifications :}
\begin{itemize}
    \item \all{Haben Sie sich verletzt?} : Perfekt d'un verbe pronominal $\rightarrow$ auxiliaire \all{haben} + \all{sich} + participe II \all{verletzt}.
    \item \all{ich trainiere ... nicht mehr} : présent avec \all{seit} (action commencée dans le passé, qui continue).
    \item \all{Sie müssen sich mehr bewegen} : modal conjugué en 2e position, pronom réfléchi puis infinitif à la fin.
    \item \all{Sie brauchen keine Tabletten} : \all{brauchen} + nom, négation avec \all{keine} (et non \all{nicht}).
\end{itemize}
\end{corrige}

\begin{corrige}[Exercice 7 — Appariement : mots composés et expressions]
\begin{enumerate}
    \item \textbf{1--b} : \all{Die Sportart ist zum Beispiel Fußball, Schwimmen oder Joggen.}
    \item \textbf{2--c} : \all{Der Hausarzt ist der Arzt für die Grundversorgung.}
    \item \textbf{3--a} : \all{Die Grippe ist eine Krankheit, die durch Viren entsteht.}
    \item \textbf{4--e} : \all{Das Rezept schreibt der Arzt für die Apotheke.}
    \item \textbf{5--d} : \all{Das Training macht man, um fit zu werden.}
    \item \textbf{6--f} : \all{Man lässt sich impfen zum Schutz vor Krankheiten.}
\end{enumerate}
\textbf{Points de vigilance :} dans les phrases produites, respecter la place du verbe (position 2) ; \all{um ... zu} + infinitif renvoie l'infinitif à la fin : \all{um fit zu werden}.
\end{corrige}

\begin{corrige}[Exercice 8 — Compréhension de texte type Baccalauréat]
\textbf{Fragen auf Deutsch :}
\begin{enumerate}
    \item Laut der Studie treiben nur \textbf{28 Prozent} der Jugendlichen in Deutschland ausreichend Sport.
    \item Die WHO empfiehlt \textbf{mindestens 60 Minuten Bewegung pro Tag}.
    \item Die negativen Folgen sind: \textbf{Übergewicht, Rückenschmerzen und häufige Erkältungen}.
    \item Prof. Wagner schlägt vor, dass \textbf{Eltern Vorbilder sein sollen}: Wenn die Eltern joggen oder Rad fahren, machen die Kinder oft mit.
    \item Die Experten fordern \textbf{mehr Schwimmunterricht und sichere Radwege}, damit sich Kinder selbstständig bewegen können.
\end{enumerate}
\textbf{Questions en français :}
\begin{enumerate}
    \setcounter{enumi}{5}
    \item L'expression signifie que la santé ne se construit pas d'abord chez le médecin (soins curatifs), mais dans la vie quotidienne : bouger régulièrement, bien manger, dormir assez. Le texte le montre par les exemples des pauses d'exercice à l'école, de l'exemple des parents et des infrastructures (piscines, pistes cyclables) : la prévention au quotidien évite la maladie.
    \item \textbf{Similitudes :} (1) Dans les deux pays, les jeunes passent beaucoup de temps devant les écrans (smartphone, jeux vidéo) et bougent moins. (2) Dans les deux cas, les écoles lancent des programmes de promotion du sport (« Fit für die Schule » en Allemagne, « Sport pour la Santé » au Cameroun). \textbf{Différences :} (1) En Allemagne, le sport est souvent organisé en clubs (\all{Vereine}) avec des infrastructures (pistes cyclables, piscines) ; au Cameroun, le sport est plus informel (football sur la place du quartier, danse bendskin). (2) L'encadrement médical et scolaire du sport est plus développé en Allemagne.
    \item \textbf{Verbes pronominaux au présent :} \all{Wer sich ... fit hält} (\all{sich fit halten}) ; \all{damit sich Kinder selbstständig bewegen können} (\all{sich bewegen}). \textbf{Verbes modaux :} \all{Kinder müssen sich wieder mehr bewegen} (\all{müssen}) ; \all{Eltern sollen Vorbilder sein} (\all{sollen}) ; \all{damit sich Kinder ... bewegen können} (\all{können}).
\end{enumerate}
\textbf{Barème indicatif (sur 20) :} questions 1--5 : 2 pts chacune (contenu 1 pt + correction linguistique 1 pt) ; question 6 : 3 pts ; question 7 : 4 pts (2 similitudes + 2 différences) ; question 8 : 3 pts (identification correcte des verbes avec sujets).
\textbf{Points de vigilance :} répondre en phrases complètes ; ne pas inventer d'informations absentes du texte ; pour la question 8, citer le verbe \emph{avec son sujet} comme demandé.
\end{corrige}

\begin{corrige}[Exercice 9 — Thème et Version]
\textbf{A. Thème (Français $\rightarrow$ Allemand)}
\begin{enumerate}
    \item Mein Freund Paul ist oft krank. Er \textbf{hat sich gestern erkältet}. (Perfekt du verbe pronominal \all{sich erkälten} : auxiliaire \all{haben} + \all{sich} + participe II \all{erkältet} ; la particule \all{er-} est inséparable, donc pas de \all{ge-}.)
    \item Er \textbf{muss} zum Arzt \textbf{gehen} und er \textbf{muss sich ausruhen}. (Modal en position 2, infinitif à la fin ; pronom \all{sich} avant l'infinitif.)
    \item Der Arzt sagt zu ihm: „Du \textbf{sollst} jede Woche Sport \textbf{machen}.“ (Conseil $\rightarrow$ \all{sollen} ; guillemets allemands „\dots“ ; majuscule à \all{Sport}.)
    \item Paul antwortet: „Ich \textbf{habe} keine Zeit, ich \textbf{muss} viel \textbf{lernen}.“ (Négation du nom : \all{keine Zeit}, pas \all{nicht Zeit}.)
    \item Aber ohne Sport \textbf{wird} er sich nicht besser \textbf{fühlen}. (Futur I : \all{werden} conjugué en position 2 + infinitif \all{fühlen} à la fin ; après \all{aber ohne Sport}, inversion sujet-verbe : \all{wird er}.)
\end{enumerate}
\textbf{B. Version (Allemand $\rightarrow$ Français)}

Ce matin, Lisa ne se sent pas bien. Elle a de la fièvre et mal à la gorge. Elle doit rester au lit et se reposer. Sa mère appelle le médecin. Le médecin vient et dit : « Lisa doit boire beaucoup de thé et ne pas faire d'efforts. Elle ne doit pas aller à l'école demain. » Lisa est triste, mais elle sait que la santé est la chose la plus importante. Elle veut se rétablir vite pour pouvoir rejouer au football.

\textbf{Points de vigilance :} \all{sich anstrengen} = faire des efforts (faux ami : pas « s'étrangler » !) ; \all{Halsschmerzen} = mal de gorge ; \all{um wieder Fußball zu spielen} = « pour rejouer au football » (but).
\end{corrige}

\begin{corrige}[Exercice 10 — Production écrite : lettre de conseils]
\begin{textebox}[Beispielbrief (ca. 75 Wörter)]
Lieber Alex,

ich habe gehört, dass du dich oft müde \textbf{fühlst} und oft krank bist. Das macht mir Sorgen. Du \textbf{sollst} dich gesünder \textbf{ernähren}: mehr Gemüse und Obst essen und weniger Süßes. Du \textbf{musst} dich auch mehr \textbf{bewegen}: Geh in einen Verein, jogge oder schwimme zweimal pro Woche! Weniger Bildschirm am Abend und mehr Schlaf sind auch wichtig. Und nach dem Training \textbf{sollst} du dich gut \textbf{ausruhen}. Dann wirst du dich bald besser fühlen!

Viele Grüße

Dein Freund Paul
\end{textebox}
\textbf{Vérification des contraintes :}
\begin{itemize}
    \item Longueur : environ 75 mots (dans la fourchette 60--80).
    \item Verbes pronominaux (au moins 3) : \all{sich fühlen}, \all{sich ernähren}, \all{sich bewegen}, \all{sich ausruhen} (4 utilisés).
    \item Modaux : \all{sollst} (2 fois) et \all{musst} (1 fois).
    \item Structure : salutation (\all{Lieber Alex}), corps avec conseils, formule finale (\all{Viele Grüße}), signature.
    \item Vocabulaire du thème : \all{das Gemüse}, \all{das Obst}, \all{der Schlaf}, \all{der Bildschirm}, \all{der Verein}, \all{joggen}, \all{schwimmen}.
\end{itemize}
\textbf{Barème indicatif (sur 10) :} respect des contraintes (longueur, 3 pronominaux, \all{sollen}/\all{müssen}) : 4 pts ; correction grammaticale (ordre des mots, cas, conjugaison) : 3 pts ; vocabulaire et orthographe (majuscules, Umlaute) : 2 pts ; structure de la lettre : 1 pt.
\textbf{Points de vigilance :} après \all{dass}, le verbe conjugué va à la fin (\all{dass du dich oft müde fühlst}) ; virgule obligatoire après la salutation et minuscule au début de la ligne suivante ; \all{Viele Grüße} avec ß.
\end{corrige}

\newpage

\coursec{Fiche bilan}

\begin{popfigure}
\begin{tikzpicture}[mindmap, grow cyclic, every node/.style={concept}, concept color=popblue!40, root concept/.append style={concept color=popblue, font=\bfseries\small, text=white}, level 1/.append style={level distance=3.4cm, sibling angle=51.4}, level 1 concept/.append style={font=\scriptsize, text=popdark}]
\node[root concept, align=center]{ALA12\\ Maladies\\ et Sport}
  child[concept color=popgreen!45]{ node[align=center]{Texte\\ \all{Sport ist die beste Medizin}} }
  child[concept color=poporange!50]{ node[align=center]{Lexique\\ santé, maladie, sport} }
  child[concept color=poppurple!45]{ node[align=center]{Verbes\\ pronominaux\\ \all{sich erkälten}} }
  child[concept color=poppink!50]{ node[align=center]{Conseil\\ \all{sollen / müssen}} }
  child[concept color=popgold!55]{ node[align=center]{Conjugaison\\ présent + pronoms\\ \all{mich, dich, sich}} }
  child[concept color=popblue!35]{ node[align=center]{Compréhension\\ globale puis détaillée} }
  child[concept color=popgreen!30]{ node[align=center]{Production\\ résumé, conseils} };
\end{tikzpicture}
\legende{Carte mentale de la leçon ALA12 : les maladies et l'importance de la pratique du sport.}
\end{popfigure}

\begin{bilan}
\textbf{1. Lexique clé (à connaître avec article et pluriel)}
\begin{itemize}
  \item \all{die Krankheit, -en} : la maladie ; \all{die Erkältung, -en} : le rhume ; \all{das Fieber} : la fièvre ; \all{der Husten} : la toux ; \all{die Grippe, -n} : la grippe ; \all{der Schmerz, -en} : la douleur ; \all{die Kopfschmerzen} (pl.) : le mal de tête.
  \item \all{der Arzt, die Ärzte} : le médecin ; \all{die Ärztin, -nen} : la femme médecin ; \all{das Krankenhaus, die Krankenhäuser} : l'hôpital ; \all{die Medizin, -en} : le médicament / la médecine ; \all{das Rezept, -e} : l'ordonnance.
  \item \all{der Sport} : le sport ; \all{das Training} : l'entraînement ; \all{die Gesundheit} : la santé ; \all{trainieren} : s'entraîner ; \all{gesund} : en bonne santé ; \all{krank} : malade ; \all{fit} : en forme ; \all{müde} : fatigué.
  \item Familles de mots : \all{krank} $\rightarrow$ \all{die Krankheit}, \all{das Krankenhaus}, \all{krank sein} ; \all{gesund} $\rightarrow$ \all{die Gesundheit}, \all{gesund leben}.
\end{itemize}

\textbf{2. Grammaire à connaître par c\oe ur}
\begin{center}
\begin{tabularx}{\linewidth}{|X|X|X|}
\hline
\rowcolor{popblueL} Règle & Forme & Exemple \\
\hline
Verbe pronominal au présent & pronom réfléchi (acc.) : \all{mich, dich, sich, uns, euch, sich} placé après le verbe conjugué & \all{Ich erkälte mich oft.} « Je m'enrhume souvent. » \\
\hline
Verbe pronominal à l'infinitif & \all{sich} + infinitif, en fin de phrase & \all{Man muss sich bewegen.} « Il faut bouger. » \\
\hline
Le conseil avec \all{sollen} & modal en position 2 + infinitif en fin & \all{Du sollst mehr Sport machen.} « Tu devrais faire plus de sport. » \\
\hline
L'obligation avec \all{müssen} & modal en position 2 + infinitif en fin & \all{Kranke müssen zum Arzt gehen.} « Les malades doivent aller chez le médecin. » \\
\hline
Négation du conseil & \all{nicht} + modal : conseil de ne pas faire & \all{Du sollst nicht rauchen.} « Tu ne devrais pas fumer. » \\
\hline
\end{tabularx}
\end{center}

\textbf{3. Conjugaison modèle : \all{sich erkälten} (s'enrhumer) au présent}
\begin{center}
\begin{tabular}{|l|l|l|}
\hline
\rowcolor{popgreenL} \all{ich erkälte mich} & \all{du erkältest dich} & \all{er/sie/es erkältet sich} \\
\hline
\all{wir erkälten uns} & \all{ihr erkältet euch} & \all{sie/Sie erkälten sich} \\
\hline
\end{tabular}
\end{center}

\textbf{4. Méthodes express}
\begin{itemize}
  \item \cle{Compréhension} : lire le titre et deviner le thème $\rightarrow$ lecture globale (idée générale) $\rightarrow$ lecture détaillée (souligner les mots-clés) $\rightarrow$ répondre par des phrases complètes en allemand.
  \item \cle{Résumé} : reprendre les idées principales avec ses propres mots, au présent, en 3 à 5 phrases (\all{Der Text handelt von...}).
  \item \cle{Donner un conseil} : \all{sollen} pour une recommandation, \all{müssen} pour une obligation ; toujours l'infinitif à la fin.
\end{itemize}
\end{bilan}

\begin{aretenir}[Checklist avant le Bac]
\begin{itemize}
  \item $\square$ Je connais le vocabulaire de la santé et du sport avec l'article et le pluriel de chaque nom.
  \item $\square$ Je sais conjuguer un verbe pronominal au présent avec \all{mich, dich, sich, uns, euch, sich}.
  \item $\square$ Je place correctement le pronom réfléchi juste après le verbe conjugué.
  \item $\square$ Je sais former un conseil avec \all{sollen} et une obligation avec \all{müssen} (modal en position 2, infinitif à la fin).
  \item $\square$ Je n'oublie pas la majuscule à tous les noms allemands (\all{die Krankheit}, \all{das Fieber}).
  \item $\square$ Je sais que \all{der Sport} ne se met pas au pluriel et que \all{die Medizin} signifie aussi « le médicament ».
  \item $\square$ Je peux résumer le texte \all{Sport ist die beste Medizin} en 3 à 5 phrases simples.
  \item $\square$ Je sais répondre aux questions de compréhension par des phrases complètes.
  \item $\square$ Je peux écrire un court paragraphe de conseils pour rester en bonne santé.
  \item $\square$ Je vérifie toujours la place du verbe : position 2 dans la phrase déclarative, infinitif en fin avec les modaux.
\end{itemize}
\end{aretenir}

\findecours

\end{document}
